\documentclass[11pt]{article}
\pdfoutput=1
\usepackage[utf8]{inputenc}
\usepackage[T1]{fontenc}
\usepackage{lmodern}
\usepackage{microtype}
\usepackage[margin=1in]{geometry}
\usepackage[shortlabels]{enumitem}
\usepackage[backgroundcolor=lightgray,linecolor=lightgray,textsize=footnotesize]{todonotes}
\usepackage{graphicx}

% Mathematics.
\usepackage{double-data-macros}
\numberwithin{equation}{section}

% References.
\usepackage[breaklinks,hidelinks,bookmarksnumbered]{hyperref}
\usepackage[capitalize,noabbrev,nameinlink]{cleveref}

% Theorems.
\usepackage{thmtools}
\declaretheorem[within=section]{theorem}
\declaretheorem[sibling=theorem]{proposition}
\declaretheorem[sibling=theorem]{corollary}
\declaretheorem[sibling=theorem]{lemma}
\declaretheorem[sibling=theorem,style=definition,qed=\qedsymbol]{definition}
\declaretheorem[sibling=theorem,style=definition,qed=\qedsymbol]{construction}
\declaretheorem[sibling=theorem,style=remark,qed=\qedsymbol]{remark}
\declaretheorem[sibling=theorem,style=remark,qed=\qedsymbol]{discussion}
\declaretheorem[sibling=theorem,style=remark,qed=\qedsymbol]{example}

% Diagrams.
\usepackage{nicematrix}
\NiceMatrixOptions{cell-space-limits=3pt}
\newenvironment{dblArray}[1]{\begin{NiceArray}{#1}[hvlines]}{\end{NiceArray}}

% Bibliography.
\usepackage[style=alphabetic,maxbibnames=10]{biblatex}
\addbibresource{double-data-bib.bib}
\renewbibmacro{in:}{}

% Title.
\title{Quantification in Double-Categorical Database Schemas}
\author{Michael Lambert}
\date{August 2026}

\begin{document}
\maketitle

\begin{abstract}
  Double-categorical database schemas are enriched with universal quantification in the form of right adjoints to substitution. This allows phrasing of the important query, \emph{relational division}. It is shown that such right adjoints together with suitable tabulators interpret modal operators, provide cartesian closed structure, and, when combined with global cocartesian structure, interpret first-order predicate logic and thus description logic. These structures are applied throughout to querying and optimization. It is seen, for example, that Frobenius reciprocity and Beck-Chevalley hold in any suitably structured double database schema and that these provide pushdown optimization rules. Likewise, negation queries are introduced and studied as a result of cocartesian and local implication structure. 
\end{abstract}

\section{Introduction}

The double-categorical approach to database schematization and knowledge representation (KR) began in \cite{lambert2025} is in this paper extended by the introduction of quantification as a further operation on double ologs. Such double ologs are a synthesis of the functional \cite{Spivak2010,spivak2012} and relational \cite{patterson2017} approaches to data schemas within the single framework of a suitably structured cartesian equipment known as a `double category of relations' \cite{lambert2022}. A data instance on such a double olog is a cartesian double functor valued in relations. Standard relational database query operations such as select, filter and join are all native operations given by the equipment and cartesian structure of the given double olog; instances preserve all such structure and compute the query results directly in relations. This has the effect of bypassing the foregoing standard of lifting to induced external adjunctions between functor categories to perform functorial data migration. Up until the present work, however, more advanced queries such as \emph{relational division} \cite{codd1972} were not known to be possible within this double-categorical framework.

To perform such queries, the present work studies the effect of asking for right adjoints to the substitution functors coming from the equipment structure in any double olog. These right adjoints are \emph{dependent products} or \emph{fibered products} in the language of \emph{fibrational semantics} \cite{lawvere1970, jacobs1999}. These dependent products have appeared in purely bicategorical contexts as \emph{Kan quantification} \cite{lawvere2002}. When combined in the present situation with suitable restrictions and tabulators they give (1) division operators, (2) modal operators, and (3) exponentials. Each of these, the queries they allow, and their interrelations with the standing modular laws, Frobenius reciprocity and Beck-Chevalley condition, are studied here in detail. The culminating result is that a double olog with dependent products, strong tabulators, and coproducts satisfying a global distributivity law is locally a model of first-order predicate logic enhanced with modal operators modelling possibility (liveness) and necessity (safety). Thus, such ologs also interpret \emph{description logic} \cite{baader2010} which forms the basis of OWL and web3. In this way, the present paper advances the project of seeing structured equipments as a models of logics over type theories \cite{jacobs1999}.

The theoretical narrative of this development is itself of interest. The external Beck-Chevalley condition is shown to hold in any `double category of relations' with strong tabulators (\cref{lemma:beck-chevalley}). This results from the theoretical advance of \cite{hoshino2025} showing that the condition of \emph{unit-purity} is automatically implied by the assumption of \emph{discreteness} \cite[Definition 2.1]{carboni1987} in any `double category of relations'. Formerly \cite{lambert2022}, unit-purity figured as a facilitating but \emph{ad hoc} assumption that is now seen to be redundant. Likewise, we show that Frobenius reciprocity holds in any `double category of relations' (\cref{prop:Frobenius-reciprocity}). This was asserted in \cite{lambert2022} and proved in appropriate references. It is seen here to be a consequence of the modular laws stemming from the fact that any `double category of relations' is compact closed with a trivial duality involution \cite[Theorem 2.4]{carboni1987}. Frobenius reciprocity is used to show that any double olog with dependent products and strong tabulators is locally cartesian closed in the sense that every hom category is cartesian closed (\cref{theo:locally-cartesian-closed}). The crucial result enabling this (\cref{lemma:identity-proarrow-exponentiable}) shows that identity proarrows are always exponentiable in any `double category of relations'. Essentially, the discreteness axiom forces the identity proarrows to act as very strict identity predicates akin to genuine diagonals. We thus recover formulas analogous to those of \cite[Theorem 10.5.4]{jacobs1999} without the need of elaborate comprehension structures. Nor do we require that the base category is cartesian closed. This establishes a direct connection between dependent (co)products, tabulators, and exponents. Finally, it is shown that global cocartesian structure with a global distributive law localizes to give local coproducts over which local products distribute. In this sense, a main result is that distributivity localizes (\cref{theorem:local-distributivity}). 

Several difficult queries in relational database theory are captured in this framework. The first is the relational division operation originating in \cite{codd1972}. The classical phrasing of this operation has a notoriously difficult syntax and is a major source of computational bottlenecking in implementations. Here it is handled by precomposing a well-chosen dependent product with a suitable restriction functor. The inevitability of this development is showcased by the first three sections (\cref{section:relational-division,section-ologs,section:varieties-of-quantification}) which have a discursive pacing and recount the narrative of the present discovery of the solution. Of course, inasmuch as right adjoints to substitution are already known to encode division (hence the very notion of a \emph{division allegory} \cite{freyd1990} and also \cite[\S 11]{lambert2022}, this solution is in hindsight probably the obvious one, but to our knowledge the connection has not been formalized for categorical database querying. Likewise, negation queries are captured using cartesian closed structure and local coproducts. These are combined with existing joins to do otherwise difficult compound queries with relative ease (\cref{section:cocartesian-negation}). Modal queries involving possibility and necessity operators are also introduced as applications of dependent products and tabulators (\cref{section:modality}). These bear a superficial resemblance to division queries but are indeed closer to liveness and safety properties \cite{lamport1977,alpern1985,alpern1987}. Disjunction queries are studied for the first time under the guise of \emph{local coproducts} which are induced from global cocartesian structure (\cref{section:cocartesian-negation}).

A major application of the present work is to query optimization and rewrite rules. Certainly, join is generally regarded as the most costly of standard queries; and select is essentially free. This is borne out for us by the fact that join is typically a complicated restriction along diagonals and select is achieved by a simple projection from a factor \cite{lambert2024}. Filter and collapse are closer overall. Filter is executed by restricting along a value or subtype; whereas collapse is best thought of as \emph{extraction of metadata} or \emph{abstraction}, that is, basically, extension along a surjection. Inasmuch as the latter is something like passing to equivalence classes, or a taking a quotient, or a colimit, we regard it as generally more expensive than filter, but still less than join. We adopt a heuristic hierarchy of relative query expense: 
  \begin{equation*}\label{equation:query-expense}
    \text{select column } \lessapprox  \text{ filter by value } \lessapprox \text{ collapse/abstract } \lessapprox \text{ join/conjunction}.
  \end{equation*}
Relative to this hierarchy, we see both the Beck-Chevalley condition and Frobenius reciprocity as rewrite and optimization rules for certain compound queries. We show (\cref{remark:Beck-Chevalley-Optimization,example:beck-chevalley}) that the former is a \emph{filter-pushdown} optimization that takes a collapse followed by filter and rewrites it as a more efficient filter followed by a collapse. In our illustrating example, the more efficient query ends up being a filter followed by a select column. Frobenius reciprocity we show (\cref{remark:Frobenius-Optimization,example:frobenius}) is a \emph{join-pushdown} which rewrites a collapse following an expensive join/conjunction as a join following a collapse. This does not eliminate the need of a join but the collapse/abstract first reduces the size of the dataset on which it is performed. The point of deriving Beck-Chevalley and Frobenius reciprocity in our ologs is that these rewrite and optimization rules are thus native features of the database schemas and preserved by the functorial and fibrational semantics of any data instance. Such optimization is thus in this sense baked-in.

This project began several years ago with the work leading up to \cite{lambert2025} where it was noticed that many of the examples of migration in \cite{Spivak2010} were queries that could be realized by pullbacks and projections. Double categories had already entered our conversation on (KR) as a result of the desirability of a general framework where genuinely relational aspects attained a first-class status. And in this context, it was realized that cartesian equipments carried exactly the derived structure to model natively select, filter and join as operations internal to a double olog generated by a presentation. Both ingredients are indispensable. Others such as the hypothesis of being a `double category of relations' and having tabulators are highly convenient inasmuch as (for example) the derived property of \emph{compactness} seems to solve a difficult bookkeeping problem introduced by the sidedness of relations. Of the two necessary ingredients, the equipment structure, that is, the fact that the source-target projection is a fibration, is the subtle and rich one. The use of fibrations to model databases is not exactly new (e.g. \cite{rosebrugh1992,islam1994}). By happy coincidence the power of a fibrational perspective was rediscovered in the present and foregoing work on streamlining the account of data querying. Genuinely novel is the introduction of the technical machinery of structured double categories presenting and formalizing ologs; and its combination with the idea that instancing via both double-categorical functorial and fibrational semantics amounts to database querying. This paper is an expos\'e on the consequences especially of the latter aspect.

To conclude this introduction, a summary of the main results of the paper is as follows:
  \begin{enumerate}
    \item Relational division is recast as a composite operation in a suitably structured double category using dependent products (\cref{example:recast-relational-division}).
    \item Beck-Chevalley and Frobenius reciprocity are realized as query optimization rewrite rules (\cref{remark:Beck-Chevalley-Optimization}, \cref{remark:Frobenius-Optimization}).
    \item Dependent products and strong tabulators are seen to give local cartesian closure and a formula for exponentials (\cref{theo:locally-cartesian-closed}).
    \item Further cocartesian structure allows negation queries (\cref{example:negation}).
    \item Distributivity localizes (\cref{theorem:local-distributivity}) without the assumption of cartesian closure or even the existence of dependent products.
    \item Any suitably structured double olog, viewed as an equipment, interprets first-order predicate logic, hence all the traditional operations of relational algebra (\cref{corollary:first-order-fibration}).
    \item What we call a \emph{FOML double olog} interprets also modal necessity and possibility operators and thus gives a semantics of \emph{description logic} (\cref{definition:interpret-description-logic}).
  \end{enumerate}
Double-categorical conventions are consistent with the previous \cite{lambert2024,lambert2025}. Our \cite{lambert2022} on `double categories of relations' is taken for granted especially for the crucial notion of \emph{discreteness}. Likewise \cite{carboni1987} is an indispensable reference for the bicategorical description of \emph{compactness}. \cite{jacobs1999} is a guiding resource for fibrational interpretations of logics over type theories. As with our previous work on double ologs, examples of data instances are taken from a non-standard source. In this instance they are taken from the fantasy RPG, The Elder Scrolls V: Skyrim.

\tableofcontents

\part{Dependent Products \& Relational Division}

This first part of the paper is devoted entirely to setting up, explaining, and recasting in appropriate double-categorical machinery the queries we are interested in. The main motivating query is again \emph{relational division} introduced in \cite{codd1972} essentially as a \emph{check inventory against list} query. The first three sections (\cref{section:relational-division,section-ologs,section:varieties-of-quantification}) are a discussion-style narrative recounting the original syntax of relational division and showing in summary a form of the analysis that led to the introduction of dependent products and the formulation of the solution as presented in \cref{example:recast-relational-division}. Dependent coproducts and products are realized as species of \emph{Kan quantification} \cite{lawvere2002}. We study Beck-Chevalley in this connection as a rewrite and optimization rule (\cref{section:existential-queries}). On this basis, universally-quantified queries are introduced using dependent products in \cref{section:universal-queries} and leading to the formulation of relation division. ologs with this kind of structure are axiomatized as $\prod$-double ologs. Along the way, we survey 6 different kinds of queries (\cref{fig:quant-queries}), each of which is derived from quantification over a relation in some form. Two of these are modal queries: one possibility and the other necessity. The latter bears a superficial resemblance to division but is better captured using suitably structured tabulators (\cref{section:modality}).

\section{Relational Division}
\label{section:relational-division}
Relational division is a query introduced in \cite[\S 2]{codd1972} that behaves essentially as a universal quantifier over given specified values in the column of a table. First we illustrate this operation with a generic example. We work with binary relations $R \colon C\proto A$ which are of course equivalently monic arrows $\langle d,c\rangle \colon R\rightarrowtail C\times A$. Recall the notation from the reference
  \begin{equation*}
    r[A] = c(r) \text{ where } r\in R
  \end{equation*}
and likewise for $r[C] = d(r)$. Likewise, for any binary relation $T$, let $g_T(x) = \lbrace y\mid xTy\rbrace$, that is, the set of all $y$ in the target of $T$ related to $x$ under $T$. Call this the \textbf{$T$-orbit} of $x$. Now, for the example, consider the two tables:
  \begin{equation*}
    \begin{tabular}{| l | l | }
        \hline\multicolumn{2}{| c |}{\bf R}\\
        \hline {\bf C }&{\bf A}\\
        \hline 1 & a \\
        \hline 1 & b \\
        \hline 1 & c \\
        \hline 2 & a \\
        \hline 2 & b \\
        \hline 3 & b \\
        \hline
    \end{tabular} \qquad\qquad\qquad
    \begin{tabular}{| l | l | }
        \hline\multicolumn{2}{| c |}{\bf S}\\
        \hline {\bf B }&{\bf D}\\
        \hline a & 2 \\
        \hline b & 1 \\
        \hline b & 2 \\
        \hline a & 3 \\
        \hline
    \end{tabular}.
  \end{equation*}
Note that $A$ and $B$ have entries from the same set $\lbrace a,b,c\rbrace$. This is required to perform division. Now, in the notation of the reference, a division query is 
    \begin{equation*}
      R[A\div B]S = \lbrace r[C]\mid r\in R \text{ and } S[B]\subset g_R(r[C])\rbrace = \lbrace 1,2\rbrace.
    \end{equation*}
We unpack this carefully to illustrate the operation. The possible values of $r[C]$ are the distinct entries of the $C$-column of $R$, namely, $1$, $2$, and $3$. Likewise, $S[B] = \lbrace a,b\rbrace$ is the set of all values of the $B$-column of $S$. So, there are three sets of the form $g_R(r[C])$, only two of which contain $S[B]$, as in: 
    \begin{align*}
      S[B] &= \lbrace a,b\rbrace \nsubseteq \lbrace b\rbrace = g_R(3) \\
      S[B] &= \lbrace a,b\rbrace \subset \lbrace a,b,c\rbrace = g_R(1)\\
      S[B] &= \lbrace a,b\rbrace = \lbrace a,b\rbrace = g_R(2).
    \end{align*}
Hence the division query returns $\lbrace 1,2\rbrace$, exactly \emph{the set of all $C$-values whose orbit contains all the values of the $B$-column of $S$}. More prosaically, a $C$-value is selected by the query iff it is related to everything on the list of $B$-values. Thus, \emph{universal quantification} enters as an essential aspect of the query. Now, in the containment relations above, the universally quantified set does not need to completely describe the orbit of any given element satisfying the division query; that is, $1$ is related to $c$, but this is not a $B$-value. Likewise $3$ is related to something on the list, but not everything.

A more concrete example illustrating the same query is the following. Consider the two relational tables below. On the left, rows pair an apothecary with an ingredient sold there; on the right, rows pair an ingredient with an effect produced by the combination of that ingredient with another in a potion. 
    \[
    \begin{tabular}{| l | l | }
        \hline\multicolumn{2}{| c |}{\bf Inventory}\\
        \hline {\bf Apothecary }&{\bf Ingredient}\\
        \hline Arcadia's Cauldron & giant's toe \\
        \hline Arcadia's Cauldron & wheat \\
        \hline Arcadia's Cauldron & void salts \\
        \hline Elgrim's Elixers & giant's toe \\
        \hline Elgrim's Elixers & wheat \\
        \hline The White Phial & wheat \\
        \hline
    \end{tabular} \qquad\qquad
    \begin{tabular}{| l | l | }
        \hline\multicolumn{2}{| c |}{\bf Potion}\\
        \hline {\bf Ingredient }&{\bf Effect}\\
        \hline giant's toe & fortify health \\
        \hline wheat & fortify health \\
        \hline wheat & lingering damage magicka \\
        \hline giant's toe & damage stamina regen \\
        \hline
    \end{tabular}
  \]
Note that the second table does not list the other ingredients required to create the potion producing the effect. Our expertise is that any two such ingredients with the same effect will create the desired potion when combined. Now, suppose we want to create a fortify health potion but are to lazy to fight a giant to get his toe. We would also really prefer to one-stop-shop (maybe the merchant will give us a deal if we buy in bulk). The orbits are then the inventories:
    \begin{align*}
      g_{\text{Inventory}}(\text{Arcadia's}) &= \lbrace \text{giant's toe, wheat, void salts}\rbrace \\
      g_{\text{Inventory}}(\text{Elgrim's}) &= \lbrace \text{giant's toe, wheat}\rbrace \\
      g_{\text{Inventory}}(\text{White Phial}) &= \lbrace \text{wheat} \rbrace
    \end{align*}
The the set of values quantified over is our shopping list: $\lbrace \text{giant's toe, wheat}\rbrace$. In the notation above, the division query is then 
    \begin{align*}
      &\text{Inventory}[\text{Ingredient}\div \text{Ingredient}]\text{Potion} \\
      =\;&\lbrace \text{Apothecary} \mid \lbrace \text{giant's toe, wheat}\rbrace\subset g_\text{Inventory}(\text{Apothecary})\rbrace
    \end{align*}
which is just to say that we are looking for the proprietors who have all the ingredients on our shopping list. The result is of course $\lbrace \text{Arcadia's, Elgrims}\rbrace$. It does not matter that Arcadia's also has void salts; but it does matter that The White Phial only has wheat but no giant's toe.

\begin{remark}
  There is something awkward about these examples inasmuch as we would probably more likely have tables better reflecting our alchemical expertise such as 
      \[
      \begin{tabular}{| l | l | }
          \hline\multicolumn{2}{| c |}{\bf Apothecary}\\
          \hline {\bf Name }&{\bf Ingredient}\\
          \hline Arcadia's Cauldron & giant's toe \\
          \hline Arcadia's Cauldron & wheat \\
          \hline Arcadia's Cauldron & void salts \\
          \hline Elgrim's Elixers & giant's toe \\
          \hline Elgrim's Elixers & wheat \\
          \hline The White Phial & wheat \\
          \hline
      \end{tabular} \qquad\qquad
      \begin{tabular}{| l | l | }
          \hline\multicolumn{2}{| c |}{\bf Potion}\\
          \hline {\bf Ingredient }&{\bf Effect}\\
          \hline giant's toe & fortify health \\
          \hline wheat & fortify health \\
          \hline wheat & lingering damage magicka \\
          \hline giant's toe & damage stamina regen \\
          \hline swamp fungal pod & restore health \\
          \hline void essence & restore health\\
          \hline
      \end{tabular}.
    \]
  However, division with quantification over the whole the Ingredient column in the Potion table now produces the empty set since no merchant has all those items on hand. And \cite[\S 2.4]{codd1972} recognizes this in the example queries. Namely, example query \#8 to find suppliers supplying at least parts 12, 13 and 15 restricts the part \# values to only 12, 13 and 15 without ranging over the whole column before executing the division query. This is exactly a \emph{shopping list} kind of example. In other words, technically, the division query is not set up to do selection from a sub-list of a column on its own. A restriction to just the values of interest from the second table is employed implicitly before doing the division query. In the example immediately above, to execute our query, we would filter the rows with the value ``fortify health'' first and then do the division query with the ingredients column. The flexibility of the division query as written is that it can be applied to any two tables with columns having the same types of values; the expense is that it might not return anything and may require restrictions to produce results in intuitive cases.
\end{remark}

\section{Ologs for Division}
\label{section-ologs}

Having reviewed relational division, we continue the discussion and now recast the situation from a double-categorical olog-perspective \cite{lambert2025}. The data is schematized by two basic relations
    \begin{equation*}
      % https://q.uiver.app/#q=WzAsMixbMCwwLCJcXGZib3h7QXBvdGhlY2FyeX0iXSxbMiwwLCJcXGZib3h7SW5ncmVkaWVudH0iXSxbMCwxLCJcXHRleHR7SW52ZW50b3J5fSIsMCx7InN0eWxlIjp7ImJvZHkiOnsibmFtZSI6ImJhcnJlZCJ9fX1dXQ==
      \begin{tikzcd}
        {\fbox{Apothecary}} && {\fbox{Ingredient}}
        \arrow["{\text{Inventory}}"{inner sep=.8ex}, "\shortmid"{marking}, from=1-1, to=1-3]
      \end{tikzcd}\qquad\qquad
      % https://q.uiver.app/#q=WzAsMixbMCwwLCJcXGZib3h7SW5ncmVkaWVudH0iXSxbMiwwLCJcXGZib3h7RWZmZWN0fSJdLFswLDEsIlxcdGV4dHtQb3Rpb259IiwwLHsic3R5bGUiOnsiYm9keSI6eyJuYW1lIjoiYmFycmVkIn19fV1d
      \begin{tikzcd}
        {\fbox{Ingredient}} && {\fbox{Effect}}
        \arrow["{\text{Potion}}"{inner sep=.8ex}, "\shortmid"{marking}, from=1-1, to=1-3]
      \end{tikzcd}.
    \end{equation*}
Suppose that these relations are instanced by the two tables at the end of the previous section, thought of as living in $\Rel$. Now, the olog generated will be a cartesian equipment. So, given an individual as a global element of a type, say, $\text{Arcadia's}\colon 1\to \fbox{Name}$, we can form the restriction
    \begin{equation*}
      % https://q.uiver.app/#q=WzAsNCxbMCwxLCJcXGZib3h7QXBvdGhlY2FyeX0iXSxbMywxLCJcXGZib3h7SW5ncmVkaWVudH0iXSxbMCwwLCIxIl0sWzMsMCwiXFxmYm94e0luZ3JlZGllbnR9Il0sWzAsMSwiXFx0ZXh0e0ludmVudG9yeX0iLDIseyJzdHlsZSI6eyJib2R5Ijp7Im5hbWUiOiJiYXJyZWQifX19XSxbMiwwLCJcXHRleHR7QXJjYWRpYSdzfSIsMl0sWzIsMywiXFx0ZXh0e0FyY2FkaWEncyBJbnZlbnRvcnl9IiwwLHsic3R5bGUiOnsiYm9keSI6eyJuYW1lIjoiYmFycmVkIn19fV0sWzMsMV0sWzYsNCwiXFx0ZXh0e3Jlc30iLDEseyJzaG9ydGVuIjp7InNvdXJjZSI6MjAsInRhcmdldCI6MjB9LCJzdHlsZSI6eyJib2R5Ijp7Im5hbWUiOiJub25lIn0sImhlYWQiOnsibmFtZSI6Im5vbmUifX19XV0=
      \begin{tikzcd}
        1 &&& {\fbox{Ingredient}} \\
        {\fbox{Apothecary}} &&& {\fbox{Ingredient}}
        \arrow[""{name=0, anchor=center, inner sep=0}, "{\text{Arcadia's Inventory}}"{inner sep=.8ex}, "\shortmid"{marking}, from=1-1, to=1-4]
        \arrow["{\text{Arcadia's}}"', from=1-1, to=2-1]
        \arrow[from=1-4, to=2-4]
        \arrow[""{name=1, anchor=center, inner sep=0}, "{\text{Inventory}}"'{inner sep=.8ex}, "\shortmid"{marking}, from=2-1, to=2-4]
        \arrow["{\text{res}}"{description}, draw=none, from=0, to=1]
      \end{tikzcd}
    \end{equation*}
which in the instance is precisely the named merchant's inventory. This would play the role of $g_R(x)$ in the previous section. Likewise, our ingredient list for a fortify health potion is schematized by a restriction:
    \begin{equation*}
      % https://q.uiver.app/#q=WzAsNCxbMCwxLCJcXGZib3h7SW5ncmVkaWVudH0iXSxbMywxLCJcXGZib3h7RWZmZWN0fSJdLFszLDAsIjEiXSxbMCwwLCJcXGZib3h7SW5ncmVkaWVudH0iXSxbMCwxLCJcXHRleHR7UG90aW9ufSIsMix7InN0eWxlIjp7ImJvZHkiOnsibmFtZSI6ImJhcnJlZCJ9fX1dLFsyLDEsIlxcdGV4dHtmb3J0aWZ5IGhlYWx0aH0iXSxbMywwXSxbMywyLCJcXHRleHR7Zm9ydGlmeSBoZWFsdGggcG90aW9ufSIsMCx7InN0eWxlIjp7ImJvZHkiOnsibmFtZSI6ImJhcnJlZCJ9fX1dLFs3LDQsIlxcdGV4dHtyZXN9IiwxLHsic2hvcnRlbiI6eyJzb3VyY2UiOjIwLCJ0YXJnZXQiOjIwfSwic3R5bGUiOnsiYm9keSI6eyJuYW1lIjoibm9uZSJ9LCJoZWFkIjp7Im5hbWUiOiJub25lIn19fV1d
      \begin{tikzcd}
        {\fbox{Ingredient}} &&& 1 \\
        {\fbox{Ingredient}} &&& {\fbox{Effect}}
        \arrow[""{name=0, anchor=center, inner sep=0}, "{\text{fortify health potion}}"{inner sep=.8ex}, "\shortmid"{marking}, from=1-1, to=1-4]
        \arrow[from=1-1, to=2-1]
        \arrow["{\text{fortify health}}", from=1-4, to=2-4]
        \arrow[""{name=1, anchor=center, inner sep=0}, "{\text{Potion}}"'{inner sep=.8ex}, "\shortmid"{marking}, from=2-1, to=2-4]
        \arrow["{\text{res}}"{description}, draw=none, from=0, to=1]
      \end{tikzcd}
    \end{equation*}
These are both examples of \emph{filter queries} discussed in \cite{lambert2025}. We are trying ultimately to compare the two restrictions, to see if the ingredients for the fortify health potion are in the various inventories. In the examples above, the composites 
    \begin{equation*}
      % https://q.uiver.app/#q=WzAsMyxbNiwwLCIxIl0sWzMsMCwiXFxmYm94e0luZ3JlZGllbnR9Il0sWzAsMCwiMSJdLFsxLDAsIlxcdGV4dHtmb3J0aWZ5IGhlYWx0aCBwb3Rpb259IiwwLHsic3R5bGUiOnsiYm9keSI6eyJuYW1lIjoiYmFycmVkIn19fV0sWzIsMSwiXFx0ZXh0e0FyY2FkaWEncyBJbnZlbnRvcnl9IiwwLHsic3R5bGUiOnsiYm9keSI6eyJuYW1lIjoiYmFycmVkIn19fV1d
      \begin{tikzcd}
        1 &&& {\fbox{Ingredient}} &&& 1
        \arrow["{\text{Arcadia's Inventory}}"{inner sep=.8ex}, "\shortmid"{marking}, from=1-1, to=1-4]
        \arrow["{\text{fortify health potion}}"{inner sep=.8ex}, "\shortmid"{marking}, from=1-4, to=1-7]
      \end{tikzcd}
    \end{equation*}
and 
    \begin{equation*}
      % https://q.uiver.app/#q=WzAsNCxbMCwxLCJcXGZib3h7QXBvdGhlY2FyeX0iXSxbMywxLCJcXGZib3h7SW5ncmVkaWVudH0iXSxbMCwwLCIxIl0sWzMsMCwiXFxmYm94e0luZ3JlZGllbnR9Il0sWzAsMSwiXFx0ZXh0e0ludmVudG9yeX0iLDIseyJzdHlsZSI6eyJib2R5Ijp7Im5hbWUiOiJiYXJyZWQifX19XSxbMiwwLCJcXHRleHR7QXJjYWRpYSdzfSIsMl0sWzIsMywiXFx0ZXh0e0FyY2FkaWEncyBJbnZlbnRvcnl9IiwwLHsic3R5bGUiOnsiYm9keSI6eyJuYW1lIjoiYmFycmVkIn19fV0sWzMsMV0sWzYsNCwiXFx0ZXh0e3Jlc30iLDEseyJzaG9ydGVuIjp7InNvdXJjZSI6MjAsInRhcmdldCI6MjB9LCJzdHlsZSI6eyJib2R5Ijp7Im5hbWUiOiJub25lIn0sImhlYWQiOnsibmFtZSI6Im5vbmUifX19XV0=
      \begin{tikzcd}
        1 &&& {\fbox{Ingredient}} \\
        {\fbox{Apothecary}} &&& {\fbox{Ingredient}}
        \arrow[""{name=0, anchor=center, inner sep=0}, "{\text{Arcadia's Inventory}}"{inner sep=.8ex}, "\shortmid"{marking}, from=1-1, to=1-4]
        \arrow["{\text{Arcadia's}}"', from=1-1, to=2-1]
        \arrow[from=1-4, to=2-4]
        \arrow[""{name=1, anchor=center, inner sep=0}, "{\text{Inventory}}"'{inner sep=.8ex}, "\shortmid"{marking}, from=2-1, to=2-4]
        \arrow["{\text{res}}"{description}, draw=none, from=0, to=1]
      \end{tikzcd}
    \end{equation*}
are natural to study but neither returns what we want. In the first case, the $\Rel$-composite will return whichever ingredient Arcadia's has on the list; in the second case, the $\Rel$-composite returns whichever merchants have an ingredient on the list. This is a result of the fact that existential quantification over the instance of the middle type is built into composition in $\Rel$. Neither of course is right: the first considers only one merchant; the second only whether a merchant has an ingredient on the list.

The better way is to view our shopping list as a subtype of the type of ingredients:
    \begin{equation*}
      % https://q.uiver.app/#q=WzAsMixbMCwxLCJcXGZib3h7SW5ncmVkaWVudH0iXSxbMCwwLCJcXGZib3h7XFx0ZXh0e2luZ3JlZGllbnQgZm9yIGZvcnRpZnkgaGVhbHRoIHBvdGlvbn19Il0sWzEsMCwiIiwwLHsic3R5bGUiOnsidGFpbCI6eyJuYW1lIjoibW9ubyJ9fX1dXQ==
      \begin{tikzcd}
        {\fbox{\text{ingredient for fortify health potion}}} \\
        {\fbox{Ingredient}}
        \arrow[tail, from=1-1, to=2-1]
      \end{tikzcd}
    \end{equation*}
This need not necessarily be monic. But monic arrows are perfectly well axiomatizable in our framework. Perhaps we use the notation to indicate we are thinking of it as a subtype obtained by comprehension \cite[\S 9.3]{jacobs1999}. Nor does it matter how the subtype comes about at this stage; it could have been obtained as a tabulator, or simply required as a primitive in the setup of the olog. Notice that this in any case is a genuine subtype that is not necessarily a global element (as was the case with all the examples in \cite{lambert2025}.

What matters is comparing this subtype with the inventories of all the merchants simultaneously and finding a way to quantify over those inventories. The individual inventories are implicit in the instancing of the Inventory relation, so we study the corner  
    \begin{equation*}
      % https://q.uiver.app/#q=WzAsMyxbMiwxLCJcXGZib3h7SW5ncmVkaWVudH0iXSxbMiwwLCJcXGZib3h7XFx0ZXh0e2luZ3JlZGllbnQgZm9yIGZvcnRpZnkgaGVhbHRoIHBvdGlvbn19Il0sWzAsMSwiXFxmYm94e05hbWV9Il0sWzEsMCwiIiwwLHsic3R5bGUiOnsidGFpbCI6eyJuYW1lIjoibW9ubyJ9fX1dLFsyLDAsIlxcdGV4dHtBcG90aGVjYXJ5fSIsMix7InN0eWxlIjp7ImJvZHkiOnsibmFtZSI6ImJhcnJlZCJ9fX1dXQ==
      \begin{tikzcd}
        && {\fbox{\text{ingredient for fortify health potion}}} \\
        {\fbox{Name}} && {\fbox{Ingredient}}
        \arrow[tail, from=1-3, to=2-3]
        \arrow["{\text{Apothecary}}"'{inner sep=.8ex}, "\shortmid"{marking}, from=2-1, to=2-3]
      \end{tikzcd}.
    \end{equation*}
It is tempting just to restrict with an identity morphism on the left. This results in a cartesian cell 
    \begin{equation*}
      % https://q.uiver.app/#q=WzAsNCxbNiwxLCJcXGZib3h7SW5ncmVkaWVudH0iXSxbNiwwLCJcXGZib3h7XFx0ZXh0e2luZ3JlZGllbnQgZm9yIGZvcnRpZnkgaGVhbHRoIHBvdGlvbn19Il0sWzAsMSwiXFxmYm94e05hbWV9Il0sWzAsMCwiXFxmYm94e05hbWV9Il0sWzEsMCwiIiwwLHsic3R5bGUiOnsidGFpbCI6eyJuYW1lIjoibW9ubyJ9fX1dLFsyLDAsIlxcdGV4dHtBcG90aGVjYXJ5fSIsMix7InN0eWxlIjp7ImJvZHkiOnsibmFtZSI6ImJhcnJlZCJ9fX1dLFszLDIsIiIsMix7ImxldmVsIjoyLCJzdHlsZSI6eyJoZWFkIjp7Im5hbWUiOiJub25lIn19fV0sWzMsMSwiXFx0ZXh0e01lcmNoYW50IHdpdGggZm9ydGlmeSBoZWFsdGggaW5ncmVkaWVudH0iLDAseyJzdHlsZSI6eyJib2R5Ijp7Im5hbWUiOiJiYXJyZWQifX19XSxbNyw1LCJcXHRleHR7cmVzfSIsMSx7InNob3J0ZW4iOnsic291cmNlIjoyMCwidGFyZ2V0IjoyMH0sInN0eWxlIjp7ImJvZHkiOnsibmFtZSI6Im5vbmUifSwiaGVhZCI6eyJuYW1lIjoibm9uZSJ9fX1dXQ==
      \begin{tikzcd}
        {\fbox{Name}} &&&&&& {\fbox{\text{ingredient for fortify health potion}}} \\
        {\fbox{Name}} &&&&&& {\fbox{Ingredient}}
        \arrow[""{name=0, anchor=center, inner sep=0}, "{\text{Merchant with fortify health ingredient}}"{inner sep=.8ex}, "\shortmid"{marking}, from=1-1, to=1-7]
        \arrow[equals, from=1-1, to=2-1]
        \arrow[tail, from=1-7, to=2-7]
        \arrow[""{name=1, anchor=center, inner sep=0}, "{\text{Apothecary}}"'{inner sep=.8ex}, "\shortmid"{marking}, from=2-1, to=2-7]
        \arrow["{\text{res}}"{description}, draw=none, from=0, to=1]
      \end{tikzcd}
    \end{equation*}
which filters the table for those rows having one of the two ingredients on the list:
    \begin{equation*}
      \begin{tabular}{| l | l | }
        \hline\multicolumn{2}{| c |}{\bf Apothecary}\\
        \hline {\bf Name }&{\bf Ingredient}\\
        \hline Arcadia's Cauldron & giant's toe \\
        \hline Arcadia's Cauldron & wheat \\
        \hline Elgrim's Elixers & giant's toe \\
        \hline Elgrim's Elixers & wheat \\
        \hline The White Phial & wheat \\
        \hline
      \end{tabular}.
    \end{equation*}
Notice that this gets rid of the unnecessary entry with the value ``void salts'' but it does not establish our query. Another operation with which our olog comes equipped is extension. In $\Rel$ this acts as an image, and formalizes a select query \cite{lambert2025}. In the present instance we might try an extension
    \begin{equation*}
      % https://q.uiver.app/#q=WzAsNCxbNiwxLCIxIl0sWzYsMCwiXFxmYm94e1xcdGV4dHtpbmdyZWRpZW50IGZvciBmb3J0aWZ5IGhlYWx0aCBwb3Rpb259fSJdLFswLDEsIlxcZmJveHtOYW1lfSJdLFswLDAsIlxcZmJveHtOYW1lfSJdLFsxLDAsIiEiXSxbMiwwLCJcXHRleHR7TWVyY2hhbnQgd2l0aCBmb3J0aWZ5IGhlYWx0aCBpbmdyZWRpZW50fSIsMix7InN0eWxlIjp7ImJvZHkiOnsibmFtZSI6ImJhcnJlZCJ9fX1dLFszLDIsIiIsMix7ImxldmVsIjoyLCJzdHlsZSI6eyJoZWFkIjp7Im5hbWUiOiJub25lIn19fV0sWzMsMSwiXFx0ZXh0e01lcmNoYW50IHdpdGggZm9ydGlmeSBoZWFsdGggaW5ncmVkaWVudH0iLDAseyJzdHlsZSI6eyJib2R5Ijp7Im5hbWUiOiJiYXJyZWQifX19XSxbNyw1LCJcXHRleHR7ZXh0fSIsMSx7InNob3J0ZW4iOnsic291cmNlIjoyMCwidGFyZ2V0IjoyMH0sInN0eWxlIjp7ImJvZHkiOnsibmFtZSI6Im5vbmUifSwiaGVhZCI6eyJuYW1lIjoibm9uZSJ9fX1dXQ==
      \begin{tikzcd}
        {\fbox{Name}} &&&&&& {\fbox{\text{ingredient for fortify health potion}}} \\
        {\fbox{Name}} &&&&&& 1
        \arrow[""{name=0, anchor=center, inner sep=0}, "{\text{Merchant with fortify health ingredient}}"{inner sep=.8ex}, "\shortmid"{marking}, from=1-1, to=1-7]
        \arrow[equals, from=1-1, to=2-1]
        \arrow["{!}", from=1-7, to=2-7]
        \arrow[""{name=1, anchor=center, inner sep=0}, "{\text{Merchant with fortify health ingredient}}"'{inner sep=.8ex}, "\shortmid"{marking}, from=2-1, to=2-7]
        \arrow["{\text{ext}}"{description}, draw=none, from=0, to=1]
      \end{tikzcd}
    \end{equation*}
which in the data instance is effectively just a list of the three merchants, since each has at least one ingredient. But if we went to the White Phial, we would be out of luck. This is owing to the fact that in $\Rel$ the extension is an image computed via \emph{existential quantification}. So, the composite query here of restriction along the subtype, followed by extension has performed an \emph{existential query}, returning each merchant whose inventory contains an ingredient on the list.

Now, this is clearly not the result we want, but it is not a dead-end, for it suggests a path forward. Extensions are left adjoints to restriction. Inasmuch as extension is thought of as existential quantification and restriction is thought of as substitution, we have an analogy of adjunctions 
    \begin{equation*}
      \text{ext}\dashv \text{res} \qquad\leftrightarrow \qquad \exists \dashv (-)^*.
    \end{equation*}
Of course in nice logical situations substitution also has a right adjoint, which is \emph{universal quantification}. This is of course exactly the kind of operation we are seeking to formalize. What we want then is to complete the adjoint analogy
    \begin{equation*}
      \text{ext}\dashv \text{res} \dashv \,\,??? \qquad\leftrightarrow \qquad \exists \dashv (-)^*\dashv \forall
    \end{equation*}
in the theater of ologs and of double categories generally. Something like this has been studied in the foundational paper \cite{shulman2008} under the guise of \emph{coextension}. It was mentioned there, but not studied in-depth due to being not as well behaved as the other two operations. And indeed this is not what we envision in this paper. For coextensions are had by asking for the equipment $\dbl{D}_1\to\dbl{D}_0\times\dbl{D}_0$ to be also a \emph{cofibration} which is too strong for what we have in mind. Rather, we take the approach of \emph{fibrational semantics} \cite{jacobs1999} and will ask merely that our equipments $\dbl{D}_1\to\dbl{D}_0\times\dbl{D}_0$ have right adjoints to restriction functors and that these satisfy a Beck-Chevalley condition. This is exactly to require \emph{fibered products}. We shall see in more detail how this recovers exactly the queries of interest below in \cref{section:universal-queries}. For this we first review quantification in more detail.

\section{Quantification and Modality}
\label{section:varieties-of-quantification}

It is well-known that quantification provides left and right adjoints to substitution \cite{lawvere2006,maclane1992}. For an ordinary set function $f\colon X\to Y$, there is the usual adjoint situation
  \begin{equation*}
    % https://q.uiver.app/#q=WzAsMixbMCwwLCJQWCJdLFsyLDAsIlBZIl0sWzAsMSwiXFxmb3JhbGxfZiIsMCx7Im9mZnNldCI6LTN9XSxbMCwxLCJcXGV4aXN0c19mIiwyLHsib2Zmc2V0IjozfV0sWzEsMCwiZl4qIiwxXV0=
    \begin{tikzcd}
      PX && PY
      \arrow["{\forall_f}", shift left=3, from=1-1, to=1-3]
      \arrow["{\exists_f}"', shift right=3, from=1-1, to=1-3]
      \arrow["{f^*}"{description}, from=1-3, to=1-1]
    \end{tikzcd}\qquad\qquad \exists_f\dashv f^*\dashv\forall_f
  \end{equation*}
where
  \begin{equation*}
    \exists_f(S) = \lbrace y\in Y\mid f(x) = y \text{ for some } x\in S\rbrace
  \end{equation*}
  \begin{equation*}
    \forall_f(S) = \lbrace y\in Y\mid \text{ for all } x\in X \text{ if } f(x) = y \text{ then } x\in S\rbrace.
  \end{equation*}
Ordinary quantification is the special case taking $f = \pi\colon X\times Y\to Y$, the usual projection to the second factor of the cartesian product. In this case, the quantified subsets are $R\subset X\times Y$, that is, binary relations. That $x$ and $y$ are $R$-related is written $R(x,y)$ or $x\leq_Ry$. Quantification is then:
  \begin{equation*}
    \exists_\pi(R) = \lbrace y\in Y\mid R(x,y) \text{ for some } x\in X\rbrace
  \end{equation*}
  \begin{equation*}
    \forall_\pi(R) = \lbrace y\in Y\mid R(x,y) \text{ for all } x\in X\rbrace.
  \end{equation*}
Typically, the left variable is bound. Note also in particular that if $Y=1$, we have closed formulas, that is, truth values.  

Now, a relation $R\colon X\proto Y$ might also be thought of as set of pairs of entities and observations or attribute values, or when $X= Y$ as a non-deterministic system and $x\leq_Rx'$ as representing a causal or temporal transition or statement of accessibility. Now, the set of all $x\in X$ for which some subsequent state or attribute value satisfies a proposition $S\subset Y$ is written 
  \begin{equation*} \label{eqn:diamond-definition}
    \diamondsuit S = \lbrace x\mid S(y) \text{ for some } x\leq_R y \rbrace.
  \end{equation*}
This is exactly the modal possibility operator familiar from the topological semantics given by the Alexandrov Topology. Likewise there is the set of those elements of $X$, all of whose subsequent states or attribute values satisfy $S\subset Y$, written 
  \begin{equation*}
    \Box S = \lbrace x\mid S(y) \text{ for all } x\leq_R y\rbrace.
  \end{equation*}
This is modal necessity in the same topological semantics. Thus, in the former case, an element of $\diamondsuit S$ is an element of $X$ for which $S$ is possible, in the sense that it is true of some later accessible state. In the latter case, an element of $\Box S$ is an element of $X$ for which $S$ is necessarily true, in the sense that any subsequent accessible state satisfies $S$. Note that this data of a relation $R$ and proposition $S\subset Y$ is exactly the set up of our olog:
  \begin{equation*}
      % https://q.uiver.app/#q=WzAsMyxbMiwxLCJcXGZib3h7SW5ncmVkaWVudH0iXSxbMiwwLCJcXGZib3h7XFx0ZXh0e2luZ3JlZGllbnQgZm9yIGZvcnRpZnkgaGVhbHRoIHBvdGlvbn19Il0sWzAsMSwiXFxmYm94e05hbWV9Il0sWzEsMCwiIiwwLHsic3R5bGUiOnsidGFpbCI6eyJuYW1lIjoibW9ubyJ9fX1dLFsyLDAsIlxcdGV4dHtBcG90aGVjYXJ5fSIsMix7InN0eWxlIjp7ImJvZHkiOnsibmFtZSI6ImJhcnJlZCJ9fX1dXQ==
      \begin{tikzcd}
        && {\fbox{\text{ingredient for fortify health potion}}} \\
        {\fbox{Name}} && {\fbox{Ingredient}}
        \arrow[tail, from=1-3, to=2-3]
        \arrow["{\text{Apothecary}}"'{inner sep=.8ex}, "\shortmid"{marking}, from=2-1, to=2-3]
      \end{tikzcd}.
    \end{equation*}
Recall too the table instancing the bottom relation:
  \begin{equation*}
    \begin{tabular}{| l | l | }
        \hline\multicolumn{2}{| c |}{\bf Apothecary}\\
        \hline {\bf Name }&{\bf Ingredient}\\
        \hline Arcadia's Cauldron & giant's toe \\
        \hline Arcadia's Cauldron & wheat \\
        \hline Arcadia's Cauldron & void salts \\
        \hline Elgrim's Elixers & giant's toe \\
        \hline Elgrim's Elixers & wheat \\
        \hline The White Phial & wheat \\
        \hline
    \end{tabular}.
  \end{equation*}
Our list $S$ again has just wheat and giant's toe. Directly from the operations above we have six queries summarized in \cref{fig:quant-queries}. In the figure, $R$ denotes the relation and $R^\dagger$ is the reverse relation. The variables $x$ and $y$ range over names and over ingredients, respectively. 
\begin{figure}[H]
\begin{center}
\renewcommand{\arraystretch}{1.5} % extra vertical padding

\newlength{\boxheight}
\setlength{\boxheight}{4cm} % adjust row height as needed

\begin{tabularx}{\linewidth}{>{\centering\arraybackslash}X >{\centering\arraybackslash}X}

% Row 1
\begin{minipage}[t][\boxheight][t]{\hsize}
\begin{tcolorbox}[title=Result of Existential Query, boxrule=0.5mm, sharp corners, height=\boxheight,colback=white,colframe=black,coltitle=white,colbacktitle=black]
  \[
  \begin{aligned}
    \exists_\pi(R) &= \lbrace y\mid R(x,y) \text{ for some } x\in X\rbrace \\
                  &= \lbrace \text{ingredients someone has}\rbrace \\
                  &= \lbrace\text{wheat}, \text{giant's toe},\text{void salts}\rbrace
  \end{aligned}
  \]
\end{tcolorbox}
\end{minipage}
&
\begin{minipage}[t][\boxheight][t]{\hsize}
\begin{tcolorbox}[title=Result of Universal Query, boxrule=0.5mm, sharp corners, height=\boxheight,colback=white,colframe=black,coltitle=white,colbacktitle=black]
  \[
  \begin{aligned}
    \forall_\pi(R) &= \lbrace y\mid R(x,y) \text{ for all } x\in X\rbrace \\
                  &= \lbrace \text{ingredients everyone has}\rbrace \\
                  &= \lbrace\text{wheat}\rbrace
  \end{aligned}
  \]
\end{tcolorbox}
\end{minipage}
\\[2ex]

% Row 2
\begin{minipage}[t][\boxheight][t]{\hsize}
\begin{tcolorbox}[title=Result of Reverse Existential Query, boxrule=0.5mm, sharp corners, height=\boxheight,colback=white,colframe=black,coltitle=white,colbacktitle=black]
  \[
  \begin{aligned}
    \exists_\pi(R^\dagger) &= \lbrace x \mid R(x,y) \text{ for some } y\in Y\rbrace \\
                  &= \lbrace \text{those with an $R$-ingredient}\rbrace \\
                  &= \lbrace\text{wheat}, \text{giant's toe},\text{void salts}\rbrace
  \end{aligned}
  \]
\end{tcolorbox}
\end{minipage}
&
\begin{minipage}[t][\boxheight][t]{\hsize}
\begin{tcolorbox}[title=Result of Reverse Universal Query, boxrule=0.5mm, sharp corners, height=\boxheight,colback=white,colframe=black,coltitle=white,colbacktitle=black]
  \[
  \begin{aligned}
    \forall_\pi(R^\dagger) &= \lbrace x\mid R(x,y) \text{ for all } y\in Y\rbrace \\
                  &= \lbrace \text{those with all $R$-ingredients}\rbrace \\
                  &= \lbrace\text{Arcadia's}\rbrace
  \end{aligned}
  \]
\end{tcolorbox}
\end{minipage}
\\[2ex]

% Row 3
\begin{minipage}[t][\boxheight][t]{\hsize}
\begin{tcolorbox}[title=Result of Possibility Query (Liveness), boxrule=0.5mm, sharp corners, height=\boxheight,colback=white,colframe=black,coltitle=white,colbacktitle=black]
  \[
  \begin{aligned}
  \diamondsuit S &=\lbrace x\mid S(y) \text{ for some } x\leq_R y \rbrace \\
  &= \lbrace x \mid x \text{ has an ingredient on } S \rbrace \\ &= \lbrace \text{Arcadia's}, \text{Elgrim's}, \text{Phial}\rbrace
  \end{aligned}
  \]
\end{tcolorbox}
\end{minipage}
&
\begin{minipage}[t][\boxheight][t]{\hsize}
\begin{tcolorbox}[title=Results of Necessity Query (Safety), boxrule=0.5mm, sharp corners, height=\boxheight,colback=white,colframe=black,coltitle=white,colbacktitle=black]
  \[
  \begin{aligned}
  \Box S &= \lbrace x\mid S(y) \text{ for all } x\leq_R y\rbrace\\ 
  &= \lbrace x \mid x \text{ has only ingredients on } S \rbrace \\ &= \lbrace \text{Elgrim's, Phial}\rbrace
  \end{aligned}
  \]
\end{tcolorbox}
\end{minipage}

\end{tabularx}
\end{center}
\caption{Results of six quantification queries.}
\label{fig:quant-queries}
\end{figure}

Notice some curiosities. The first is that $R$ is supposed to be given. Thus, the existential queries are just the select column queries returning either all the merchants or all the ingredients specified by $R$. These are already studied in \cite{lambert2025}. The universal queries are new and more interesting, but neither returns our list query since $S$ does not even enter into the construction. The modal queries as well are new. Although they seemingly superficial resembling the division query, they are actually quite different. The possibility query returns the same as the existential query only as an accident of $R$. If The White Phial was paired with something other than wheat, for example, the results would be different. Note as well that the necessity query is not the list query we are looking for. It returns the set of those merchants having only the ingredients on our list. This is a query with utility, but as far as our shopping list goes it does not matter whether the returned merchant has other items. In fact, suppose that one closer has our ingredients but also a lot of others. We would naturally prefer to go there. The possibility construction is very close to presenting a \emph{liveness property} whereas the necessity construction is essentially phrasing a \emph{safety property} \cite{lamport1977}, \cite{alpern1985}, \cite{alpern1987}. The former means \emph{something good eventually happens} (that is, eventually values are on the list $S$); and the latter means \emph{nothing bad happens} (i.e. values never leave the list $S$). Combined with negation (\cref{section:cocartesian-negation}) we have something that looks more explicitly like safety: an execution never enters the defined \emph{bad region} in that it always remains in its complement.

It will turn out that the list query is given by a pair of operations, one of which is universal quantification, applied in sequence. And so, to enrich the data schema to handle this, the goal now is to formulate all of the above operations purely double-categorically and axiomatize them in our generated ologs. Ultimately, we will allow the data-instancing and double-categorical functorial semantics to perform the desired queries.

\section{Existential Queries}
\label{section:existential-queries}

Before casting universal queries in double-categorical formalism, we look more closely at existential-type queries. The main consideration is the possibility query in \cref{fig:quant-queries}. In this section we work in the setting of a generic cartesian equipment $\dbl D$ and gradually add further assumptions as needed. The discussion of existential operations in terms of \emph{Kan quantification} (\cref{discussion:kan-quantification}) motivates the eventual adoption of right adjoints to substitution as that machinery capturing division.

\begin{construction}
  The set-up we are entertaining is that of having a subobject or subtype $\phi\colon s\rightarrowtail y$ thought of as specifying in the associated data some subcollection or sublist of items in the ambient list instancing $y$. We have already seen that in $\Rel$, merely composing has the effect of an existential query. This can be formalized $\dbl{D}$ by first viewing $s$ as an extension and then composition with $r$ on the left as in 
    \begin{equation*}
      % https://q.uiver.app/#q=WzAsNSxbMSwwLCJzIl0sWzEsMSwieSJdLFswLDEsIngiXSxbMiwwLCJzIl0sWzIsMSwiMSJdLFswLDEsIlxccGhpIiwyLHsic3R5bGUiOnsidGFpbCI6eyJuYW1lIjoibW9ubyJ9fX1dLFsyLDEsInIiLDIseyJzdHlsZSI6eyJib2R5Ijp7Im5hbWUiOiJiYXJyZWQifX19XSxbMCwzLCJcXGlkX1MiLDAseyJzdHlsZSI6eyJib2R5Ijp7Im5hbWUiOiJiYXJyZWQifX19XSxbMyw0LCIhIl0sWzEsNCwiXFxoYXRcXHBoaSIsMix7InN0eWxlIjp7ImJvZHkiOnsibmFtZSI6ImJhcnJlZCJ9fX1dLFs3LDksIlxcdGV4dHtleHR9IiwxLHsic2hvcnRlbiI6eyJzb3VyY2UiOjIwLCJ0YXJnZXQiOjIwfSwic3R5bGUiOnsiYm9keSI6eyJuYW1lIjoibm9uZSJ9LCJoZWFkIjp7Im5hbWUiOiJub25lIn19fV1d
      \begin{tikzcd}
        & s & s \\
        x & y & 1
        \arrow[""{name=0, anchor=center, inner sep=0}, "{\id_S}"{inner sep=.8ex}, "\shortmid"{marking}, from=1-2, to=1-3]
        \arrow["\phi"', tail, from=1-2, to=2-2]
        \arrow["{!}", from=1-3, to=2-3]
        \arrow["r"'{inner sep=.8ex}, "\shortmid"{marking}, from=2-1, to=2-2]
        \arrow[""{name=1, anchor=center, inner sep=0}, "{\hat\phi}"'{inner sep=.8ex}, "\shortmid"{marking}, from=2-2, to=2-3]
        \arrow["{\text{ext}}"{description}, draw=none, from=0, to=1]
      \end{tikzcd}.
    \end{equation*}
  Here we are using the notation $\hat\phi$ to stand for the codomain of the extension above which is a composite of the conjoint of $\phi$ with the companion of the unique $s\to 1$. This is essentially converting a subtype or term to a one-sided proposition in the logic and semantically it is reversing the operation of taking a tabulator. In $\Rel$, the composite of proarrows on the bottom is precisely modal possibility, $\diamondsuit_r\phi$. An equivalent $\Rel$-construction is to form a pullback along $\phi$ and then take an image along $s\to 1$. This is formalized in any cartesian equipment as 
    \begin{equation*}
      % https://q.uiver.app/#q=WzAsNCxbMCwwLCJ4Il0sWzEsMCwicyJdLFsxLDEsIngiXSxbMCwxLCJ4Il0sWzAsMSwiclxcb2RvdCBcXHBoaV4qIiwwLHsic3R5bGUiOnsiYm9keSI6eyJuYW1lIjoiYmFycmVkIn19fV0sWzEsMiwiXFxwaGkiLDAseyJzdHlsZSI6eyJ0YWlsIjp7Im5hbWUiOiJtb25vIn19fV0sWzAsMywiIiwyLHsibGV2ZWwiOjIsInN0eWxlIjp7ImhlYWQiOnsibmFtZSI6Im5vbmUifX19XSxbMywyLCJyIiwyLHsic3R5bGUiOnsiYm9keSI6eyJuYW1lIjoiYmFycmVkIn19fV0sWzQsNywiXFx0ZXh0e3Jlc30iLDEseyJzaG9ydGVuIjp7InNvdXJjZSI6MjAsInRhcmdldCI6MjB9LCJzdHlsZSI6eyJib2R5Ijp7Im5hbWUiOiJub25lIn0sImhlYWQiOnsibmFtZSI6Im5vbmUifX19XV0=
      \begin{tikzcd}
        x & s \\
        x & y
        \arrow[""{name=0, anchor=center, inner sep=0}, "{r\odot \phi^*}"{inner sep=.8ex}, "\shortmid"{marking}, from=1-1, to=1-2]
        \arrow[equals, from=1-1, to=2-1]
        \arrow["\phi", tail, from=1-2, to=2-2]
        \arrow[""{name=1, anchor=center, inner sep=0}, "r"'{inner sep=.8ex}, "\shortmid"{marking}, from=2-1, to=2-2]
        \arrow["{\text{res}}"{description}, draw=none, from=0, to=1]
      \end{tikzcd} \qquad\leadsto\qquad 
      % https://q.uiver.app/#q=WzAsNCxbMCwwLCJ4Il0sWzEsMCwicyJdLFsxLDEsIjEiXSxbMCwxLCJ4Il0sWzAsMSwiclxcb2RvdCBcXHBoaV4qIiwwLHsic3R5bGUiOnsiYm9keSI6eyJuYW1lIjoiYmFycmVkIn19fV0sWzEsMiwiISJdLFswLDMsIiIsMix7ImxldmVsIjoyLCJzdHlsZSI6eyJoZWFkIjp7Im5hbWUiOiJub25lIn19fV0sWzMsMiwiXFxkaWFtb25kc3VpdF9yIFxccGhpIiwyLHsic3R5bGUiOnsiYm9keSI6eyJuYW1lIjoiYmFycmVkIn19fV0sWzQsNywiXFx0ZXh0e2V4dH0iLDEseyJzaG9ydGVuIjp7InNvdXJjZSI6MjAsInRhcmdldCI6MjB9LCJzdHlsZSI6eyJib2R5Ijp7Im5hbWUiOiJub25lIn0sImhlYWQiOnsibmFtZSI6Im5vbmUifX19XV0=
      \begin{tikzcd}
        x & s \\
        x & 1
        \arrow[""{name=0, anchor=center, inner sep=0}, "{r\odot \phi^*}"{inner sep=.8ex}, "\shortmid"{marking}, from=1-1, to=1-2]
        \arrow[equals, from=1-1, to=2-1]
        \arrow["{!}", from=1-2, to=2-2]
        \arrow[""{name=1, anchor=center, inner sep=0}, "{\diamondsuit_r \phi}"'{inner sep=.8ex}, "\shortmid"{marking}, from=2-1, to=2-2]
        \arrow["{\text{ext}}"{description}, draw=none, from=0, to=1]
      \end{tikzcd}.
    \end{equation*}
  In $\Rel$, the top of the restriction on the right is exactly the set of pairs $(x,y)$ where $x\leq_Ry$ under $R$. The image then binds over $y$ via existential quantification. Now, in $\dbl{D}$ there is a globular comparison cell $\diamondsuit_r \phi \Rightarrow r\odot \hat\phi$ since $\diamondsuit_r \phi$ is constructed as an extension: 
    \begin{equation} \label{equation:globular-cell}
      % https://q.uiver.app/#q=WzAsOCxbMCwxLCJ4Il0sWzEsMSwicyJdLFsyLDEsInMiXSxbMiwyLCIxIl0sWzAsMiwieCJdLFsxLDIsInkiXSxbMCwwLCJ4Il0sWzIsMCwicyJdLFsxLDIsIlxcaWQiLDAseyJzdHlsZSI6eyJib2R5Ijp7Im5hbWUiOiJiYXJyZWQifX19XSxbMiwzLCIhIl0sWzAsNCwiMSIsMix7ImxldmVsIjoyLCJzdHlsZSI6eyJoZWFkIjp7Im5hbWUiOiJub25lIn19fV0sWzQsNSwiciIsMix7InN0eWxlIjp7ImJvZHkiOnsibmFtZSI6ImJhcnJlZCJ9fX1dLFs1LDMsIlxcaGF0XFxwaGkiLDIseyJzdHlsZSI6eyJib2R5Ijp7Im5hbWUiOiJiYXJyZWQifX19XSxbMSw1LCJcXHBoaSJdLFs2LDAsIiIsMCx7ImxldmVsIjoyLCJzdHlsZSI6eyJoZWFkIjp7Im5hbWUiOiJub25lIn19fV0sWzYsNywiclxcb2RvdCBcXHBoaV4qIiwwLHsic3R5bGUiOnsiYm9keSI6eyJuYW1lIjoiYmFycmVkIn19fV0sWzcsMl0sWzAsMSwiclxcb2RvdCBcXHBoaV4qIiwwLHsic3R5bGUiOnsiYm9keSI6eyJuYW1lIjoiYmFycmVkIn19fV0sWzgsMTIsIlxcdGV4dHtleHR9IiwxLHsic2hvcnRlbiI6eyJzb3VyY2UiOjIwLCJ0YXJnZXQiOjIwfSwic3R5bGUiOnsiYm9keSI6eyJuYW1lIjoibm9uZSJ9LCJoZWFkIjp7Im5hbWUiOiJub25lIn19fV0sWzE3LDExLCJcXHRleHR7cmVzfSIsMSx7InNob3J0ZW4iOnsic291cmNlIjoyMCwidGFyZ2V0IjoyMH0sInN0eWxlIjp7ImJvZHkiOnsibmFtZSI6Im5vbmUifSwiaGVhZCI6eyJuYW1lIjoibm9uZSJ9fX1dLFsxNSwxLCJcXGNvbmciLDEseyJzaG9ydGVuIjp7InNvdXJjZSI6MjB9LCJzdHlsZSI6eyJib2R5Ijp7Im5hbWUiOiJub25lIn0sImhlYWQiOnsibmFtZSI6Im5vbmUifX19XV0=
      \begin{tikzcd}
        x && s \\
        x & s & s \\
        x & y & 1
        \arrow[""{name=0, anchor=center, inner sep=0}, "{r\odot \phi^*}"{inner sep=.8ex}, "\shortmid"{marking}, from=1-1, to=1-3]
        \arrow[equals, from=1-1, to=2-1]
        \arrow[from=1-3, to=2-3]
        \arrow[""{name=1, anchor=center, inner sep=0}, "{r\odot \phi^*}"{inner sep=.8ex}, "\shortmid"{marking}, from=2-1, to=2-2]
        \arrow["1"', equals, from=2-1, to=3-1]
        \arrow[""{name=2, anchor=center, inner sep=0}, "\id"{inner sep=.8ex}, "\shortmid"{marking}, from=2-2, to=2-3]
        \arrow["\phi", from=2-2, to=3-2]
        \arrow["{!}", from=2-3, to=3-3]
        \arrow[""{name=3, anchor=center, inner sep=0}, "r"'{inner sep=.8ex}, "\shortmid"{marking}, from=3-1, to=3-2]
        \arrow[""{name=4, anchor=center, inner sep=0}, "{\hat\phi}"'{inner sep=.8ex}, "\shortmid"{marking}, from=3-2, to=3-3]
        \arrow["\cong"{description}, draw=none, from=0, to=2-2]
        \arrow["{\text{res}}"{description}, draw=none, from=1, to=3]
        \arrow["{\text{ext}}"{description}, draw=none, from=2, to=4]
      \end{tikzcd} \qquad = \qquad 
      % https://q.uiver.app/#q=WzAsNixbMCwwLCJ4Il0sWzEsMCwicyJdLFsxLDEsIjEiXSxbMCwxLCJ4Il0sWzAsMiwieCJdLFsxLDIsIjEiXSxbMCwxLCJyXFxvZG90IFxccGhpXioiLDAseyJzdHlsZSI6eyJib2R5Ijp7Im5hbWUiOiJiYXJyZWQifX19XSxbMSwyLCIhIl0sWzAsMywiIiwyLHsibGV2ZWwiOjIsInN0eWxlIjp7ImhlYWQiOnsibmFtZSI6Im5vbmUifX19XSxbMywyLCJcXGRpYW1vbmRzdWl0X3IgXFxwaGkiLDIseyJzdHlsZSI6eyJib2R5Ijp7Im5hbWUiOiJiYXJyZWQifX19XSxbMyw0LCIiLDIseyJsZXZlbCI6Miwic3R5bGUiOnsiaGVhZCI6eyJuYW1lIjoibm9uZSJ9fX1dLFs0LDUsInJcXG9kb3QgXFxoYXRcXHBoaSIsMix7InN0eWxlIjp7ImJvZHkiOnsibmFtZSI6ImJhcnJlZCJ9fX1dLFsyLDUsIiIsMix7ImxldmVsIjoyLCJzdHlsZSI6eyJoZWFkIjp7Im5hbWUiOiJub25lIn19fV0sWzksMTEsIlxcZXhpc3RzXFwsISIsMSx7ImxhYmVsX3Bvc2l0aW9uIjo2MCwic2hvcnRlbiI6eyJzb3VyY2UiOjIwLCJ0YXJnZXQiOjIwfSwic3R5bGUiOnsiYm9keSI6eyJuYW1lIjoibm9uZSJ9LCJoZWFkIjp7Im5hbWUiOiJub25lIn19fV0sWzYsOSwiXFx0ZXh0e2V4dH0iLDEseyJzaG9ydGVuIjp7InNvdXJjZSI6MjAsInRhcmdldCI6MjB9LCJzdHlsZSI6eyJib2R5Ijp7Im5hbWUiOiJub25lIn0sImhlYWQiOnsibmFtZSI6Im5vbmUifX19XV0=
      \begin{tikzcd}
        x & s \\
        x & 1 \\
        x & 1
        \arrow[""{name=0, anchor=center, inner sep=0}, "{r\odot \phi^*}"{inner sep=.8ex}, "\shortmid"{marking}, from=1-1, to=1-2]
        \arrow[equals, from=1-1, to=2-1]
        \arrow["{!}", from=1-2, to=2-2]
        \arrow[""{name=1, anchor=center, inner sep=0}, "{\diamondsuit_r \phi}"'{inner sep=.8ex}, "\shortmid"{marking}, from=2-1, to=2-2]
        \arrow[equals, from=2-1, to=3-1]
        \arrow[equals, from=2-2, to=3-2]
        \arrow[""{name=2, anchor=center, inner sep=0}, "{r\odot \hat\phi}"'{inner sep=.8ex}, "\shortmid"{marking}, from=3-1, to=3-2]
        \arrow["{\text{ext}}"{description}, draw=none, from=0, to=1]
        \arrow["{\exists\,!}"{description, pos=0.6}, draw=none, from=1, to=2]
      \end{tikzcd}.
    \end{equation}
  Now, it can be seen directly that this is an invertible cell when $\dbl{D}$ is $\Mat(\cat{V})$ for suitably structured monoidal $\cat V$. This covers relations and spans. Now, in fact they agree in general for any equipment.
\end{construction}

\begin{proposition} \label{prop:possibility-rewrite-rule-1}
  The globular cell $\diamondsuit_r \phi \Rightarrow r\odot \hat\phi$ induced in \cref{equation:globular-cell} is an isomorphism in any equipment. Accordingly, the two potential definitions of possibility agree.
\end{proposition}
\begin{proof}
  It suffices to show that the two sides of \cref{equation:globular-cell} compute the same extension. But this is essentially an exercise is rearranging the left-hand side via the properties of restrictions and extensions and their construction via companions and conjoints \cite[\S 4]{shulman2008}.
\end{proof}

So, although the constructions are thus essentially the same, we make a choice as to which is basic. We insist on starting with the subtype $\phi\colon s\to y$ consistent with the set-up of \cref{section-ologs}. And so in either case there are two operations to perform the query. There is a slight semantic advantage to the latter. For in the case of $\Span$, although the appropriate element $x\in X$ is implicit in the composite as $d(r) = x$, our preference is to keep the data as a triple $((x,a),r) = (x,a,r)$ as in the second way of computing the possibility operator. This way all the information is recorded in the apex of the span and the two legs could potentially be forgotten.

\begin{definition}[Possibility] \label{def:possibility-operator}
  Let $\phi\colon s\rightarrowtail y$ denote any monic arrow in the equipment $\dbl{D}$. The proposition \textbf{possibly} $\phi$ is then formed as an restriction followed by an extension:
    \begin{equation*}
      % https://q.uiver.app/#q=WzAsNCxbMCwwLCJ4Il0sWzEsMCwicyJdLFsxLDEsIngiXSxbMCwxLCJ4Il0sWzAsMSwiclxcb2RvdCBcXHBoaV4qIiwwLHsic3R5bGUiOnsiYm9keSI6eyJuYW1lIjoiYmFycmVkIn19fV0sWzEsMiwiXFxwaGkiLDAseyJzdHlsZSI6eyJ0YWlsIjp7Im5hbWUiOiJtb25vIn19fV0sWzAsMywiIiwyLHsibGV2ZWwiOjIsInN0eWxlIjp7ImhlYWQiOnsibmFtZSI6Im5vbmUifX19XSxbMywyLCJyIiwyLHsic3R5bGUiOnsiYm9keSI6eyJuYW1lIjoiYmFycmVkIn19fV0sWzQsNywiXFx0ZXh0e3Jlc30iLDEseyJzaG9ydGVuIjp7InNvdXJjZSI6MjAsInRhcmdldCI6MjB9LCJzdHlsZSI6eyJib2R5Ijp7Im5hbWUiOiJub25lIn0sImhlYWQiOnsibmFtZSI6Im5vbmUifX19XV0=
      \begin{tikzcd}
        x & s \\
        x & y
        \arrow[""{name=0, anchor=center, inner sep=0}, "{r\odot \phi^*}"{inner sep=.8ex}, "\shortmid"{marking}, from=1-1, to=1-2]
        \arrow[equals, from=1-1, to=2-1]
        \arrow["\phi", tail, from=1-2, to=2-2]
        \arrow[""{name=1, anchor=center, inner sep=0}, "r"'{inner sep=.8ex}, "\shortmid"{marking}, from=2-1, to=2-2]
        \arrow["{\text{res}}"{description}, draw=none, from=0, to=1]
      \end{tikzcd} \qquad\leadsto\qquad 
      % https://q.uiver.app/#q=WzAsNCxbMCwwLCJ4Il0sWzEsMCwicyJdLFsxLDEsIjEiXSxbMCwxLCJ4Il0sWzAsMSwiclxcb2RvdCBcXHBoaV4qIiwwLHsic3R5bGUiOnsiYm9keSI6eyJuYW1lIjoiYmFycmVkIn19fV0sWzEsMiwiISJdLFswLDMsIiIsMix7ImxldmVsIjoyLCJzdHlsZSI6eyJoZWFkIjp7Im5hbWUiOiJub25lIn19fV0sWzMsMiwiXFxkaWFtb25kc3VpdF9yIFxccGhpIiwyLHsic3R5bGUiOnsiYm9keSI6eyJuYW1lIjoiYmFycmVkIn19fV0sWzQsNywiXFx0ZXh0e2V4dH0iLDEseyJzaG9ydGVuIjp7InNvdXJjZSI6MjAsInRhcmdldCI6MjB9LCJzdHlsZSI6eyJib2R5Ijp7Im5hbWUiOiJub25lIn0sImhlYWQiOnsibmFtZSI6Im5vbmUifX19XV0=
      \begin{tikzcd}
        x & s \\
        x & 1
        \arrow[""{name=0, anchor=center, inner sep=0}, "{r\odot \phi^*}"{inner sep=.8ex}, "\shortmid"{marking}, from=1-1, to=1-2]
        \arrow[equals, from=1-1, to=2-1]
        \arrow["{!}", from=1-2, to=2-2]
        \arrow[""{name=1, anchor=center, inner sep=0}, "{\diamondsuit_r \phi}"'{inner sep=.8ex}, "\shortmid"{marking}, from=2-1, to=2-2]
        \arrow["{\text{ext}}"{description}, draw=none, from=0, to=1]
      \end{tikzcd}.
    \end{equation*}
  If $r$ is understood, write just `$\diamondsuit\phi$' for this proposition.
\end{definition}

Note that then in $\Rel$-semantics, this operation actually recovers all the existential queries in \cref{fig:quant-queries} by combining with restrictions or reverse operators; that is, each is either an image or a composition. But the content of \cref{prop:possibility-rewrite-rule-1} is that images and compositions have essentially the same effect in suitable cases in general equipments. We close the section now with a high-level description of the common setting embracing all of these cases. This leads to our abstract formulation of universal queries in \cref{section:universal-queries} below.

\begin{discussion}[Quantification Motivation] \label{discussion:kan-quantification}
  All the instances of existential quantification discussed above arise from the \emph{fibrational semantics} of the equipment $\dbl{D}$. We have the general picture:
    \begin{equation} \label{eqn:left-adjoint-to-restriction}
      \coprod_{f,g}\colon\dbl{D}(a,b)\rightleftarrows\dbl{D}(x,y)\colon (f,g)^*\qquad\qquad\coprod_{f,g}\dashv (f,g)^*
    \end{equation}
  for any arrows $f\colon a\to x$ and $g\colon b\to y$ of $\dbl{D}$. The hom categories are the fibers of source-target fibration $\langle\src,\tgt\rangle\colon\dbl{D}_1\to\dbl{D}_0\times\dbl{D}_0$ over the pairs of objects $(a,b)$ and $(x,y)$. And the coproduct denotes the left adjoint to restriction along $f$ and $g$. All instances of existential quantification in \cref{fig:quant-queries} are recoverable in this framework. On the one hand, for $\phi\colon s\to y$ we have the functors
    \begin{equation} \label{eqn:possibility-via-doctrines}
      % https://q.uiver.app/#q=WzAsMyxbMCwwLCJcXGRibHtEfSh4LHkpIl0sWzEsMCwiXFxkYmx7RH0oeCxzKSJdLFsyLDAsIlxcZGJse0R9KHgsMSkiXSxbMCwxLCIoMV9YLHMpXioiXSxbMSwyLCJcXGNvcHJvZF97MV94LCF9Il1d
      \begin{tikzcd}
        {\dbl{D}(x,y)} & {\dbl{D}(x,s)} & {\dbl{D}(x,1)}
        \arrow["{(x,\phi)^*}", from=1-1, to=1-2]
        \arrow["{\coprod_{x,!}}", from=1-2, to=1-3]
      \end{tikzcd}\qquad\qquad \coprod_{x,!}(x,s)^*(r) = \diamondsuit_r \phi
    \end{equation}
  as in \cref{def:possibility-operator}. On the other hand, ordinary existential quantification is just an application of a left adjoint without the subtyping and restriction:
    \begin{equation*}
      % https://q.uiver.app/#q=WzAsMixbMCwwLCJcXGRibHtEfSh4LHkpIl0sWzEsMCwiXFxkYmx7RH0oeCwxKSJdLFswLDEsIlxcY29wcm9kX3sxX1gsIX0iXV0=
      \begin{tikzcd}
        {\dbl{D}(x,y)} & {\dbl{D}(x,1)}
        \arrow["{\coprod_{x,!}}", from=1-1, to=1-2]
      \end{tikzcd} \qquad\qquad \coprod_{x,!}(r) =: \exists y . r
    \end{equation*}
  as has already been observed in the case of $\Rel$. A reverse operator thus covers the last existential query in \cref{fig:quant-queries}. Now, in each case, the left adjoints used are special cases of \emph{Kan quantification} \cite[\S 3]{lawvere2002}. These are also fragments of the bicategorical closed structure discussed in \cite[\S 5]{shulman2008}. That is, these are species of one-sided adjunctions 
    \begin{equation} \label{eqn:one-sided-left-adjoint-to-restriction}
      \coprod_{x,h}\colon\dbl{D}(x,z)\rightleftarrows\dbl{D}(x,w)\colon (x,h)^*\qquad\coprod_{x,h}\dashv (x,h)^*
    \end{equation}
  for various choices of $h\colon z\to w$. Owning to the nature of the computation of restrictions and extensions via companions and conjoints \cite[Theorem 4.1]{shulman2008}, the two functors in the adjunction immediately above are simply: \emph{postcompose with a companion or conjoint} as in 
    \begin{equation} \label{eqn:one-sided-left-adjoint-to-restriction2}
      -\odot h_!\colon\dbl{D}(x,z)\rightleftarrows\dbl{D}(x,w)\colon -\odot h^* \qquad -\odot h_! \dashv -\odot h^*.
    \end{equation}
  Retaining the coproduct notation for the left adjoint, the adjunction thus expresses an existential quantification introduction rule. That is, it presents a bijection between globular cells
    \begin{equation*}
      % https://q.uiver.app/#q=WzAsNCxbMCwwLCJYIl0sWzEsMCwiWiJdLFsxLDEsIlciXSxbMCwxLCJYIl0sWzAsMSwiXFxjb3Byb2Rfe1gsaH0obSkiLDAseyJzdHlsZSI6eyJib2R5Ijp7Im5hbWUiOiJiYXJyZWQifX19XSxbMSwyLCJoIl0sWzAsMywiIiwyLHsibGV2ZWwiOjIsInN0eWxlIjp7ImhlYWQiOnsibmFtZSI6Im5vbmUifX19XSxbMywyLCJuIiwyLHsic3R5bGUiOnsiYm9keSI6eyJuYW1lIjoiYmFycmVkIn19fV0sWzQsNywiIiwwLHsic2hvcnRlbiI6eyJzb3VyY2UiOjMwLCJ0YXJnZXQiOjMwfX1dXQ==
      \begin{tikzcd}
        x & z \\
        x & w
        \arrow[""{name=0, anchor=center, inner sep=0}, "{\coprod_{x,h}(m)}"{inner sep=.8ex}, "\shortmid"{marking}, from=1-1, to=1-2]
        \arrow[equals, from=1-1, to=2-1]
        \arrow["h", from=1-2, to=2-2]
        \arrow[""{name=1, anchor=center, inner sep=0}, "n"'{inner sep=.8ex}, "\shortmid"{marking}, from=2-1, to=2-2]
        \arrow[between={0.3}{0.7}, Rightarrow, from=0, to=1]
      \end{tikzcd} \quad\leftrightarrow\quad 
      % https://q.uiver.app/#q=WzAsNSxbMCwwLCJYIl0sWzIsMCwiWiJdLFsyLDEsIloiXSxbMCwxLCJYIl0sWzEsMSwiVyJdLFswLDEsIm0iLDAseyJzdHlsZSI6eyJib2R5Ijp7Im5hbWUiOiJiYXJyZWQifX19XSxbMSwyLCIiLDAseyJsZXZlbCI6Miwic3R5bGUiOnsiaGVhZCI6eyJuYW1lIjoibm9uZSJ9fX1dLFswLDMsIiIsMix7ImxldmVsIjoyLCJzdHlsZSI6eyJoZWFkIjp7Im5hbWUiOiJub25lIn19fV0sWzMsNCwibiIsMix7InN0eWxlIjp7ImJvZHkiOnsibmFtZSI6ImJhcnJlZCJ9fX1dLFs0LDIsImheKiIsMix7InN0eWxlIjp7ImJvZHkiOnsibmFtZSI6ImJhcnJlZCJ9fX1dLFs1LDQsIiIsMCx7InNob3J0ZW4iOnsic291cmNlIjo0MCwidGFyZ2V0IjoyMH19XV0=
      \begin{tikzcd}
        x && z \\
        x & w & z
        \arrow[""{name=0, anchor=center, inner sep=0}, "m"{inner sep=.8ex}, "\shortmid"{marking}, from=1-1, to=1-3]
        \arrow[equals, from=1-1, to=2-1]
        \arrow[equals, from=1-3, to=2-3]
        \arrow["n"'{inner sep=.8ex}, "\shortmid"{marking}, from=2-1, to=2-2]
        \arrow["{h^*}"'{inner sep=.8ex}, "\shortmid"{marking}, from=2-2, to=2-3]
        \arrow[between={0.4}{0.8}, Rightarrow, from=0, to=2-2]
      \end{tikzcd}.
    \end{equation*}
  Specializing to the locally posetal case and where $h$ is the unique $y\to 1$ as above and writing the traditional `$\exists$' for the left adjoint, the bijection expressing the universal property is thus precisely the usual existential two-way rule:
    \begin{equation*}
      % https://q.uiver.app/#q=WzAsNCxbMCwwLCJYIl0sWzEsMCwiWSJdLFsxLDEsIjEiXSxbMCwxLCJYIl0sWzAsMSwiXFxleGlzdHNfWShtKSIsMCx7InN0eWxlIjp7ImJvZHkiOnsibmFtZSI6ImJhcnJlZCJ9fX1dLFsxLDIsIiEiXSxbMCwzLCIiLDIseyJsZXZlbCI6Miwic3R5bGUiOnsiaGVhZCI6eyJuYW1lIjoibm9uZSJ9fX1dLFszLDIsIm4iLDIseyJzdHlsZSI6eyJib2R5Ijp7Im5hbWUiOiJiYXJyZWQifX19XSxbNCw3LCJcXGxlcSIsMSx7InNob3J0ZW4iOnsic291cmNlIjozMCwidGFyZ2V0IjozMH0sInN0eWxlIjp7ImJvZHkiOnsibmFtZSI6Im5vbmUifSwiaGVhZCI6eyJuYW1lIjoibm9uZSJ9fX1dXQ==
      \begin{tikzcd}
        x & y \\
        x & 1
        \arrow[""{name=0, anchor=center, inner sep=0}, "{\exists_y(m)}"{inner sep=.8ex}, "\shortmid"{marking}, from=1-1, to=1-2]
        \arrow[equals, from=1-1, to=2-1]
        \arrow["{!}", from=1-2, to=2-2]
        \arrow[""{name=1, anchor=center, inner sep=0}, "n"'{inner sep=.8ex}, "\shortmid"{marking}, from=2-1, to=2-2]
        \arrow["\leq"{description}, draw=none, from=0, to=1]
      \end{tikzcd} \quad\leftrightarrow\quad
      % https://q.uiver.app/#q=WzAsNSxbMCwwLCJ4Il0sWzIsMCwieSJdLFsyLDEsInkiXSxbMCwxLCJ4Il0sWzEsMSwiMSJdLFswLDEsIm0iLDAseyJzdHlsZSI6eyJib2R5Ijp7Im5hbWUiOiJiYXJyZWQifX19XSxbMSwyLCIiLDAseyJsZXZlbCI6Miwic3R5bGUiOnsiaGVhZCI6eyJuYW1lIjoibm9uZSJ9fX1dLFswLDMsIiIsMix7ImxldmVsIjoyLCJzdHlsZSI6eyJoZWFkIjp7Im5hbWUiOiJub25lIn19fV0sWzMsNCwibiIsMix7InN0eWxlIjp7ImJvZHkiOnsibmFtZSI6ImJhcnJlZCJ9fX1dLFs0LDIsInleKiIsMix7InN0eWxlIjp7ImJvZHkiOnsibmFtZSI6ImJhcnJlZCJ9fX1dLFs1LDQsIlxcbGVxIiwxLHsibGFiZWxfcG9zaXRpb24iOjYwLCJzaG9ydGVuIjp7InNvdXJjZSI6NDAsInRhcmdldCI6MjB9LCJzdHlsZSI6eyJib2R5Ijp7Im5hbWUiOiJub25lIn0sImhlYWQiOnsibmFtZSI6Im5vbmUifX19XV0=
      \begin{tikzcd}
        x && y \\
        x & 1 & y
        \arrow[""{name=0, anchor=center, inner sep=0}, "m"{inner sep=.8ex}, "\shortmid"{marking}, from=1-1, to=1-3]
        \arrow[equals, from=1-1, to=2-1]
        \arrow[equals, from=1-3, to=2-3]
        \arrow["n"'{inner sep=.8ex}, "\shortmid"{marking}, from=2-1, to=2-2]
        \arrow["{y^*}"'{inner sep=.8ex}, "\shortmid"{marking}, from=2-2, to=2-3]
        \arrow["\leq"{description, pos=0.6}, draw=none, from=0, to=2-2]
      \end{tikzcd}\qquad \qquad
      \inferrule{\exists_y(m)\leq n}{m\leq n(y^*)}
    \end{equation*}
  where we have written $y^*$ for the conjoint of the unique $y\to 1$. This existential two-way rule is the basis motivating the introduction of dependent products as right adjoints to substitution in order to capture the implicit universal quantification in relational division queries.
\end{discussion}

\section{Universal Queries}
\label{section:universal-queries}

Now we turn to universal quantification starting from fibrational semantics and Kan quantification as framed in \cref{discussion:kan-quantification}. That the restriction functor $(f,g)^*\colon \dbl{D}(x,y)\to\dbl{D}(a,b)$ has a right adjoint is the statement that there exists a functor 
  \begin{equation*}
    \prod_{f,g}\colon \dbl{D}(a,b)\to\dbl{D}(x,y)
  \end{equation*}
and natural isomorphisms
  \begin{equation*}
    \dbl{D}(a,b)(f_!mg^*,s)\cong \dbl{D}(x,y)(m,\prod_{f,g}(s))
  \end{equation*}
establishing natural families of bijections of globular cells 
    \begin{equation*}
      \inferrule{f_!mg^*\Rightarrow s}{m\Rightarrow \prod_{f,g}(s)}
    \end{equation*}
given in each direction by application of the appropriate functor and then composition with a unit or counit as the case may be. Notice again that this specializes to a bidirectional universal quantification rule. The elementary characterization of this right adjoint as a universal in this context is the following.

\begin{proposition}
  The restriction functor $(f,g)^*\colon \dbl{D}(x,y)\to\dbl{D}(a,b)$ has a right adjoint if, and only if, for each $s\colon a\proto b$ there is a specified proarrow $\forall_{f,g}(s)$ together with a globular cell $\epsilon_s$ of the form 
    \begin{equation*}
      % https://q.uiver.app/#q=WzAsNCxbMCwwLCJhIl0sWzAsMSwieCJdLFsyLDAsImIiXSxbMiwxLCJ5Il0sWzAsMSwiZiIsMl0sWzAsMiwiZl8hKFxcZm9yYWxsX3tmLGd9KHMpKWdeKiIsMCx7InN0eWxlIjp7ImJvZHkiOnsibmFtZSI6ImJhcnJlZCJ9fX1dLFsxLDMsIlxcZm9yYWxsX3tmLGd9KHMpIiwyLHsic3R5bGUiOnsiYm9keSI6eyJuYW1lIjoiYmFycmVkIn19fV0sWzIsMywiZyJdLFs1LDYsIlxcdGV4dHtyZXN9IiwxLHsic2hvcnRlbiI6eyJzb3VyY2UiOjIwLCJ0YXJnZXQiOjIwfSwic3R5bGUiOnsiYm9keSI6eyJuYW1lIjoibm9uZSJ9LCJoZWFkIjp7Im5hbWUiOiJub25lIn19fV1d
      \begin{tikzcd}
        a && b \\
        x && y
        \arrow[""{name=0, anchor=center, inner sep=0}, "{f_!(\forall_{f,g}(s))g^*}"{inner sep=.8ex}, "\shortmid"{marking}, from=1-1, to=1-3]
        \arrow["f"', from=1-1, to=2-1]
        \arrow["g", from=1-3, to=2-3]
        \arrow[""{name=1, anchor=center, inner sep=0}, "{\forall_{f,g}(s)}"'{inner sep=.8ex}, "\shortmid"{marking}, from=2-1, to=2-3]
        \arrow["{\text{res}}"{description}, draw=none, from=0, to=1]
      \end{tikzcd} \qquad \leadsto \qquad
      % https://q.uiver.app/#q=WzAsNCxbMCwwLCJhIl0sWzAsMSwiYSJdLFsyLDEsImIiXSxbMiwwLCJiIl0sWzAsMSwiIiwyLHsibGV2ZWwiOjIsInN0eWxlIjp7ImhlYWQiOnsibmFtZSI6Im5vbmUifX19XSxbMSwyLCJzIiwyLHsic3R5bGUiOnsiYm9keSI6eyJuYW1lIjoiYmFycmVkIn19fV0sWzMsMiwiIiwwLHsibGV2ZWwiOjIsInN0eWxlIjp7ImhlYWQiOnsibmFtZSI6Im5vbmUifX19XSxbMCwzLCJmXyEoXFxmb3JhbGxfe2YsZ30ocykpZ14qIiwwLHsic3R5bGUiOnsiYm9keSI6eyJuYW1lIjoiYmFycmVkIn19fV0sWzcsNSwiXFxlcHNpbG9uX3MiLDEseyJzaG9ydGVuIjp7InNvdXJjZSI6MjAsInRhcmdldCI6MjB9LCJzdHlsZSI6eyJib2R5Ijp7Im5hbWUiOiJub25lIn0sImhlYWQiOnsibmFtZSI6Im5vbmUifX19XV0=
      \begin{tikzcd}
        a && b \\
        a && b
        \arrow[""{name=0, anchor=center, inner sep=0}, "{f_!(\forall_{f,g}(s))g^*}"{inner sep=.8ex}, "\shortmid"{marking}, from=1-1, to=1-3]
        \arrow[equals, from=1-1, to=2-1]
        \arrow[equals, from=1-3, to=2-3]
        \arrow[""{name=1, anchor=center, inner sep=0}, "s"'{inner sep=.8ex}, "\shortmid"{marking}, from=2-1, to=2-3]
        \arrow["{\epsilon_s}"{description}, draw=none, from=0, to=1]
      \end{tikzcd}
    \end{equation*}
  such that for any other globular cell from a restriction along $f$ and $g$, say, of the form 
    \begin{equation*}
      % https://q.uiver.app/#q=WzAsNCxbMCwwLCJhIl0sWzAsMSwieCJdLFsyLDEsInkiXSxbMiwwLCJiIl0sWzAsMSwiZiIsMl0sWzEsMiwidyIsMix7InN0eWxlIjp7ImJvZHkiOnsibmFtZSI6ImJhcnJlZCJ9fX1dLFszLDIsImciXSxbMCwzLCJmXyF3Z14qIiwwLHsic3R5bGUiOnsiYm9keSI6eyJuYW1lIjoiYmFycmVkIn19fV0sWzcsNSwiXFx0ZXh0e3Jlc30iLDEseyJzaG9ydGVuIjp7InNvdXJjZSI6MjAsInRhcmdldCI6MjB9LCJzdHlsZSI6eyJib2R5Ijp7Im5hbWUiOiJub25lIn0sImhlYWQiOnsibmFtZSI6Im5vbmUifX19XV0=
      \begin{tikzcd}
        a && b \\
        x && y
        \arrow[""{name=0, anchor=center, inner sep=0}, "{f_!wg^*}"{inner sep=.8ex}, "\shortmid"{marking}, from=1-1, to=1-3]
        \arrow["f"', from=1-1, to=2-1]
        \arrow["g", from=1-3, to=2-3]
        \arrow[""{name=1, anchor=center, inner sep=0}, "w"'{inner sep=.8ex}, "\shortmid"{marking}, from=2-1, to=2-3]
        \arrow["{\text{res}}"{description}, draw=none, from=0, to=1]
      \end{tikzcd} \qquad\leadsto\qquad
      % https://q.uiver.app/#q=WzAsNCxbMCwwLCJhIl0sWzAsMSwiYSJdLFsyLDAsImIiXSxbMiwxLCJiIl0sWzAsMSwiIiwyLHsibGV2ZWwiOjIsInN0eWxlIjp7ImhlYWQiOnsibmFtZSI6Im5vbmUifX19XSxbMCwyLCJmXyF3Z14qIiwwLHsic3R5bGUiOnsiYm9keSI6eyJuYW1lIjoiYmFycmVkIn19fV0sWzEsMywicyIsMix7InN0eWxlIjp7ImJvZHkiOnsibmFtZSI6ImJhcnJlZCJ9fX1dLFsyLDMsIiIsMCx7ImxldmVsIjoyLCJzdHlsZSI6eyJoZWFkIjp7Im5hbWUiOiJub25lIn19fV0sWzUsNiwiXFx4aSIsMSx7InNob3J0ZW4iOnsic291cmNlIjoyMCwidGFyZ2V0IjoyMH0sInN0eWxlIjp7ImJvZHkiOnsibmFtZSI6Im5vbmUifSwiaGVhZCI6eyJuYW1lIjoibm9uZSJ9fX1dXQ==
      \begin{tikzcd}
        a && b \\
        a && b
        \arrow[""{name=0, anchor=center, inner sep=0}, "{f_!wg^*}"{inner sep=.8ex}, "\shortmid"{marking}, from=1-1, to=1-3]
        \arrow[equals, from=1-1, to=2-1]
        \arrow[equals, from=1-3, to=2-3]
        \arrow[""{name=1, anchor=center, inner sep=0}, "s"'{inner sep=.8ex}, "\shortmid"{marking}, from=2-1, to=2-3]
        \arrow["\xi"{description}, draw=none, from=0, to=1]
      \end{tikzcd}
    \end{equation*}
  there is a unique cell $\hat \xi$ for which $h$ factors through $\epsilon_s$ as in 
    \begin{equation*}
      % https://q.uiver.app/#q=WzAsNixbMCwwLCJhIl0sWzAsMiwiYSJdLFsyLDAsImIiXSxbMiwyLCJiIl0sWzAsMSwiYSJdLFsyLDEsImIiXSxbMCwyLCJmXyF3Z14qIiwwLHsic3R5bGUiOnsiYm9keSI6eyJuYW1lIjoiYmFycmVkIn19fV0sWzEsMywicyIsMix7InN0eWxlIjp7ImJvZHkiOnsibmFtZSI6ImJhcnJlZCJ9fX1dLFswLDQsIiIsMix7ImxldmVsIjoyLCJzdHlsZSI6eyJoZWFkIjp7Im5hbWUiOiJub25lIn19fV0sWzIsNSwiIiwwLHsibGV2ZWwiOjIsInN0eWxlIjp7ImhlYWQiOnsibmFtZSI6Im5vbmUifX19XSxbNCwxLCIiLDIseyJsZXZlbCI6Miwic3R5bGUiOnsiaGVhZCI6eyJuYW1lIjoibm9uZSJ9fX1dLFs1LDMsIiIsMCx7ImxldmVsIjoyLCJzdHlsZSI6eyJoZWFkIjp7Im5hbWUiOiJub25lIn19fV0sWzQsNSwiZl8hKFxcZm9yYWxsX3tmLGd9KHMpKWdeKiIsMCx7InN0eWxlIjp7ImJvZHkiOnsibmFtZSI6ImJhcnJlZCJ9fX1dLFs2LDEyLCJcXGhhdFxceGkiLDEseyJsYWJlbF9wb3NpdGlvbiI6NDAsInNob3J0ZW4iOnsic291cmNlIjoyMCwidGFyZ2V0IjoyMH0sInN0eWxlIjp7ImJvZHkiOnsibmFtZSI6Im5vbmUifSwiaGVhZCI6eyJuYW1lIjoibm9uZSJ9fX1dLFsxMiw3LCJcXGVwc2lsb25fcyIsMSx7InNob3J0ZW4iOnsic291cmNlIjoyMCwidGFyZ2V0IjoyMH0sInN0eWxlIjp7ImJvZHkiOnsibmFtZSI6Im5vbmUifSwiaGVhZCI6eyJuYW1lIjoibm9uZSJ9fX1dXQ==
      \begin{tikzcd}
        a && b \\
        a && b \\
        a && b
        \arrow[""{name=0, anchor=center, inner sep=0}, "{f_!wg^*}"{inner sep=.8ex}, "\shortmid"{marking}, from=1-1, to=1-3]
        \arrow[equals, from=1-1, to=2-1]
        \arrow[equals, from=1-3, to=2-3]
        \arrow[""{name=1, anchor=center, inner sep=0}, "{f_!(\forall_{f,g}(s))g^*}"{inner sep=.8ex}, "\shortmid"{marking}, from=2-1, to=2-3]
        \arrow[equals, from=2-1, to=3-1]
        \arrow[equals, from=2-3, to=3-3]
        \arrow[""{name=2, anchor=center, inner sep=0}, "s"'{inner sep=.8ex}, "\shortmid"{marking}, from=3-1, to=3-3]
        \arrow["{\hat\xi}"{description, pos=0.4}, draw=none, from=0, to=1]
        \arrow["{\epsilon_s}"{description}, draw=none, from=1, to=2]
      \end{tikzcd}  
      \qquad = \qquad
      % https://q.uiver.app/#q=WzAsNCxbMCwwLCJhIl0sWzAsMiwiYSJdLFsxLDAsImIiXSxbMSwyLCJiIl0sWzAsMSwiIiwyLHsibGV2ZWwiOjIsInN0eWxlIjp7ImhlYWQiOnsibmFtZSI6Im5vbmUifX19XSxbMCwyLCJmXyF3Z14qIiwwLHsic3R5bGUiOnsiYm9keSI6eyJuYW1lIjoiYmFycmVkIn19fV0sWzEsMywicyIsMix7InN0eWxlIjp7ImJvZHkiOnsibmFtZSI6ImJhcnJlZCJ9fX1dLFsyLDMsIiIsMCx7ImxldmVsIjoyLCJzdHlsZSI6eyJoZWFkIjp7Im5hbWUiOiJub25lIn19fV0sWzUsNiwiXFx4aSIsMSx7InNob3J0ZW4iOnsic291cmNlIjoyMCwidGFyZ2V0IjoyMH0sInN0eWxlIjp7ImJvZHkiOnsibmFtZSI6Im5vbmUifSwiaGVhZCI6eyJuYW1lIjoibm9uZSJ9fX1dXQ==
      \begin{tikzcd}
        a & b \\
        \\
        a & b
        \arrow[""{name=0, anchor=center, inner sep=0}, "{f_!wg^*}"{inner sep=.8ex}, "\shortmid"{marking}, from=1-1, to=1-2]
        \arrow[equals, from=1-1, to=3-1]
        \arrow[equals, from=1-2, to=3-2]
        \arrow[""{name=1, anchor=center, inner sep=0}, "s"'{inner sep=.8ex}, "\shortmid"{marking}, from=3-1, to=3-2]
        \arrow["\xi"{description}, draw=none, from=0, to=1]
      \end{tikzcd}
    \end{equation*}
\end{proposition}
\begin{proof}
  This is just the characterization of an adjunction in terms of the existence of a counit as a universal \cite[\S IV.1, Theorem 2(iv)]{maclane1998}. Note that naturality of $\epsilon_s$ in $s$ follows automatically.
\end{proof}

\begin{remark}
  There is a crucial conceptual and technical difference here in the fibrational semantics \emph{vis-\'a-vis} the prior case of existential quantification. This is that universal quantification is not neatly expressible via loose composition with a companion or conjoint; nor is there some obvious third type of proarrow $f_*$ associated to a morphism $f$ with its own accompanying base-change cells and universal properties that would do this. For although equipments $\dbl{D}$ involve both a fibration and opfibration $\langle \src,\tgt\rangle\colon \dbl{D}_1\to\dbl{D}_0\times\dbl{D}_0$, this same functor is not usually also a cofibration \cite{shulman2008}. This would the structure required to provide some kind of \emph{corestriction} as a right adjoint to restriction. Universal quantification thus must appear under some different guise. Our approach of course has been that of fibrational semantics: adjoints to substitution are the basic notion; that in one case they happen to be computed by companions or conjoints is a nice feature of equipments, but evidently does not carry over the others. Now, our axiomatized notion is the following, based on \cite[\S1.9]{jacobs1999}. This is a geometric/fibrational analogue of a \emph{quantification hyperdoctrine} \cite{lawvere2006}, \cite{lawvere1970} instantiated here for double categories. As defined in \cite{lambert2025}, a double olog is a small and finitely-presented `double category of relations' \cite{lambert2022}. This is in particular a locally posetal cartesian equipment.
\end{remark}

\begin{definition} \label{def:prod-double-olog}
  A $\prod$-\textbf{double olog} is a double olog where each substitution functor has a right adjoint satisfying Beck-Chevalley.
\end{definition}

\begin{remark}
  Of course Beck-Chevalley holds automatically if the mere existence of dependent products is assumed \emph{and} dependent coproducts are assumed to satisfy Beck-Chevalley. This is by the usual mate calculus \cite[\S 1.9]{jacobs1999}. In particular, a `double category of relations' with strong tabulators and dependent products is a $\prod$-double olog (\cref{lemma:beck-chevalley}). Certainly $\Rel$ has such adjoints since they are inherited from $\Set$ and computed in the manner described in the introduction. Beck-Chevalley follows or can be verified directly. Thus, given an appropriate additional \emph{reverse} operator on proarrows, we have an axiomatization of the two universal queries of \cref{fig:quant-queries} given by application of the appropriate right adjoint.
\end{remark}

Our main application of universal quantification is, however, the list-query described \cref{section:relational-division} where we look for those proprietors having all the items on our shopping list. This can be done by a composition of a restriction followed by a universal quantifier right adjoint. Start with a given proarrow $r\colon x\proto y$ and a subtype $\phi\colon s\to y$ representing our list. Consider then the composite functor 
  \begin{equation} \label{eqn:list-via-doctrines}
    % https://q.uiver.app/#q=WzAsMyxbMCwwLCJcXGRibHtEfSh4LHkpIl0sWzEsMCwiXFxkYmx7RH0oeCxzKSJdLFsyLDAsIlxcZGJse0R9KHgsMSkiXSxbMCwxLCIoeCxzKV4qIl0sWzEsMiwiXFxwcm9kX3t4LCF9Il1d
    \begin{tikzcd}
      {\dbl{D}(x,y)} & {\dbl{D}(x,s)} & {\dbl{D}(x,1)}
      \arrow["{(x,\phi)^*}", from=1-1, to=1-2]
      \arrow["{\prod_{x,!}}", from=1-2, to=1-3]
    \end{tikzcd}
  \end{equation}
assuming of course that such right adjoints exist. Note that this is a direct analogue of the composite \cref{eqn:possibility-via-doctrines} which produced the modal query retrieving exactly those proprietors with some item on the list $S$. In the case of relations, this has the following effect.

\begin{example} \label{example:recast-relational-division}
  Specializing \cref{eqn:list-via-doctrines} to $\Rel$, recall that the restriction $R\odot \phi^*$ is the set of those pairs $(x,y)$ such that $x\leq_R y$ and $y\in S$. Now, the map $X\times S\to X\times 1$ is up to bijection between codomains the same as the first projection $X\times S\to X$. So, using the computation of $\forall_{\pi_X}$ in \cref{section:varieties-of-quantification}, we have 
    \begin{align*}
      \forall_{\pi_X}(R\odot \phi^*) &= \lbrace x\in X\mid \text{ for all } (x',y)\in X\times S \, (\pi(x',y) = x \text{ implies } (x',y)\in R\odot \phi^*)\rbrace \\
      &= \lbrace x\in X\mid \text{ for all } (x,y)\in X\times S \, (x,y)\in R\odot \phi^*\rbrace \\ 
      &= \lbrace x\in X\mid \text{ for all } y \in S \, (x,y)\in R\odot \phi^*\rbrace \\
      &= \lbrace x\in X\mid \text{ for all } y \in S \, (x\leq_R y)\rbrace
    \end{align*}
  That is, the result of the application of the composite of functors to $R$ is the set of those $X$-elements $x$ for which everything in $S$ is $R$-related to $x$. This is precisely to say that if $S$ is a list of ingredients and $X$ is our list of proprietors, then the query returns the list of those who have everything on the list (but could have other items).
\end{example}

\begin{definition} \label{def:data-instance}
  A \emph{data instance} for a $\prod$-double olog $\dbl D$ is a double functor $I\colon \dbl D\to\Rel$ that is also a morphism of fibrations. In particular, $I$ preserves substitution and its left and right adjoints.
\end{definition}

\begin{remark}
  Although this construction thus achieves the desired effect and in this sense the objective of the paper too, it leaves a gap in the narrative inasmuch as it is not a modal query like the existential query relative to $S$ that it mimics. The explanation is the relative flexibility of the existential operator. That is, existential validity asks just that \emph{something} satisfy the quantified formula, relation or proposition. This is a relatively weak requirement. And so the result of the possibility query in \cref{fig:quant-queries} is just the existential query $\exists$ applied to $R\odot \phi^*$. Universal queries, by their nature surveying all entities under the scope of the quantifier, are more stringent. Thus, $\forall(R\odot \phi^*)$ and $\Box S$ produce dramatically different results: the first case that of the shops with all ingredients on the list; and second those with only ingredients on the list. That is, the first case quantifies over list items and checks against inventory; the second quantifies over inventory and checks against list items. The difficulty consists in the fact that what needs to be returned in each case is the list of proprietors. The built-in for $\forall$ as a right adjoint must match variances without the mixing this would demand in one or the other case. That is, the comprehension of $\forall$ must always match the codomain typing of operator as a right adjoint. The object $R\odot \phi^*$ must be formed to compare shops with out list items. But just, for example, switching the projection or reversing the relation results in the codomain typing of $\forall$ to $S$-objects which will return ingredients, not shops in the query. To recover this operator we turn to tabulators in \cref{section:modality}.
\end{remark}

\section{Beck-Chevalley Condition}
\label{section:Beck-Chevalley}

As in \cite[\S 5.2]{aleiferi2018}, the (internal) Beck-Chevalley condition for an equipment $\dbl D$ says that given any pullback square on the left, the composite on the right is an isomorphism:
  \begin{equation*}
    % https://q.uiver.app/#q=WzAsNCxbMCwwLCJ4Il0sWzEsMCwieSJdLFsxLDEsInciXSxbMCwxLCJ6Il0sWzAsMSwiZiJdLFsxLDIsImsiXSxbMCwzLCJoIiwyXSxbMywyLCJnIiwyXSxbMCwyLCIiLDEseyJzdHlsZSI6eyJuYW1lIjoiY29ybmVyIn19XV0=
    \begin{tikzcd}
      x & y \\
      z & w
      \arrow["f", from=1-1, to=1-2]
      \arrow["h"', from=1-1, to=2-1]
      \arrow["\lrcorner"{anchor=center, pos=0.125}, draw=none, from=1-1, to=2-2]
      \arrow["k", from=1-2, to=2-2]
      \arrow["g"', from=2-1, to=2-2]
    \end{tikzcd}\qquad\qquad\qquad
    % https://q.uiver.app/#q=WzAsMTIsWzEsMCwieCJdLFsyLDAsIngiXSxbMiwxLCJ5Il0sWzEsMSwieiJdLFsxLDIsInciXSxbMiwyLCJ3Il0sWzAsMCwieiJdLFswLDEsInoiXSxbMCwyLCJ6Il0sWzMsMCwieSJdLFszLDEsInkiXSxbMywyLCJ5Il0sWzAsMSwiXFxpZCIsMCx7InN0eWxlIjp7ImJvZHkiOnsibmFtZSI6ImJhcnJlZCJ9fX1dLFsxLDIsImYiLDJdLFswLDMsImgiXSxbMyw0LCJnIl0sWzQsNSwiXFxpZCIsMix7InN0eWxlIjp7ImJvZHkiOnsibmFtZSI6ImJhcnJlZCJ9fX1dLFsyLDUsImsiLDJdLFs2LDAsImheKiIsMCx7InN0eWxlIjp7ImJvZHkiOnsibmFtZSI6ImJhcnJlZCJ9fX1dLFs2LDcsIiIsMCx7ImxldmVsIjoyLCJzdHlsZSI6eyJoZWFkIjp7Im5hbWUiOiJub25lIn19fV0sWzcsMywiIiwwLHsic3R5bGUiOnsiYm9keSI6eyJuYW1lIjoiYmFycmVkIn19fV0sWzcsOCwiIiwwLHsibGV2ZWwiOjIsInN0eWxlIjp7ImhlYWQiOnsibmFtZSI6Im5vbmUifX19XSxbOCw0LCJnXyEiLDIseyJzdHlsZSI6eyJib2R5Ijp7Im5hbWUiOiJiYXJyZWQifX19XSxbMSw5LCJmXyEiLDAseyJzdHlsZSI6eyJib2R5Ijp7Im5hbWUiOiJiYXJyZWQifX19XSxbOSwxMCwiIiwwLHsibGV2ZWwiOjIsInN0eWxlIjp7ImhlYWQiOnsibmFtZSI6Im5vbmUifX19XSxbMiwxMCwiIiwwLHsic3R5bGUiOnsiYm9keSI6eyJuYW1lIjoiYmFycmVkIn19fV0sWzUsMTEsImteKiIsMix7InN0eWxlIjp7ImJvZHkiOnsibmFtZSI6ImJhcnJlZCJ9fX1dLFsxMCwxMSwiIiwwLHsibGV2ZWwiOjIsInN0eWxlIjp7ImhlYWQiOnsibmFtZSI6Im5vbmUifX19XSxbMTgsMjAsIlxcZXBzaWxvbiIsMSx7InNob3J0ZW4iOnsic291cmNlIjoyMCwidGFyZ2V0IjoyMH0sInN0eWxlIjp7ImJvZHkiOnsibmFtZSI6Im5vbmUifSwiaGVhZCI6eyJuYW1lIjoibm9uZSJ9fX1dLFsyMCwyMiwiXFxldGEiLDEseyJzaG9ydGVuIjp7InNvdXJjZSI6MjAsInRhcmdldCI6MjB9LCJzdHlsZSI6eyJib2R5Ijp7Im5hbWUiOiJub25lIn0sImhlYWQiOnsibmFtZSI6Im5vbmUifX19XSxbMjMsMjUsIlxccmhvIiwxLHsic2hvcnRlbiI6eyJzb3VyY2UiOjIwLCJ0YXJnZXQiOjIwfSwic3R5bGUiOnsiYm9keSI6eyJuYW1lIjoibm9uZSJ9LCJoZWFkIjp7Im5hbWUiOiJub25lIn19fV0sWzI1LDI2LCJcXGRlbHRhIiwxLHsic2hvcnRlbiI6eyJzb3VyY2UiOjIwLCJ0YXJnZXQiOjIwfSwic3R5bGUiOnsiYm9keSI6eyJuYW1lIjoibm9uZSJ9LCJoZWFkIjp7Im5hbWUiOiJub25lIn19fV1d
    \begin{tikzcd}
      z & x & x & y \\
      z & z & y & y \\
      z & w & w & y
      \arrow[""{name=0, anchor=center, inner sep=0}, "{h^*}"{inner sep=.8ex}, "\shortmid"{marking}, from=1-1, to=1-2]
      \arrow[equals, from=1-1, to=2-1]
      \arrow["\id"{inner sep=.8ex}, "\shortmid"{marking}, from=1-2, to=1-3]
      \arrow["h", from=1-2, to=2-2]
      \arrow[""{name=1, anchor=center, inner sep=0}, "{f_!}"{inner sep=.8ex}, "\shortmid"{marking}, from=1-3, to=1-4]
      \arrow["f"', from=1-3, to=2-3]
      \arrow[equals, from=1-4, to=2-4]
      \arrow[""{name=2, anchor=center, inner sep=0}, "\shortmid"{marking}, from=2-1, to=2-2]
      \arrow[equals, from=2-1, to=3-1]
      \arrow["g", from=2-2, to=3-2]
      \arrow[""{name=3, anchor=center, inner sep=0}, "\shortmid"{marking}, from=2-3, to=2-4]
      \arrow["k"', from=2-3, to=3-3]
      \arrow[equals, from=2-4, to=3-4]
      \arrow[""{name=4, anchor=center, inner sep=0}, "{g_!}"'{inner sep=.8ex}, "\shortmid"{marking}, from=3-1, to=3-2]
      \arrow["\id"'{inner sep=.8ex}, "\shortmid"{marking}, from=3-2, to=3-3]
      \arrow[""{name=5, anchor=center, inner sep=0}, "{k^*}"'{inner sep=.8ex}, "\shortmid"{marking}, from=3-3, to=3-4]
      \arrow["\epsilon"{description}, draw=none, from=0, to=2]
      \arrow["\rho"{description}, draw=none, from=1, to=3]
      \arrow["\eta"{description}, draw=none, from=2, to=4]
      \arrow["\delta"{description}, draw=none, from=3, to=5]
    \end{tikzcd}.
  \end{equation*}
The cells labeled with Greek letters are the canonical base change cells coming with companions and conjoints. The main result of this section is that internal Beck-Chevalley implies the usual (external) Beck-Chevalley condition for the left adjoints to restriction in the context of fibrations \cite[\S 1.9]{jacobs1999}. As an upshot, this implies that any `double category of relations' with \emph{strong tabulators} (\cref{def:strong-tabulator}) satisfies external Beck-Chevalley automatically.

\begin{lemma} \label{lemma:beck-chevalley}
  If $\dbl D$ is a cartesian equipment satisfying internal Beck-Chevalley, then the source-target fibration $\langle\src,\tgt\rangle\colon\dbl{D}_1\to\dbl{D}_0\times\dbl{D}_0$ satisfies the usual Beck-Chevalley condition for the left adjoints to restriction along pullback squares. In particular, any `double category of relations' with strong tabulators satisfies the usual external Beck-Chevalley condition.
\end{lemma}
\begin{proof}
  The statement of the result is that the canonical cell 
    \begin{equation*}
      \coprod_{f,g}(h,k)^*(m)\Rightarrow (u,v)^*\coprod_{p,q}(m)
    \end{equation*}
  arising from an appropriate pullback square in $\dbl D_0\times\dbl D_0$ is an iso. But this is just a matter of the definition of the construction of restrictions and extensions using companions and conjoints and the assumed internal Beck-Chevalley condition summarized above. The last statement is a consequence of \cite[Lemma 5.2.3]{aleiferi2018} which shows that any unit-pure equipment with strong tabulators has pullback satisfying internal Beck-Chevalley. The key assumption for this argument is \emph{unit-purity}. But this is shown in \cite[Lemma 4.1.9]{hoshino2025} always to hold in any double category where every object is discrete.
\end{proof}

\begin{remark}[Optimization Rewrite Rule] \label{remark:Beck-Chevalley-Optimization}
  External Beck-Chevalley as above is a rewrite rule for compound queries. This is a \emph{filter-pushdown} optimization rule. The extension first on the right above acts as an abstraction/collapse or extraction of metadata before computing the filter/restriction. The rule says that this can more efficiently be done by \emph{filtering first} to restrict the number of rows in the table before the collapse. This is a standard and highly effective optimization in query planning in relational databases which is thus available in any double olog with strong tabulators by \cref{lemma:beck-chevalley}.
\end{remark}

\begin{example} \label{example:beck-chevalley}
  This rule is illustrated by the following. Start with two tables: one assigning to each individual his or her clan membership and another viewing each clan with its factional loyalty:
    \begin{equation*}
      \begin{tabular}{| l | l | }
          \hline\multicolumn{2}{| c |}{\bf Membership}\\
          \hline {\bf Individual }&{\bf Clan}\\
          \hline  Asgeir & Snow-Shod  \\
          \hline  Eorlund & Gray-Mane  \\
          \hline  Galmar & Stone-Fist  \\
          \hline  Idgrod & Battle-Born  \\
          \hline  Nilsine & Shatter-Shield \\
          \hline  Torbjorn & Shatter-Shield  \\
          \hline  Vignar & Gray-Mane  \\
          \hline
      \end{tabular} \qquad\qquad
      \begin{tabular}{| l | l | }
          \hline\multicolumn{2}{| c |}{\bf Loyalty}\\
          \hline {\bf Clan }&{\bf Faction}\\
          \hline  Battle-Born & Imperial  \\
          \hline  Gray-Mane & Stormcloak  \\
          \hline  Shatter-Shield & Imperial  \\
          \hline  Snow-Shod & Stormcloak  \\
          \hline  Stone-Fist & Stormcloack \\
          \hline
      \end{tabular}.
    \end{equation*}
  Suppose that the Membership table is a relation extracted perhaps from some larger dataset of various affiliations of relevant individuals; suppose that the Loyalty table is functional metadata. There are then evidently two schematized morphisms in the olog; one is a proarrow and one is an ordinary arrow:
    \begin{equation*}
      % https://q.uiver.app/#q=WzAsMixbMCwwLCJcXGZib3h7SW5kaXZpZHVhbH0iXSxbMiwwLCJcXGZib3h7Q2xhbn0iXSxbMCwxLCJcXHRleHR7TWVtYmVyc2hpcH0iLDAseyJzdHlsZSI6eyJib2R5Ijp7Im5hbWUiOiJiYXJyZWQifX19XV0=
      \begin{tikzcd}
        {\fbox{Individual}} && {\fbox{Clan}}
        \arrow["{\text{Membership}}"{inner sep=.8ex}, "\shortmid"{marking}, from=1-1, to=1-3]
      \end{tikzcd}\qquad\qquad
      % https://q.uiver.app/#q=WzAsMixbMCwwLCJcXGZib3h7Q2xhbn0iXSxbMiwwLCJcXGZib3h7RmFjdGlvbn0iXSxbMCwxLCJcXHRleHR7TG95YWx0eX0iXV0=
      \begin{tikzcd}
        {\fbox{Clan}} && {\fbox{Faction}}
        \arrow["{\text{Loyalty}}", from=1-1, to=1-3]
      \end{tikzcd}.
    \end{equation*}
  Now, to find the Empire loyalists, we can posit a global element and, given sufficient structure in the olog, form a pullback in the base of the source-target fibration:
    \begin{equation*}
      % https://q.uiver.app/#q=WzAsNCxbMCwwLCJcXGZib3h7SW5kaXZpZHVhbH0iXSxbMiwwLCJcXGZib3h7SW5kaXZpZHVhbH0iXSxbMiwxLCJcXGZib3h7SW5kaXZpZHVhbH0iXSxbMCwxLCJcXGZib3h7SW5kaXZpZHVhbH0iXSxbMCwxLCIxIl0sWzEsMiwiMSJdLFswLDMsIjEiLDJdLFszLDIsIjEiLDJdLFswLDcsIiIsMix7ImxldmVsIjoxLCJzdHlsZSI6eyJuYW1lIjoiY29ybmVyIn19XV0=
      \begin{tikzcd}
        {\fbox{Individual}} && {\fbox{Individual}} \\
        {\fbox{Individual}} && {\fbox{Individual}}
        \arrow["1", from=1-1, to=1-3]
        \arrow["1"', from=1-1, to=2-1]
        \arrow["1", from=1-3, to=2-3]
        \arrow[""{name=0, anchor=center, inner sep=0}, "1"', from=2-1, to=2-3]
        \arrow["\lrcorner"{anchor=center, pos=0.125}, draw=none, from=1-1, to=0]
      \end{tikzcd}\qquad\qquad
      % https://q.uiver.app/#q=WzAsNCxbMiwwLCJcXGZib3h7Q2xhbn0iXSxbMiwxLCJcXGZib3h7RmFjdGlvbn0iXSxbMCwxLCIxIl0sWzAsMCwiXFxmYm94e0ltcGVyaWFsIENsYW59Il0sWzAsMSwiXFx0ZXh0e0xveWFsdHl9Il0sWzIsMSwiXFx0ZXh0e0ltcGVyaWFsfSIsMl0sWzMsMiwiISIsMl0sWzMsMCwiXFx0ZXh0e2lzIGF9Il0sWzMsNSwiIiwwLHsibGV2ZWwiOjEsInN0eWxlIjp7Im5hbWUiOiJjb3JuZXIifX1dXQ==
      \begin{tikzcd}
        {\fbox{Imperial Clan}} && {\fbox{Clan}} \\
        1 && {\fbox{Faction}}
        \arrow["{\text{is a}}", from=1-1, to=1-3]
        \arrow["{!}"', from=1-1, to=2-1]
        \arrow["{\text{Loyalty}}", from=1-3, to=2-3]
        \arrow[""{name=0, anchor=center, inner sep=0}, "{\text{Imperial}}"', from=2-1, to=2-3]
        \arrow["\lrcorner"{anchor=center, pos=0.125}, draw=none, from=1-1, to=0]
      \end{tikzcd}
    \end{equation*}
  Now, there are two ways to do the Loyalist query using Beck-Chevalley. The naive way to start is to abstract Membership first along $(1,\text{Loyalty})$ and then filter for the Imperials as in 
    \begin{equation*}
      % https://q.uiver.app/#q=WzAsNCxbMCwwLCJcXGZib3h7SW5kaXZpZHVhbH0iXSxbMiwwLCJcXGZib3h7Q2xhbn0iXSxbMiwxLCJcXGZib3h7RmFjdGlvbn0iXSxbMCwxLCJcXGZib3h7SW5kaXZpZHVhbH0iXSxbMCwxLCJcXHRleHR7TWVtYmVyc2hpcH0iLDAseyJzdHlsZSI6eyJib2R5Ijp7Im5hbWUiOiJiYXJyZWQifX19XSxbMSwyLCJcXHRleHR7TG95YWx0eX0iXSxbMCwzLCIiLDAseyJsZXZlbCI6Miwic3R5bGUiOnsiaGVhZCI6eyJuYW1lIjoibm9uZSJ9fX1dLFszLDIsIlxcdGV4dHtMb3lhbHR5fSIsMix7InN0eWxlIjp7ImJvZHkiOnsibmFtZSI6ImJhcnJlZCJ9fX1dLFs0LDcsIlxcbWF0aHJte2V4dH0iLDEseyJzaG9ydGVuIjp7InNvdXJjZSI6MjAsInRhcmdldCI6MjB9LCJzdHlsZSI6eyJib2R5Ijp7Im5hbWUiOiJub25lIn0sImhlYWQiOnsibmFtZSI6Im5vbmUifX19XV0=
      \begin{tikzcd}
        {\fbox{Individual}} && {\fbox{Clan}} \\
        {\fbox{Individual}} && {\fbox{Faction}}
        \arrow[""{name=0, anchor=center, inner sep=0}, "{\text{Membership}}"{inner sep=.8ex}, "\shortmid"{marking}, from=1-1, to=1-3]
        \arrow[equals, from=1-1, to=2-1]
        \arrow["{\text{Loyalty}}", from=1-3, to=2-3]
        \arrow[""{name=1, anchor=center, inner sep=0}, "{\text{Loyalty}}"'{inner sep=.8ex}, "\shortmid"{marking}, from=2-1, to=2-3]
        \arrow["{\mathrm{ext}}"{description}, draw=none, from=0, to=1]
      \end{tikzcd}\qquad\qquad
      % https://q.uiver.app/#q=WzAsNCxbMCwwLCJcXGZib3h7SW5kaXZpZHVhbH0iXSxbMiwwLCIxIl0sWzIsMSwiXFxmYm94e0ZhY3Rpb259Il0sWzAsMSwiXFxmYm94e0luZGl2aWR1YWx9Il0sWzAsMSwiXFx0ZXh0e0ltcGVyaWFsIExveWFsaXN0c30iLDAseyJzdHlsZSI6eyJib2R5Ijp7Im5hbWUiOiJiYXJyZWQifX19XSxbMSwyLCJcXHRleHR7SW1wZXJpYWx9Il0sWzAsMywiIiwwLHsibGV2ZWwiOjIsInN0eWxlIjp7ImhlYWQiOnsibmFtZSI6Im5vbmUifX19XSxbMywyLCJcXHRleHR7TG95YWx0eX0iLDIseyJzdHlsZSI6eyJib2R5Ijp7Im5hbWUiOiJiYXJyZWQifX19XSxbNCw3LCJcXG1hdGhybXtyZXN9IiwxLHsic2hvcnRlbiI6eyJzb3VyY2UiOjIwLCJ0YXJnZXQiOjIwfSwic3R5bGUiOnsiYm9keSI6eyJuYW1lIjoibm9uZSJ9LCJoZWFkIjp7Im5hbWUiOiJub25lIn19fV1d
      \begin{tikzcd}
        {\fbox{Individual}} && 1 \\
        {\fbox{Individual}} && {\fbox{Faction}}
        \arrow[""{name=0, anchor=center, inner sep=0}, "{\text{Imperial Loyalists}}"{inner sep=.8ex}, "\shortmid"{marking}, from=1-1, to=1-3]
        \arrow[equals, from=1-1, to=2-1]
        \arrow["{\text{Imperial}}", from=1-3, to=2-3]
        \arrow[""{name=1, anchor=center, inner sep=0}, "{\text{Loyalty}}"'{inner sep=.8ex}, "\shortmid"{marking}, from=2-1, to=2-3]
        \arrow["{\mathrm{res}}"{description}, draw=none, from=0, to=1]
      \end{tikzcd}.
    \end{equation*}
  This has the resulting tables
    \begin{equation*}
      \begin{tabular}{| l | l | }
          \hline\multicolumn{2}{| c |}{\bf Loyalty}\\
          \hline {\bf Individual }&{\bf Faction}\\
          \hline  Asgeir & Stormcloak  \\
          \hline  Eorlund & Stormcloak  \\
          \hline  Galmar & Stormcloak  \\
          \hline  Idgrod & Imperial  \\
          \hline  Nilsine & Imperial \\
          \hline  Torbjorn & Imperial  \\
          \hline  Vignar & Stormcloack  \\
          \hline
      \end{tabular} \qquad\qquad
      \begin{tabular}{| l | }
          \hline\multicolumn{1}{| c |}{\bf Imperial Loyalists}\\
          \hline  Idgrod \\
          \hline  Nilsine\\
          \hline  Torbjorn \\
          \hline
      \end{tabular}.
    \end{equation*} 
  Basically, we looked to see which faction each clan supports, associated each individual to the appropriate faction and then took just the names. This is possibly the straightforward and naive way one might proceed at first. On the other hand, we can do a filter on the Membership relation by $(1,\text{is a})$ first and then select individuals:
    \begin{equation*}
      % https://q.uiver.app/#q=WzAsNCxbMCwxLCJcXGZib3h7SW5kaXZpZHVhbH0iXSxbMiwxLCJcXGZib3h7Q2xhbn0iXSxbMCwwLCJcXGZib3h7SW5kaXZpZHVhbH0iXSxbMiwwLCJcXGZib3h7SW1wZXJpYWwgQ2xhbn0iXSxbMCwxLCJcXHRleHR7TWVtYmVyc2hpcH0iLDIseyJzdHlsZSI6eyJib2R5Ijp7Im5hbWUiOiJiYXJyZWQifX19XSxbMiwwLCIiLDAseyJsZXZlbCI6Miwic3R5bGUiOnsiaGVhZCI6eyJuYW1lIjoibm9uZSJ9fX1dLFszLDEsIlxcdGV4dHtpcyBhfSJdLFsyLDMsIlxcdGV4dHtNZW1iZXJzaGlwfSIsMCx7InN0eWxlIjp7ImJvZHkiOnsibmFtZSI6ImJhcnJlZCJ9fX1dLFs3LDQsIlxcbWF0aHJte3Jlc30iLDEseyJzaG9ydGVuIjp7InNvdXJjZSI6MjAsInRhcmdldCI6MjB9LCJzdHlsZSI6eyJib2R5Ijp7Im5hbWUiOiJub25lIn0sImhlYWQiOnsibmFtZSI6Im5vbmUifX19XV0=
      \begin{tikzcd}
        {\fbox{Individual}} && {\fbox{Imperial Clan}} \\
        {\fbox{Individual}} && {\fbox{Clan}}
        \arrow[""{name=0, anchor=center, inner sep=0}, "{\text{Membership}}"{inner sep=.8ex}, "\shortmid"{marking}, from=1-1, to=1-3]
        \arrow[equals, from=1-1, to=2-1]
        \arrow["{\text{is a}}", from=1-3, to=2-3]
        \arrow[""{name=1, anchor=center, inner sep=0}, "{\text{Membership}}"'{inner sep=.8ex}, "\shortmid"{marking}, from=2-1, to=2-3]
        \arrow["{\mathrm{res}}"{description}, draw=none, from=0, to=1]
      \end{tikzcd}\qquad\qquad
      % https://q.uiver.app/#q=WzAsNCxbMCwwLCJcXGZib3h7SW5kaXZpZHVhbH0iXSxbMiwwLCJcXGZib3h7SW1wZXJpYWwgQ2xhbn0iXSxbMiwxLCIxIl0sWzAsMSwiXFxmYm94e0luZGl2aWR1YWx9Il0sWzAsMSwiXFx0ZXh0e01lbWJlcnNoaXB9IiwwLHsic3R5bGUiOnsiYm9keSI6eyJuYW1lIjoiYmFycmVkIn19fV0sWzEsMiwiISJdLFswLDMsIiIsMSx7ImxldmVsIjoyLCJzdHlsZSI6eyJoZWFkIjp7Im5hbWUiOiJub25lIn19fV0sWzMsMiwiXFx0ZXh0e0ltcGVyaWFsIExveWFsaXN0c30iLDIseyJzdHlsZSI6eyJib2R5Ijp7Im5hbWUiOiJiYXJyZWQifX19XSxbNCw3LCJcXG1hdGhybXtleHR9IiwxLHsic2hvcnRlbiI6eyJzb3VyY2UiOjIwLCJ0YXJnZXQiOjIwfSwic3R5bGUiOnsiYm9keSI6eyJuYW1lIjoibm9uZSJ9LCJoZWFkIjp7Im5hbWUiOiJub25lIn19fV1d
      \begin{tikzcd}
        {\fbox{Individual}} && {\fbox{Imperial Clan}} \\
        {\fbox{Individual}} && 1
        \arrow[""{name=0, anchor=center, inner sep=0}, "{\text{Membership}}"{inner sep=.8ex}, "\shortmid"{marking}, from=1-1, to=1-3]
        \arrow[equals, from=1-1, to=2-1]
        \arrow["{!}", from=1-3, to=2-3]
        \arrow[""{name=1, anchor=center, inner sep=0}, "{\text{Imperial Loyalists}}"'{inner sep=.8ex}, "\shortmid"{marking}, from=2-1, to=2-3]
        \arrow["{\mathrm{ext}}"{description}, draw=none, from=0, to=1]
      \end{tikzcd}.
    \end{equation*}
  This is to filter for individuals with their Imperial-supporting clan and then to project to the individuals:
    \begin{equation*}
      \begin{tabular}{| l | l | }
          \hline\multicolumn{2}{| c |}{\bf Membership}\\
          \hline {\bf Individual }&{\bf Imperial Clan}\\
          \hline  Idgrod & Battle-Born  \\
          \hline  Nilsine & Shatter-Shield \\
          \hline  Torbjorn & Shatter-Shield  \\
          \hline
      \end{tabular} \qquad\qquad
      \begin{tabular}{| l | }
          \hline\multicolumn{1}{| c |}{\bf Imperial Loyalists}\\
          \hline  Idgrod \\
          \hline  Nilsine\\
          \hline  Torbjorn \\
          \hline
      \end{tabular}.
    \end{equation*}
  Notice that the latter ends up with a much smaller intermediate table by filtering first. And there is actually no aggregation or abstraction, just a select query by taking the projection. This is basically a free query from a computational expense standpoint. These computations are of course carried out in $\Rel$ where we have Beck-Chevalley and pullbacks in the underlying category $\Set$. The point of axiomatizing Beck-Chevalley (perhaps through the assumption of strong tabulators) in the olog is that the structure-preserving data instance will, by functorial and fibrational semantics, can execute the query by either path allowed by the rewrite rule.
\end{example}

To conclude, briefly we note that coproducts satisfying Beck-Chevalley automatically gives an interpretation of fibered equality in the usual fibrational semantics of equational type theories.

\begin{corollary} \label{cor:equality}
  Any `double category of relations' for which dependent coproducts satisfy Beck-Chevalley is a Eq-fibration in the sense of \cite[\S\S 3.4-3.5]{jacobs1999}.
\end{corollary}
\begin{proof}
  Such a double category always has all dependent coproducts from the equipment structure. Hence it has fibered equality as in \cite[\S 3.4]{jacobs1999} since Beck-Chevalley follows as a special case. Any `double category of relations' has local products (reviewed below in \cref{section:local-products-exponentiability}).
\end{proof}

\begin{remark}
  A native treatment of equality is somewhat beside the point in the present work as relational database equality is closer to comparing common columns along diagonals, that is, basically a join operation (already studied in \cite{lambert2025}). Equality here is substitution of terms into fibered equality predicates derived from the dependent coproducts applied to local terminals.
\end{remark}

\section{Modality \& Tabulators}
\label{section:modality}

We are left with accounting for the modal queries in 
\cref{fig:quant-queries}. Viewing $r\colon x\proto y$ as a fixed 
proarrow, or perhaps a fixed relational attribute in an olog, we 
have the operation
  \begin{equation} \label{eqn:possibility-operator}
    \diamondsuit_r:=r\odot - \colon \dbl D(y,1)\to\dbl D(x,1) \qquad\qquad p\colon y\proto 1 \quad \mapsto \quad r\odot p
  \end{equation}
by precomposition with $r$. This is a functor by interchange. And 
we will drop the subscripted `$r$' whenever it is convenient to do 
so. The point of the previous discussion and 
\cref{prop:possibility-rewrite-rule-1} is that a subtype of $y$ 
induces an effect / proposition $y\proto 1$ and then loose 
composition carries the force of existential quantification over 
the subtype thus realizing the modal operator.

\begin{definition}[Adjoint Modality] \label{def:adjoint-modality}
  The operator $\diamondsuit_r$ of \cref{eqn:possibility-operator}
  is the left adjoint of an \textbf{adjoint modality} if it has a 
  right adjoint 
    \begin{equation*}
      \Box\colon\dbl D(x,1)\to\dbl D(y,1) \qquad \diamondsuit_r\dashv\Box
    \end{equation*}
  called \emph{(reverse) necessity}. The pair is referred to as an adjoint modality where $\Box$ is the right adjoint of the modality.
\end{definition}

The reason for the qualifier \emph{reverse} will be made clear below. Notice that we have not constructed $\Box$ at any point so far in an elementary way in a general cartesian equipment. The double category of relations, however, certainly has such adjoint modalities for each fixed relation. Under the definition above, these operators justly bear the standard notations not only on account of the $\Rel$-semantics. There is also the following observation, namely, that these operators satisfy the usual monad and comonad laws, as well as the characteristic (B) axiom of modal S5. In this sense, the definition is perhaps to strong, since S5 is a relatively refined system; however, our construction of the necessity operators using dependent products will satisfy these laws anyway.

\begin{proposition} \label{prop:modal-properties}
  Let $r\colon x\proto y$ induce an adjoint modality (\cref{def:adjoint-modality}). If $r$ is a preorder, then the possibility and necessity operators associated to $r$ satisfy the monad and comonad axioms 
    \begin{enumerate}
      \item $\id\leq\diamondsuit$ and $\diamondsuit\diamondsuit\leq\diamondsuit$
      \item $\Box\leq \id$ and $\Box\Box\leq\Box$
    \end{enumerate}
  as well as the further (B) axioms, namely, 
    \begin{enumerate}
      \item $\id\leq \Box\diamondsuit$
      \item $\diamondsuit\Box\leq \id$.
    \end{enumerate}
\end{proposition}
\begin{proof}
  Since $\diamondsuit\dashv\Box$, it suffices to show that $\diamondsuit$ satisfies the monad axioms as $\Box$ will automatically be a comonad. Likewise the (B) axioms are just a restatement of the adjunction assumption. Reflexivity of $r$ means that there is a canonical cell 
    \begin{equation*}
      % https://q.uiver.app/#q=WzAsNCxbMCwwLCJ4Il0sWzEsMCwieCJdLFsxLDEsIngiXSxbMCwxLCJ4Il0sWzAsMSwiXFxpZF94IiwwLHsic3R5bGUiOnsiYm9keSI6eyJuYW1lIjoiYmFycmVkIn19fV0sWzEsMiwiIiwwLHsibGV2ZWwiOjIsInN0eWxlIjp7ImhlYWQiOnsibmFtZSI6Im5vbmUifX19XSxbMCwzLCIiLDIseyJsZXZlbCI6Miwic3R5bGUiOnsiaGVhZCI6eyJuYW1lIjoibm9uZSJ9fX1dLFszLDIsInIiLDIseyJzdHlsZSI6eyJib2R5Ijp7Im5hbWUiOiJiYXJyZWQifX19XSxbNCw3LCJcXGRlbHRhX3giLDEseyJzaG9ydGVuIjp7InNvdXJjZSI6MjAsInRhcmdldCI6MjB9LCJzdHlsZSI6eyJib2R5Ijp7Im5hbWUiOiJub25lIn0sImhlYWQiOnsibmFtZSI6Im5vbmUifX19XV0=
      \begin{tikzcd}
        x & x \\
        x & x
        \arrow[""{name=0, anchor=center, inner sep=0}, "{\id_x}"{inner sep=.8ex}, "\shortmid"{marking}, from=1-1, to=1-2]
        \arrow[equals, from=1-1, to=2-1]
        \arrow[equals, from=1-2, to=2-2]
        \arrow[""{name=1, anchor=center, inner sep=0}, "r"'{inner sep=.8ex}, "\shortmid"{marking}, from=2-1, to=2-2]
        \arrow["{\delta_x}"{description}, draw=none, from=0, to=1]
      \end{tikzcd}
    \end{equation*}
  expressing the fact that the diagonal factors through $r$. But then composing we have 
    \begin{equation*}
      % https://q.uiver.app/#q=WzAsNixbMCwwLCIxIl0sWzEsMCwieCJdLFsxLDEsIngiXSxbMCwxLCIxIl0sWzIsMCwieCJdLFsyLDEsIngiXSxbMCwxLCJwIiwwLHsic3R5bGUiOnsiYm9keSI6eyJuYW1lIjoiYmFycmVkIn19fV0sWzEsMiwiIiwwLHsibGV2ZWwiOjIsInN0eWxlIjp7ImhlYWQiOnsibmFtZSI6Im5vbmUifX19XSxbMCwzLCIiLDIseyJsZXZlbCI6Miwic3R5bGUiOnsiaGVhZCI6eyJuYW1lIjoibm9uZSJ9fX1dLFszLDIsInAiLDIseyJzdHlsZSI6eyJib2R5Ijp7Im5hbWUiOiJiYXJyZWQifX19XSxbMSw0LCJcXGlkIiwwLHsic3R5bGUiOnsiYm9keSI6eyJuYW1lIjoiYmFycmVkIn19fV0sWzIsNSwiUiIsMix7InN0eWxlIjp7ImJvZHkiOnsibmFtZSI6ImJhcnJlZCJ9fX1dLFs0LDUsIiIsMSx7ImxldmVsIjoyLCJzdHlsZSI6eyJoZWFkIjp7Im5hbWUiOiJub25lIn19fV0sWzYsOSwiMV9wIiwxLHsic2hvcnRlbiI6eyJzb3VyY2UiOjIwLCJ0YXJnZXQiOjIwfSwic3R5bGUiOnsiYm9keSI6eyJuYW1lIjoibm9uZSJ9LCJoZWFkIjp7Im5hbWUiOiJub25lIn19fV0sWzEwLDExLCJcXGxlcSIsMSx7InNob3J0ZW4iOnsic291cmNlIjoyMCwidGFyZ2V0IjoyMH0sInN0eWxlIjp7ImJvZHkiOnsibmFtZSI6Im5vbmUifSwiaGVhZCI6eyJuYW1lIjoibm9uZSJ9fX1dXQ==
      \begin{tikzcd}
        1 & x & x \\
        1 & x & x
        \arrow[""{name=0, anchor=center, inner sep=0}, "p"{inner sep=.8ex}, "\shortmid"{marking}, from=1-1, to=1-2]
        \arrow[equals, from=1-1, to=2-1]
        \arrow[""{name=1, anchor=center, inner sep=0}, "\id"{inner sep=.8ex}, "\shortmid"{marking}, from=1-2, to=1-3]
        \arrow[equals, from=1-2, to=2-2]
        \arrow[equals, from=1-3, to=2-3]
        \arrow[""{name=2, anchor=center, inner sep=0}, "p"'{inner sep=.8ex}, "\shortmid"{marking}, from=2-1, to=2-2]
        \arrow[""{name=3, anchor=center, inner sep=0}, "r"'{inner sep=.8ex}, "\shortmid"{marking}, from=2-2, to=2-3]
        \arrow["{1_p}"{description}, draw=none, from=0, to=2]
        \arrow["\leq"{description}, draw=none, from=1, to=3]
      \end{tikzcd}
    \end{equation*}
  showing that $p\leq \diamondsuit p$ by loose composition and \cref{prop:possibility-rewrite-rule-1}. Likewise, using reflexivity and the same proposition, we have 
    \begin{equation*}
      % https://q.uiver.app/#q=WzAsNyxbMCwwLCIxIl0sWzEsMCwieCJdLFsxLDEsIngiXSxbMCwxLCIxIl0sWzIsMCwieCJdLFszLDEsIngiXSxbMywwLCJ4Il0sWzAsMSwicCIsMCx7InN0eWxlIjp7ImJvZHkiOnsibmFtZSI6ImJhcnJlZCJ9fX1dLFsxLDIsIiIsMCx7ImxldmVsIjoyLCJzdHlsZSI6eyJoZWFkIjp7Im5hbWUiOiJub25lIn19fV0sWzAsMywiIiwyLHsibGV2ZWwiOjIsInN0eWxlIjp7ImhlYWQiOnsibmFtZSI6Im5vbmUifX19XSxbMywyLCJwIiwyLHsic3R5bGUiOnsiYm9keSI6eyJuYW1lIjoiYmFycmVkIn19fV0sWzEsNCwiciIsMCx7InN0eWxlIjp7ImJvZHkiOnsibmFtZSI6ImJhcnJlZCJ9fX1dLFsyLDUsInIiLDIseyJzdHlsZSI6eyJib2R5Ijp7Im5hbWUiOiJiYXJyZWQifX19XSxbNCw2LCJyIiwwLHsic3R5bGUiOnsiYm9keSI6eyJuYW1lIjoiYmFycmVkIn19fV0sWzYsNSwiIiwwLHsibGV2ZWwiOjIsInN0eWxlIjp7ImhlYWQiOnsibmFtZSI6Im5vbmUifX19XSxbNywxMCwiMV9wIiwxLHsic2hvcnRlbiI6eyJzb3VyY2UiOjIwLCJ0YXJnZXQiOjIwfSwic3R5bGUiOnsiYm9keSI6eyJuYW1lIjoibm9uZSJ9LCJoZWFkIjp7Im5hbWUiOiJub25lIn19fV0sWzQsMTIsIlxcbGVxIiwxLHsic2hvcnRlbiI6eyJzb3VyY2UiOjIwLCJ0YXJnZXQiOjIwfSwic3R5bGUiOnsiYm9keSI6eyJuYW1lIjoibm9uZSJ9LCJoZWFkIjp7Im5hbWUiOiJub25lIn19fV1d
      \begin{tikzcd}
        1 & x & x & x \\
        1 & x && x
        \arrow[""{name=0, anchor=center, inner sep=0}, "p"{inner sep=.8ex}, "\shortmid"{marking}, from=1-1, to=1-2]
        \arrow[equals, from=1-1, to=2-1]
        \arrow["r"{inner sep=.8ex}, "\shortmid"{marking}, from=1-2, to=1-3]
        \arrow[equals, from=1-2, to=2-2]
        \arrow["r"{inner sep=.8ex}, "\shortmid"{marking}, from=1-3, to=1-4]
        \arrow[equals, from=1-4, to=2-4]
        \arrow[""{name=1, anchor=center, inner sep=0}, "p"'{inner sep=.8ex}, "\shortmid"{marking}, from=2-1, to=2-2]
        \arrow[""{name=2, anchor=center, inner sep=0}, "r"'{inner sep=.8ex}, "\shortmid"{marking}, from=2-2, to=2-4]
        \arrow["{1_p}"{description}, draw=none, from=0, to=1]
        \arrow["\leq"{description}, draw=none, from=1-3, to=2]
      \end{tikzcd}
    \end{equation*}
  so that $\diamondsuit \diamondsuit p \leq \diamondsuit p$ holds, as required.
\end{proof}

\begin{remark}
  There are pecularities of this set-up. First, although classical set-theoretic Kripke semantics require the underlying relation to be not only a preorder but also \emph{symmetric} to ensure that the characteristic (B) axioms of modal S5 are satisfied. In fact, equivalence relations are precisely those accessibility relations modeling classical S5. In the present instance, the assumption of the existence of the adjoint modality of course just axiomatizes this axiom scheme, so it suffices to ask that the underlying relation is a preorder to get the characteristic monad and comonad laws. Now, this adjointness is warranted as although classical possibility and necessity are interdefinable using negation, no such relationship exists in general categorical semantics due to the general failure of double negation. Essentially, that is, we ask for some kind of governing interaction between the two operators, and those laws of S5 in the form of the (B) axioms are the clearest, cleanest and most intuitive and well-established. Whereas on the other hand, the various proposed weak interactions axiomatized for intuitionistic S4 are not universally agreed upon and do not imply interdefinability anyway. Second, the tabulator models of necessity in the next section arise from such adjunctions anyway and will thus satisfy the (B) axiom scheme.
\end{remark}

Since possibility can be viewed as loose composition, the modal necessity operator could certainly just be asked for in the form of something like \emph{closedness} \cite[\S 5]{shulman2008}. Our interest, however, is in recovering necessity from the assumed structure of dependent products, since, as we have seen, necessity is universal quantification over accessible states. \cref{prop:possibility-rewrite-rule-1} shows that possibility is a composite operation, namely, a substitution followed by a dependent coproduct further confirming that a right adjoint to subsitution would then plausibly recover the necessity operator too. That this is so in $\Rel$ can easily be calculated by hand. This calculation suggests the general case. In fact, this perspective has already been considered in \cite[\S 3]{hermida2011} where the modal operators arise as composite substitutions followed by quantification functors. This approach is both bicategorical in nature and starts from the bicategories of spans or of relations where every proarrow is essentially its own tabulator. Our approach applies the same general philosophy to double categories in the following way. First recall the standard notion \cite{grandis1999}.

\begin{definition} \label{def:tabulator}
  A double category $\dbl D$ \textbf{has tabulators} if the external identity functor has a right adjoint 
    \begin{equation*}
      % https://q.uiver.app/#q=WzAsMixbMCwwLCJcXGRibCBEXzAiXSxbMSwwLCJcXGRibHtEfV8xIl0sWzAsMSwiXFxpZF97LX0iLDAseyJvZmZzZXQiOi0yfV0sWzEsMCwiVCIsMCx7Im9mZnNldCI6LTJ9XV0=
      \begin{tikzcd}
        {\dbl D_0} & {\dbl{D}_1}
        \arrow["{\id_{-}}", shift left=2, from=1-1, to=1-2]
        \arrow["T", shift left=2, from=1-2, to=1-1]
      \end{tikzcd} \qquad \qquad \id_{-}\dashv T.
    \end{equation*}
  The \textbf{tabulator} of a proarrow $m\colon x\proto y$ is a cell $\tau_m$ having the universal property that there exists a unique morphism $h$ making an equality of diagrams
    \begin{equation*}
      % https://q.uiver.app/#q=WzAsNCxbMCwwLCJUbSJdLFswLDEsIngiXSxbMSwwLCJUbSJdLFsxLDEsInkiXSxbMCwxLCJkIiwyXSxbMCwyLCIiLDIseyJsZXZlbCI6Miwic3R5bGUiOnsiYm9keSI6eyJuYW1lIjoiYmFycmVkIn0sImhlYWQiOnsibmFtZSI6Im5vbmUifX19XSxbMiwzLCJjIl0sWzEsMywibSIsMix7InN0eWxlIjp7ImJvZHkiOnsibmFtZSI6ImJhcnJlZCJ9fX1dLFs1LDcsIlxcdGF1X20iLDEseyJzaG9ydGVuIjp7InNvdXJjZSI6MjAsInRhcmdldCI6MjB9LCJzdHlsZSI6eyJib2R5Ijp7Im5hbWUiOiJub25lIn0sImhlYWQiOnsibmFtZSI6Im5vbmUifX19XV0=
      \begin{tikzcd}
        Tm & Tm \\
        x & y
        \arrow[""{name=0, anchor=center, inner sep=0}, "\shortmid"{marking}, equals, from=1-1, to=1-2]
        \arrow["d"', from=1-1, to=2-1]
        \arrow["c", from=1-2, to=2-2]
        \arrow[""{name=1, anchor=center, inner sep=0}, "m"'{inner sep=.8ex}, "\shortmid"{marking}, from=2-1, to=2-2]
        \arrow["{\tau_m}"{description}, draw=none, from=0, to=1]
      \end{tikzcd} \qquad = \qquad 
      % https://q.uiver.app/#q=WzAsNixbMCwwLCJUbSJdLFswLDEsInoiXSxbMSwwLCJUbSJdLFsxLDEsInoiXSxbMCwyLCJ4Il0sWzEsMiwieSJdLFswLDEsImgiLDJdLFswLDIsIiIsMix7ImxldmVsIjoyLCJzdHlsZSI6eyJib2R5Ijp7Im5hbWUiOiJiYXJyZWQifSwiaGVhZCI6eyJuYW1lIjoibm9uZSJ9fX1dLFsyLDMsImgiXSxbMSwzLCIiLDIseyJsZXZlbCI6Miwic3R5bGUiOnsiYm9keSI6eyJuYW1lIjoiYmFycmVkIn0sImhlYWQiOnsibmFtZSI6Im5vbmUifX19XSxbMSw0LCJmIiwyXSxbNCw1LCJtIiwyLHsic3R5bGUiOnsiYm9keSI6eyJuYW1lIjoiYmFycmVkIn19fV0sWzMsNSwiZyJdLFs3LDksIlxcaWRfaCIsMSx7InNob3J0ZW4iOnsic291cmNlIjoyMCwidGFyZ2V0IjoyMH0sInN0eWxlIjp7ImJvZHkiOnsibmFtZSI6Im5vbmUifSwiaGVhZCI6eyJuYW1lIjoibm9uZSJ9fX1dLFs5LDExLCJcXHRoZXRhIiwxLHsic2hvcnRlbiI6eyJzb3VyY2UiOjIwLCJ0YXJnZXQiOjIwfSwic3R5bGUiOnsiYm9keSI6eyJuYW1lIjoibm9uZSJ9LCJoZWFkIjp7Im5hbWUiOiJub25lIn19fV1d
      \begin{tikzcd}
        Tm & Tm \\
        z & z \\
        x & y
        \arrow[""{name=0, anchor=center, inner sep=0}, "\shortmid"{marking}, equals, from=1-1, to=1-2]
        \arrow["h"', from=1-1, to=2-1]
        \arrow["h", from=1-2, to=2-2]
        \arrow[""{name=1, anchor=center, inner sep=0}, "\shortmid"{marking}, equals, from=2-1, to=2-2]
        \arrow["f"', from=2-1, to=3-1]
        \arrow["g", from=2-2, to=3-2]
        \arrow[""{name=2, anchor=center, inner sep=0}, "m"'{inner sep=.8ex}, "\shortmid"{marking}, from=3-1, to=3-2]
        \arrow["{\id_h}"{description}, draw=none, from=0, to=1]
        \arrow["\theta"{description}, draw=none, from=1, to=2]
      \end{tikzcd}
    \end{equation*}
  for each such cell $\theta$.
\end{definition}

Tabulators viewed as generalized elements constructions ought always to roundtrip in one direction, namely, that from proarrows to spans and then back to proarrows. This is to say that the elements construction of the base-valued functor representing the original fibration object ought to return the original fibration object. Note that the other roundtrip works in the case of ordinary fibrations and discrete fibrations but fails in \emph{2-toposes} \cite{weber2007} due to issues of \emph{size} and \emph{admissibility}. For this reason we only axiomatize the former direction. The definition formalizing this is then the following one \cite[\S 5.1]{aleiferi2018}.

\begin{definition} \label{def:strong-tabulator}
  The tabulator of a proarrow $p\colon X\proto Y$ is \textbf{strong} if the canonical morphism 
    \begin{equation*}
      % https://q.uiver.app/#q=WzAsNCxbMCwwLCJUcCJdLFsxLDAsIlRwIl0sWzEsMSwieSJdLFswLDEsIngiXSxbMCwxLCIiLDAseyJsZXZlbCI6Miwic3R5bGUiOnsiYm9keSI6eyJuYW1lIjoiYmFycmVkIn0sImhlYWQiOnsibmFtZSI6Im5vbmUifX19XSxbMSwyLCJjIl0sWzAsMywiZCIsMl0sWzMsMiwicCIsMix7InN0eWxlIjp7ImJvZHkiOnsibmFtZSI6ImJhcnJlZCJ9fX1dLFs0LDcsIlxcdGF1X3AiLDEseyJzaG9ydGVuIjp7InNvdXJjZSI6MjAsInRhcmdldCI6MjB9LCJzdHlsZSI6eyJib2R5Ijp7Im5hbWUiOiJub25lIn0sImhlYWQiOnsibmFtZSI6Im5vbmUifX19XV0=
      \begin{tikzcd}
        Tp & Tp \\
        x & y
        \arrow[""{name=0, anchor=center, inner sep=0}, "\shortmid"{marking}, equals, from=1-1, to=1-2]
        \arrow["d"', from=1-1, to=2-1]
        \arrow["c", from=1-2, to=2-2]
        \arrow[""{name=1, anchor=center, inner sep=0}, "p"'{inner sep=.8ex}, "\shortmid"{marking}, from=2-1, to=2-2]
        \arrow["{\tau_p}"{description}, draw=none, from=0, to=1]
      \end{tikzcd} \qquad = \qquad
      % https://q.uiver.app/#q=WzAsNixbMCwwLCJUcCJdLFsxLDAsIlRwIl0sWzEsMSwieSJdLFswLDEsIngiXSxbMCwyLCJ4Il0sWzEsMiwieSJdLFswLDEsIiIsMCx7ImxldmVsIjoyLCJzdHlsZSI6eyJib2R5Ijp7Im5hbWUiOiJiYXJyZWQifSwiaGVhZCI6eyJuYW1lIjoibm9uZSJ9fX1dLFsxLDIsImMiXSxbMCwzLCJkIiwyXSxbMywyLCJkXipcXG9kb3QgY18hIiwyLHsic3R5bGUiOnsiYm9keSI6eyJuYW1lIjoiYmFycmVkIn19fV0sWzMsNCwiIiwyLHsibGV2ZWwiOjIsInN0eWxlIjp7ImhlYWQiOnsibmFtZSI6Im5vbmUifX19XSxbMiw1LCIiLDIseyJsZXZlbCI6Miwic3R5bGUiOnsiaGVhZCI6eyJuYW1lIjoibm9uZSJ9fX1dLFs0LDUsInAiLDIseyJzdHlsZSI6eyJib2R5Ijp7Im5hbWUiOiJiYXJyZWQifX19XSxbNiw5LCJcXHRhdV9wIiwxLHsic2hvcnRlbiI6eyJzb3VyY2UiOjIwLCJ0YXJnZXQiOjIwfSwic3R5bGUiOnsiYm9keSI6eyJuYW1lIjoibm9uZSJ9LCJoZWFkIjp7Im5hbWUiOiJub25lIn19fV0sWzksMTIsIlxcZXhpc3RzXFwsISIsMSx7ImxhYmVsX3Bvc2l0aW9uIjo2MCwic2hvcnRlbiI6eyJzb3VyY2UiOjIwLCJ0YXJnZXQiOjIwfSwic3R5bGUiOnsiYm9keSI6eyJuYW1lIjoibm9uZSJ9LCJoZWFkIjp7Im5hbWUiOiJub25lIn19fV1d
      \begin{tikzcd}
        Tp & Tp \\
        x & y \\
        x & y
        \arrow[""{name=0, anchor=center, inner sep=0}, "\shortmid"{marking}, equals, from=1-1, to=1-2]
        \arrow["d"', from=1-1, to=2-1]
        \arrow["c", from=1-2, to=2-2]
        \arrow[""{name=1, anchor=center, inner sep=0}, "{d^*\odot c_!}"'{inner sep=.8ex}, "\shortmid"{marking}, from=2-1, to=2-2]
        \arrow[equals, from=2-1, to=3-1]
        \arrow[equals, from=2-2, to=3-2]
        \arrow[""{name=2, anchor=center, inner sep=0}, "p"'{inner sep=.8ex}, "\shortmid"{marking}, from=3-1, to=3-2]
        \arrow["{\tau_p}"{description}, draw=none, from=0, to=1]
        \arrow["{\exists\,!}"{description, pos=0.6}, draw=none, from=1, to=2]
      \end{tikzcd}
    \end{equation*}
    is an isomorphism. A double category has \textbf{strong tabulators} if it has tabulators and each is strong.
\end{definition}

\begin{remark}
  This means, equivalently, that each such tabulator cell $\tau_p$ is an extension cell. Consequently, a double category has strong tabulators if each of its proarrows is the extension of its tabulator. The double category of relations is an example.
\end{remark}

\begin{lemma}
  For any such strong tabulator, the three constructions of $\diamondsuit \phi$ coincide in the sense that
    \begin{equation*}
      \coprod_{d}c^*\phi \cong r\odot\phi 
    \end{equation*}
  holds for any proposition $\phi\colon y\proto 1$ and $r\colon x\proto y$ whose tabulator is strong.
\end{lemma}
\begin{proof}
  Consider 
    \begin{equation*}
      \coprod_{d}c^*\phi \cong d^*\odot c_!\odot \phi \cong r\odot\phi
    \end{equation*}
  by the definition of the action $c^*(-) = c_!\odot -$ and that $r$ has a strong tabulator. If $\phi$ arises as the extension of a subtype or other arrow into $y$, apply \cref{prop:possibility-rewrite-rule-1} to see that the other possible definition agrees with these two.
\end{proof}

\begin{definition} \label{def:necessity-from-dep-products}
  Let $\dbl D$ denote a cartesian equipment with strong tabulators and dependent products. The \textbf{modal necessity} operator associated to $r\colon x\proto y$ is the composite functor 
    \begin{equation*}
      \Box_r:=\prod_{d}c^*(-)\colon\dbl D(y,1)\to\dbl D(x,1)
    \end{equation*}
  where $\langle d,c\rangle\colon Tr\to x\times y$ is the tabulator of $r$.
\end{definition}

\begin{remark}
  \cref{def:necessity-from-dep-products} is to priviledge modalities in the \emph{same direction} as in \cite[Remark 3.3]{hermida2011}. That is, both necessity and possibility for a given $r$ return values of its domain. In the measurement interpretation, this is to return the states for which all measurements fall inside the \emph{good} values specified by the proposition in question. Now, in any such double category, there is a reverse operation $r\mapsto r^\dagger$ by taking an extension along the legs of the tabulator of $r$ in reversed order. This coincides with the operation $r\mapsto r^\circ$ from \cite{carboni1987} as discussed above in good cases \cite{lambert2022}. Thus, there are actually four operators
    \begin{enumerate}
      \item $\Box_r:=\prod_{d}c^*(-)$
      \item $\Box_{r^\dagger}:=\prod_{c}d^*(-)$
      \item $\diamondsuit_r:=\coprod_dc^*(-)$
      \item $\diamondsuit_{r^\dagger}:=\coprod_cd^*(-)$
    \end{enumerate}
  This four-fold modal set-up is a feature of such relational and hyperdoctrine accounts of modality. The \emph{up-modalities} are the first and third above, returning states where outcomes stay within the specified allowed values; whereas the \emph{down-modalities} are the second and fourth, returing the outcomes whose states are among the specified ones. When $r$ is an endorelation, up- or down-modalities simply specify subsequent or antecedent states. In any case, the two types are inter-related in the following way.
\end{remark}

\begin{lemma}
  Under the hypotheses of \cref{def:necessity-from-dep-products} above, the (B) axioms are satisfied and the pairs $\diamondsuit_r\dashv\Box_{r^\dagger}$ and $\diamondsuit_{r^\dagger}\dashv \Box_r$ are adjoint modalities as in \cref{def:adjoint-modality}.  If $r\colon x\proto x$ is a preorder, then all the (co)monad axioms for both up- and down-modalities are satisfied too.
\end{lemma}
\begin{proof}
  See \cref{prop:modal-properties} for the last part. For the adjoint modality statement, consider the diagram of adjunctions
    \begin{equation*}
      % https://q.uiver.app/#q=WzAsMyxbMCwwLCJcXGRibCBEKHksMSkiXSxbMSwwLCJcXGRibCBEKFRyLDEpIl0sWzIsMCwiXFxkYmwgRCh4LDEpIl0sWzAsMSwiY14qKC0pIiwwLHsib2Zmc2V0IjotMn1dLFsxLDAsIlxccHJvZF9jIiwwLHsib2Zmc2V0IjotMn1dLFsxLDIsIlxcY29wcm9kX2QiLDAseyJvZmZzZXQiOi0yfV0sWzIsMSwiZF4qKC0pIiwwLHsib2Zmc2V0IjotMn1dXQ==
      \begin{tikzcd}
        {\dbl D(y,1)} & {\dbl D(Tr,1)} & {\dbl D(x,1)}
        \arrow["{c^*(-)}", shift left=2, from=1-1, to=1-2]
        \arrow["{\prod_c}", shift left=2, from=1-2, to=1-1]
        \arrow["{\coprod_d}", shift left=2, from=1-2, to=1-3]
        \arrow["{d^*(-)}", shift left=2, from=1-3, to=1-2]
      \end{tikzcd}
    \end{equation*}
  showing that $\diamondsuit_r\dashv\Box_{r^\dagger}$, as required. The other statement is dual.
\end{proof}

\begin{example}
  When specialized to $\Rel$, the necessity operator here returns the necessity query in \cref{fig:quant-queries}, thus completing our account of those six queries.
\end{example}

\begin{remark}
  Evidently symmetry in the form $r^\dagger\leq r$ collapses the distinction between the possibility and necessity pairs: $\diamondsuit_r = \diamondsuit_{r^\dagger}$ and $\Box_r = \Box_{r^\dagger}$ so that directly $\diamondsuit_r\dashv\Box_r$. In this way, an equivalence relation $r$ induces an adjoint modality between the now just \emph{two} modal operators without the foregoing up- vs. down- distinction.
\end{remark}

\part{Cartesian Closedness \& FOML}

This second part concerns ultimately the way in which double ologs model first-order (modal) logic (FOML) locally. Inasmuch as the classical operations of relational databases are framable in FOML this means that suitably structured double ologs thus provide a complete synthetic framework for database querying. There are two technical aspects of this development. First is to connect cartesian closedness with the assumed dependent products coming with any $\prod$-double olog. The connection is made via the additional structure of tabulators. To show that this works, we study in some detail the modular laws and Frobenius reciprocity which are valid in any `double category of relations'. Our analysis of local products with identity proarrows leads to directly to our formulas for the general exponential. Secondly, the evident utility of a hypothetical negation operator leads to the natural question of cocartesian structure and whether, how and to what extent this results in local disjunction and how it interacts with local products. It turns out that, assuming a global distributive law, conjunction and disjunction distribute locally as well. These propositional and first-order laws are then compared with the various canonical modal statements such as the usual (K) axiom. Finally, it is shown that such first-order models interpret description logic.

\section{Local Products \& Exponentiability}
\label{section:local-products-exponentiability}

To continue the development, we work first with the given cartesian structure assumed in our ambient set-ups without assuming further structures as in later sections. If $\dbl D$ is a cartesian equipment, then each hom-category 
$\dbl D(x,y)$ of course has cartesian products, given by the 
formula 
  \begin{equation*} \label{eqn:local-products-construction}
    % https://q.uiver.app/#q=WzAsNCxbMCwwLCJ4Il0sWzEsMCwieSJdLFsxLDEsIjEiXSxbMCwxLCIxIl0sWzAsMSwiMV97eCx5fSIsMCx7InN0eWxlIjp7ImJvZHkiOnsibmFtZSI6ImJhcnJlZCJ9fX1dLFsxLDIsIiEiXSxbMCwzLCIhIiwyXSxbMywyLCIiLDIseyJzdHlsZSI6eyJib2R5Ijp7Im5hbWUiOiJiYXJyZWQifX19XSxbNCw3LCJcXG1hdGhybXtyZXN9IiwxLHsic2hvcnRlbiI6eyJzb3VyY2UiOjIwLCJ0YXJnZXQiOjIwfSwic3R5bGUiOnsiYm9keSI6eyJuYW1lIjoibm9uZSJ9LCJoZWFkIjp7Im5hbWUiOiJub25lIn19fV1d
    \begin{tikzcd}
      x & y \\
      1 & 1
      \arrow[""{name=0, anchor=center, inner sep=0}, "{1_{x,y}}"{inner sep=.8ex}, "\shortmid"{marking}, from=1-1, to=1-2]
      \arrow["{!}"', from=1-1, to=2-1]
      \arrow["{!}", from=1-2, to=2-2]
      \arrow[""{name=1, anchor=center, inner sep=0}, "\shortmid"{marking}, from=2-1, to=2-2]
      \arrow["{\mathrm{res}}"{description}, draw=none, from=0, to=1]
    \end{tikzcd}\qquad\qquad
    % https://q.uiver.app/#q=WzAsNCxbMCwxLCJ4XFx0aW1lcyB4Il0sWzIsMSwieVxcdGltZXMgeSJdLFswLDAsIngiXSxbMiwwLCJ5Il0sWzAsMSwibVxcdGltZXMgbiIsMix7InN0eWxlIjp7ImJvZHkiOnsibmFtZSI6ImJhcnJlZCJ9fX1dLFsyLDAsIlxcRGVsdGFfeCIsMl0sWzIsMywibVxcd2VkZ2UgbiIsMCx7InN0eWxlIjp7ImJvZHkiOnsibmFtZSI6ImJhcnJlZCJ9fX1dLFszLDEsIlxcRGVsdGFfeSJdLFs2LDQsIlxcbWF0aHJte3Jlc30iLDEseyJzaG9ydGVuIjp7InNvdXJjZSI6MjAsInRhcmdldCI6MjB9LCJzdHlsZSI6eyJib2R5Ijp7Im5hbWUiOiJub25lIn0sImhlYWQiOnsibmFtZSI6Im5vbmUifX19XV0=
    \begin{tikzcd}
      x && y \\
      {x\times x} && {y\times y}
      \arrow[""{name=0, anchor=center, inner sep=0}, "{m\wedge n}"{inner sep=.8ex}, "\shortmid"{marking}, from=1-1, to=1-3]
      \arrow["{\Delta_x}"', from=1-1, to=2-1]
      \arrow["{\Delta_y}", from=1-3, to=2-3]
      \arrow[""{name=1, anchor=center, inner sep=0}, "{m\times n}"'{inner sep=.8ex}, "\shortmid"{marking}, from=2-1, to=2-3]
      \arrow["{\mathrm{res}}"{description}, draw=none, from=0, to=1]
    \end{tikzcd}
  \end{equation*}
as described in \cite[Proposition 4.3.2]{aleiferi2018}. Note that the wedge operator is thus understood also as a loose composite 
  \begin{equation*}
    m\wedge n \cong \Delta_!\odot (m\times n) \odot \Delta^*
  \end{equation*}
by the usual construction of restrictions. In any case, this leads 
naturally to the following definitions. 

\begin{definition}\label{def:exponentiable-and-cartesianclosed}
  A proarrow $m$ in a cartesian equipment $\dbl D$ is \textbf{exponentiable} if each product functor 
    \begin{equation*}
      m\wedge-\colon \dbl{D}(x,y)\to\dbl{D}(x,y)
    \end{equation*}
  has a right adjoint $(-)^m$. The object $n^m$ is called the \textbf{exponential} of $n$ by $m$. A cartesian equipment is \textbf{locally cartesian closed} if each proarrow is exponentiable.
\end{definition}

\begin{example}
  For relations $R,S\colon X\proto Y$, the exponential is given by 
    \begin{equation*}
      S^R = \lbrace (x,y)\mid xRy \text{ implies } xSy\rbrace
    \end{equation*}
  and computed via adjoints as 
    \begin{equation*}
      \forall_RR^*(S)
    \end{equation*}
  thinking of $R$ as its own tabulator.
\end{example}

Of course we may simply ask that our cartesian equipments are locally cartesian closed in this sense. It is a point of interest, however, to connect this structure to the right adjoints to substitution assumed in previous sections inasmuch as, (1) these are supposed to provide \emph{generalized hom objects} and (2) there is an evident universal quantification over pairs $(x,y)$ in the example above. The answer, roughly, is that these right adjoints are computable from right adjoints to substitution whenever the given cartesian equipment has strong tabulators. This development showcases the relationships between compactness, Frobenius, modularity, comprehension schemes, and exponentiability all in terms of the requested equipment and cartesian structure on the ambient double category.

\section{Modularity \& Frobenius}
\label{section:modularity-frobenius}

It is asserted in \cite[\S 3]{lambert2022} and proved in the cited references that any `double category of relations' satisfies the usual \emph{Frobenius reciprocity Law} for the accompanying doctrine. This is proved here in detail since the proof is interesting in its own right and the law is the main ingredient in proving that right adjoints to substitution do give cartesian closed structure, locally.

Recall from \cite{carboni1987} that any `bicategory of relations' is compact closed and has an involution operation 
  \begin{equation*}
    f \mapsto f^\circ
  \end{equation*}
satisfying 
  \begin{equation*}
    1^\circ = 1\qquad (f^\circ)^\circ = f\qquad (gf)^\circ = f^\circ g^\circ
  \end{equation*}
with $(-)^\circ$ defined using the unit and counit in the compact closed structure as in the proof of \cite[Theorem 2.4]{carboni1987}. Now, if $f$ is a \emph{map} then $f^\circ$ is that right adjoint $f\dashv f^\circ$ \cite[Lemma 2.5]{carboni1987}. This extends to `double categories of relations' in the following way.

\begin{lemma}
  In any `double category of relations', the equations
    \begin{enumerate}
      \item $(f_!)^\circ = f^*$
      \item $(f^*)^\circ = f_!$
    \end{enumerate}
  each hold.
\end{lemma}
\begin{proof}
  Start with the usual equipment adjunction $f_!\dashv f^*$ \cite[\S 5]{shulman2008}. The underlying bicategory of any `double category of relations' is a `bicategory of relations' \cite{lambert2022} and bicategorical adjoints are unique up to isomorphism.
\end{proof}

The proof of the modular law (cf. \cite{freyd1990}) is now an application. It is phrased here in the manner needed in proving Frobenius.

\begin{proposition}[modular laws] \label{prop:modular-laws}
  In any `double category of relations', the two inequalities 
    \begin{enumerate}
      \item $(f^*\odot m)\wedge n \leq f^*\odot (m\wedge f_!\odot n)$
      \item $(p\odot f_!)\wedge q \leq (p\wedge q\odot f^*)\odot f_!$
    \end{enumerate}
  each hold.
\end{proposition}
\begin{proof}
  Two inequalities follow directly from \cite[Theorem 2.4]{carboni1987}. These are 
    \begin{enumerate}
      \item $\Delta_!\odot (f^*\times 1) \leq f^*\odot\Delta_!\odot (1\times f_!)$
      \item $(f_!\times 1)\odot \Delta^*\leq (1\times f^*)\odot \Delta^*\odot f_!$
    \end{enumerate}
  using $(f^*)^\circ = f_!$ as in the lemma above for the first. The two modular laws are direct applications. For the first one, we have 
    \begin{align*}
      (f^*\odot m)\wedge n &= \Delta_!\odot (f^*\odot m\times n)\odot\Delta^* \\
                           &= \Delta_!\odot (f^*\times 1\odot m\times n)\odot\Delta^* \\
                           &\leq f^*\odot\Delta_!\odot (1\times f_! \odot m\times n)\odot\Delta^* \\
                           &= f^*\odot\Delta_!\odot (m\times f_! \odot n)\odot\Delta^* \\
                           &= f^*\odot (m\wedge f_!\odot n)
    \end{align*}
  using the interaction of $\odot$ and $\times$ as well as the definition of local products. The other inequality is analogous.
\end{proof}

\begin{construction}[Frobenius Comparison Cell]\label{construction:Frob-comparison-cell}
  Let $\dbl D$ denote any cartesian equipment. Local products make sense in this context. Let $f\colon x\to z$ and $g\colon y\to w$ denote morphisms. First form the coproduct extension as the third step in the sequence:
    \begin{equation*}
      % https://q.uiver.app/#q=WzAsNCxbMCwwLCJ4Il0sWzAsMSwieiJdLFsxLDAsInkiXSxbMSwxLCJ3Il0sWzAsMSwiZiIsMl0sWzIsMywiZyJdLFsxLDMsInAiLDIseyJzdHlsZSI6eyJib2R5Ijp7Im5hbWUiOiJiYXJyZWQifX19XSxbMCwyLCJmXyFwZ14qIiwwLHsic3R5bGUiOnsiYm9keSI6eyJuYW1lIjoiYmFycmVkIn19fV0sWzcsNiwiXFxtYXRocm17cmVzfSIsMSx7InNob3J0ZW4iOnsic291cmNlIjoyMCwidGFyZ2V0IjoyMH0sInN0eWxlIjp7ImJvZHkiOnsibmFtZSI6Im5vbmUifSwiaGVhZCI6eyJuYW1lIjoibm9uZSJ9fX1dXQ==
      \begin{tikzcd}
        x & y \\
        z & w
        \arrow[""{name=0, anchor=center, inner sep=0}, "{f_!pg^*}"{inner sep=.8ex}, "\shortmid"{marking}, from=1-1, to=1-2]
        \arrow["f"', from=1-1, to=2-1]
        \arrow["g", from=1-2, to=2-2]
        \arrow[""{name=1, anchor=center, inner sep=0}, "p"'{inner sep=.8ex}, "\shortmid"{marking}, from=2-1, to=2-2]
        \arrow["{\mathrm{res}}"{description}, draw=none, from=0, to=1]
      \end{tikzcd}\quad\leadsto\quad
      % https://q.uiver.app/#q=WzAsNCxbMCwwLCJ4Il0sWzIsMCwieSJdLFswLDEsInhcXHRpbWVzIHgiXSxbMiwxLCJ5XFx0aW1lcyB5Il0sWzAsMSwibVxcd2VkZ2UgZl8hcGdeKiIsMCx7InN0eWxlIjp7ImJvZHkiOnsibmFtZSI6ImJhcnJlZCJ9fX1dLFswLDIsIlxcRGVsdGFfeCIsMl0sWzEsMywiXFxEZWx0YV95Il0sWzIsMywibVxcdGltZXMgZl8hcGdeKiIsMix7InN0eWxlIjp7ImJvZHkiOnsibmFtZSI6ImJhcnJlZCJ9fX1dLFs0LDcsIlxcbWF0aHJte3Jlc30iLDEseyJzaG9ydGVuIjp7InNvdXJjZSI6MjAsInRhcmdldCI6MjB9LCJzdHlsZSI6eyJib2R5Ijp7Im5hbWUiOiJub25lIn0sImhlYWQiOnsibmFtZSI6Im5vbmUifX19XV0=
      \begin{tikzcd}
        x && y \\
        {x\times x} && {y\times y}
        \arrow[""{name=0, anchor=center, inner sep=0}, "{m\wedge f_!pg^*}"{inner sep=.8ex}, "\shortmid"{marking}, from=1-1, to=1-3]
        \arrow["{\Delta_x}"', from=1-1, to=2-1]
        \arrow["{\Delta_y}", from=1-3, to=2-3]
        \arrow[""{name=1, anchor=center, inner sep=0}, "{m\times f_!pg^*}"'{inner sep=.8ex}, "\shortmid"{marking}, from=2-1, to=2-3]
        \arrow["{\mathrm{res}}"{description}, draw=none, from=0, to=1]
      \end{tikzcd}\quad\leadsto\quad
      % https://q.uiver.app/#q=WzAsNCxbMCwwLCJ4Il0sWzIsMCwieSJdLFswLDEsInoiXSxbMiwxLCJ3Il0sWzAsMSwibVxcd2VkZ2UgZl8hcGdeKiIsMCx7InN0eWxlIjp7ImJvZHkiOnsibmFtZSI6ImJhcnJlZCJ9fX1dLFswLDIsImYiLDJdLFsxLDMsImciXSxbMiwzLCJcXGNvcHJvZF97ZixnfShtXFx3ZWRnZSBmXyFwZ14qKSIsMix7InN0eWxlIjp7ImJvZHkiOnsibmFtZSI6ImJhcnJlZCJ9fX1dLFs0LDcsIlxcbWF0aHJte2V4dH0iLDEseyJzaG9ydGVuIjp7InNvdXJjZSI6MjAsInRhcmdldCI6MjB9LCJzdHlsZSI6eyJib2R5Ijp7Im5hbWUiOiJub25lIn0sImhlYWQiOnsibmFtZSI6Im5vbmUifX19XV0=
      \begin{tikzcd}
        x && y \\
        z && w
        \arrow[""{name=0, anchor=center, inner sep=0}, "{m\wedge f_!pg^*}"{inner sep=.8ex}, "\shortmid"{marking}, from=1-1, to=1-3]
        \arrow["f"', from=1-1, to=2-1]
        \arrow["g", from=1-3, to=2-3]
        \arrow[""{name=1, anchor=center, inner sep=0}, "{\coprod_{f,g}(m\wedge f_!pg^*)}"'{inner sep=.8ex}, "\shortmid"{marking}, from=2-1, to=2-3]
        \arrow["{\mathrm{ext}}"{description}, draw=none, from=0, to=1]
      \end{tikzcd}.
    \end{equation*}
  Note that we are suppressing some of the loose composition notation for readability. Now, on the other hand, form the local product as the second step in the sequence:
    \begin{equation*}
      % https://q.uiver.app/#q=WzAsNCxbMCwwLCJ4Il0sWzAsMSwieiJdLFsxLDAsInkiXSxbMSwxLCJ3Il0sWzAsMSwiZiIsMl0sWzIsMywiZyJdLFsxLDMsIlxcY29wcm9kX3tmLGd9bSIsMix7InN0eWxlIjp7ImJvZHkiOnsibmFtZSI6ImJhcnJlZCJ9fX1dLFswLDIsIm0iLDAseyJzdHlsZSI6eyJib2R5Ijp7Im5hbWUiOiJiYXJyZWQifX19XSxbNyw2LCJcXG1hdGhybXtleHR9IiwxLHsic2hvcnRlbiI6eyJzb3VyY2UiOjIwLCJ0YXJnZXQiOjIwfSwic3R5bGUiOnsiYm9keSI6eyJuYW1lIjoibm9uZSJ9LCJoZWFkIjp7Im5hbWUiOiJub25lIn19fV1d
      \begin{tikzcd}
        x & y \\
        z & w
        \arrow[""{name=0, anchor=center, inner sep=0}, "m"{inner sep=.8ex}, "\shortmid"{marking}, from=1-1, to=1-2]
        \arrow["f"', from=1-1, to=2-1]
        \arrow["g", from=1-2, to=2-2]
        \arrow[""{name=1, anchor=center, inner sep=0}, "{\coprod_{f,g}m}"'{inner sep=.8ex}, "\shortmid"{marking}, from=2-1, to=2-2]
        \arrow["{\mathrm{ext}}"{description}, draw=none, from=0, to=1]
      \end{tikzcd}\qquad\leadsto\qquad 
      % https://q.uiver.app/#q=WzAsNCxbMCwwLCJ6Il0sWzIsMCwidyJdLFsyLDEsIndcXHRpbWVzIHciXSxbMCwxLCJ6XFx0aW1lcyB6Il0sWzAsMSwiKFxcY29wcm9kX3tmLGd9bSlcXHdlZGdlIHAiLDAseyJzdHlsZSI6eyJib2R5Ijp7Im5hbWUiOiJiYXJyZWQifX19XSxbMSwyLCJcXERlbHRhX3ciXSxbMCwzLCJcXERlbHRhX3oiLDJdLFszLDIsIihcXGNvcHJvZF97ZixnfW0pXFx0aW1lcyBwIiwyLHsic3R5bGUiOnsiYm9keSI6eyJuYW1lIjoiYmFycmVkIn19fV0sWzQsNywiXFxtYXRocm17cmVzfSIsMSx7InNob3J0ZW4iOnsic291cmNlIjoyMCwidGFyZ2V0IjoyMH0sInN0eWxlIjp7ImJvZHkiOnsibmFtZSI6Im5vbmUifSwiaGVhZCI6eyJuYW1lIjoibm9uZSJ9fX1dXQ==
      \begin{tikzcd}
        z && w \\
        {z\times z} && {w\times w}
        \arrow[""{name=0, anchor=center, inner sep=0}, "{(\coprod_{f,g}m)\wedge p}"{inner sep=.8ex}, "\shortmid"{marking}, from=1-1, to=1-3]
        \arrow["{\Delta_z}"', from=1-1, to=2-1]
        \arrow["{\Delta_w}", from=1-3, to=2-3]
        \arrow[""{name=1, anchor=center, inner sep=0}, "{(\coprod_{f,g}m)\times p}"'{inner sep=.8ex}, "\shortmid"{marking}, from=2-1, to=2-3]
        \arrow["{\mathrm{res}}"{description}, draw=none, from=0, to=1]
      \end{tikzcd}
    \end{equation*}
  By the universal property of the local product immediately above, there is thus a canonical comparison cell 
    \begin{equation*}
      m\wedge f_!pg^* \Rightarrow (\coprod_{f,g}m)\wedge p
    \end{equation*}
  and likewise a comparison cell 
    \begin{equation*}
      \coprod_{f,g}(m\wedge f_!pg^*) \Rightarrow (\coprod_{f,g}m)\wedge p
    \end{equation*}
  by the universal property of the extension cell defining the coproduct in the preantepenultimate display. This is the required \emph{Frobenius comparison cell}. Note that it exists even if, as assumed, $\dbl D$ is merely a cartesian equipment.
\end{construction}

\begin{proposition} \label{prop:Frobenius-reciprocity}
  Any `double category of relations' satisfies the Frobenius reciprocity Law. That is, the canonical cell 
    \begin{equation*}
      \coprod_{f,g}(m\wedge f_!pg^*) \leq (\coprod_{f,g}m)\wedge p
    \end{equation*}
  described above is invertible.
\end{proposition}
\begin{proof}
  For the proof, we produce an reverse inequality. Uniqueness then forces the inequality above to be an equality. Starting from the left side and using the construction of extensions and both modular laws (\cref{prop:modular-laws}), we have 
    \begin{align*}
      (\coprod_{f,g}m)\wedge p &=(f^*\odot m \odot g_!)\wedge p \\
                               &\leq f^*\odot (m\odot g^*\wedge f_!\odot p) \\
                               &\leq f^*\odot (m \wedge (f_!\odot p\odot g^*))\odot g_! \\
                               &=\coprod_{f,g}(m\wedge f_!pg^*) 
    \end{align*}
  as required.
\end{proof}

\begin{corollary}
  Any `double category of relations', hence any double olog, viewed as an equipment, is a regular fibration \cite[Definition 4.2.1]{jacobs1999}, hence a fibrational model of regular logic.
\end{corollary}
\begin{proof}
  It needs only to be observed that such a double olog has equality in the sense of \cite[\S 3.4]{jacobs1999} since it has \emph{all} coproducts satisfying Frobenius reciprocity by the above. This was noted above in \cref{cor:equality}.
\end{proof}

\begin{remark}[Query Optimization] \label{remark:Frobenius-Optimization}
  Frobenius reciprocity is a \emph{join pushdown} optimization rewrite rule, saying that a extension following an expensive join can be equivalently executed by first collapsing using an extension and then performing a smaller join. This has the effect of restricting to a smaller table via a collapse operation, hence surveying a smaller totality, \emph{before} computing the product. That is, the right hand side of the Frobenius reciprocity identity in \cref{prop:Frobenius-reciprocity} above is the optimized query and the result says it produces the same result. 
\end{remark}

\begin{example} \label{example:frobenius}
  We illustrate this with a simple example. This scenario is one of \emph{abstraction} and \emph{querying metadata}. We have a surjective function associating to each of several individuals his or her class:
    \begin{equation*}
      \begin{tabular}{| l | l | }
          \hline\multicolumn{2}{| c |}{\bf Belongs to}\\
          \hline {\bf Individual }&{\bf Class}\\
          \hline Belrand & spellsword \\
          \hline Farcas & warrior \\
          \hline Illia & mage \\
          \hline J'zargo & mage \\
          \hline Marcurio & mage \\
          \hline Teldryn Sero & spellsword \\
          \hline
      \end{tabular}. 
    \end{equation*}
  Suppose we have the following subsets
    \begin{equation*}
      S = \lbrace \text{Marcurio}, \text{ Illia}, \text{ Teldryn Sero}\rbrace
      \qquad T = \lbrace \text{spellsword}, \text{ warrior}\rbrace
    \end{equation*}
  and suppose we wish to know which of the classes on the latter list are represented by the individuals on the former list. Of course the example is simple enough that we can just look at the table. But for large datasets a process is of course required. A naive way to proceed is see who on the list falls into any of the specified classes and then just take the class information. That is, more formally, restrict $T$ and intersect with $S$ to see who on the list falls into the classes; and then to project back to the codomain to see which classes we ended up with:
    \begin{align*}
      &\exists_{\text{Belongs}}(S\wedge (\text{Belongs}^*(T)))\\ 
        = \;&\exists_{\text{Belongs}}(\lbrace \text{Marcurio}, \text{ Illia}, \text{ 
        Teldryn Sero}\rbrace\cap \lbrace \text{Belrand}, \text{ Farcas}, \text{ Teldryn Sero}\rbrace )\\
        = \;&\exists_{\text{Belongs}}(\lbrace \text{Teldryn Sero}\rbrace) \\
        = \;&\lbrace \text{spellsword}\rbrace.
    \end{align*}
  Note that three operations are required. The restriction and intersection are where most of the computation is done. But since we are working in relations where Frobenius reciprocity holds, the exact same effect is had by \emph{first} abstracting or aggregating the metadata and \text{then} intersecting:
    \begin{equation*}
      \exists_{\text{Belongs}}(S)\wedge T = \lbrace\text{spellsword}\rbrace\cap\lbrace \text{spellsword},\text{ warrior}\rbrace = \lbrace\text{spellsword}\rbrace.
    \end{equation*}
  Only two operations are required. The point is that Frobenius reciprocity formalizes this relationship for double ologs and the functorial and fibrational semantics of data instances control the computation of the query in the receiving data structure. Notably, the proposition above shows that it emerges as a derivable \emph{property} of a `double categories of relations'.
\end{example}

\section{Cartesian Closedness via Products \& Tabulators}
\label{section:cartesian-closedness}

As a preliminary, the first result concerns exponentiability. This is an argument that local products with identity proarrows are essentially restrictions along the relevant diagonal. In this sense they are somewhat like identity predicates \cite[\S 3.4]{jacobs1999} and in paricular are left adjoints. 

\begin{lemma} \label{lemma:identity-proarrow-exponentiable}
  Any identity proarrow in a $\prod$-double olog is exponentiable.
\end{lemma}
\begin{proof}
  The compact closed transpose of $\id\wedge p$ to a proarrow $x\times x\proto 1$ is the composite
    \begin{equation*}
      \widehat{p\wedge\id}  = ((p\wedge\id) \times \id) \odot\Delta^*\odot x_!
    \end{equation*}
  as in \cite[Theorem 2.4]{carboni1987}. Now, this can be rewritten into a more convenient form via the following calculation:
    \begin{align*}
      \widehat{p\wedge\id} &= (\Delta_!\times \id)\odot (p\times\id\times\id)\odot(\Delta^*\times\id)\odot\Delta^*\odot x_! \qquad &\text{(\text{Constr. Local Product})}\\
      &=(\Delta_!\times \id)\odot (p\times\id\times\id)\odot(\id\times\Delta^*)\odot\Delta^*\odot x_! \qquad &\text{(\text{Comonoid Axiom})}\\
      &=(\Delta_!\times \id)\odot (\id\times\Delta^*)\odot (p\times\id)\odot\Delta^*\odot x_! \qquad &\text{(\text{Functoriality of $\times$})}\\
      &=\Delta^*\odot \Delta_!\odot (p\times\id) \odot\Delta^*\odot x_! \qquad &\text{(\text{Discreteness})}
    \end{align*}
  which shows that the compact closed transpose of the local product is precisely the following extension following a restriction:
    \begin{equation} \label{eqn:transpose-local-product-with-identity}
      \widehat{p\wedge\id} = \coprod_{\Delta,1} (\Delta,1)^*((p\times \id)\odot\Delta\odot x_!).
    \end{equation}
  Now, using this identity, we have a bijection of cells in the following sequence:
    \begin{align*}
      % https://q.uiver.app/#q=WzAsNCxbMCwwLCJ4Il0sWzEsMCwieCJdLFsxLDEsIngiXSxbMCwxLCJ4Il0sWzAsMSwicFxcd2VkZ2VcXGlkIiwwLHsic3R5bGUiOnsiYm9keSI6eyJuYW1lIjoiYmFycmVkIn19fV0sWzEsMiwiIiwwLHsibGV2ZWwiOjIsInN0eWxlIjp7ImhlYWQiOnsibmFtZSI6Im5vbmUifX19XSxbMCwzLCIiLDIseyJsZXZlbCI6Miwic3R5bGUiOnsiaGVhZCI6eyJuYW1lIjoibm9uZSJ9fX1dLFszLDIsIm4iLDIseyJzdHlsZSI6eyJib2R5Ijp7Im5hbWUiOiJiYXJyZWQifX19XSxbNCw3LCJcXGxlcSIsMSx7InNob3J0ZW4iOnsic291cmNlIjoyMCwidGFyZ2V0IjoyMH0sInN0eWxlIjp7ImJvZHkiOnsibmFtZSI6Im5vbmUifSwiaGVhZCI6eyJuYW1lIjoibm9uZSJ9fX1dXQ==
      \begin{tikzcd}[ampersand replacement=\&]
        x \& x \\
        x \& x
        \arrow[""{name=0, anchor=center, inner sep=0}, "{p\wedge\id}"{inner sep=.8ex}, "\shortmid"{marking}, from=1-1, to=1-2]
        \arrow[equals, from=1-1, to=2-1]
        \arrow[equals, from=1-2, to=2-2]
        \arrow[""{name=1, anchor=center, inner sep=0}, "n"'{inner sep=.8ex}, "\shortmid"{marking}, from=2-1, to=2-2]
        \arrow["\leq"{description}, draw=none, from=0, to=1]
      \end{tikzcd} \quad &\leftrightarrow\quad 
      % https://q.uiver.app/#q=WzAsNCxbMCwwLCJ4XFx0aW1lcyB4Il0sWzIsMCwiMSJdLFsyLDEsIjEiXSxbMCwxLCJ4XFx0aW1lcyB4Il0sWzAsMSwiXFx3aWRlaGF0e3BcXHdlZGdlXFxpZH0iLDAseyJzdHlsZSI6eyJib2R5Ijp7Im5hbWUiOiJiYXJyZWQifX19XSxbMSwyLCIiLDAseyJsZXZlbCI6Miwic3R5bGUiOnsiaGVhZCI6eyJuYW1lIjoibm9uZSJ9fX1dLFswLDMsIiIsMix7ImxldmVsIjoyLCJzdHlsZSI6eyJoZWFkIjp7Im5hbWUiOiJub25lIn19fV0sWzMsMiwiXFx3aWRlaGF0IG4iLDIseyJzdHlsZSI6eyJib2R5Ijp7Im5hbWUiOiJiYXJyZWQifX19XSxbNCw3LCJcXGxlcSIsMSx7InNob3J0ZW4iOnsic291cmNlIjoyMCwidGFyZ2V0IjoyMH0sInN0eWxlIjp7ImJvZHkiOnsibmFtZSI6Im5vbmUifSwiaGVhZCI6eyJuYW1lIjoibm9uZSJ9fX1dXQ==
      \begin{tikzcd}[ampersand replacement=\&]
        {x\times x} \&\& 1 \\
        {x\times x} \&\& 1
        \arrow[""{name=0, anchor=center, inner sep=0}, "{\widehat{p\wedge\id}}"{inner sep=.8ex}, "\shortmid"{marking}, from=1-1, to=1-3]
        \arrow[equals, from=1-1, to=2-1]
        \arrow[equals, from=1-3, to=2-3]
        \arrow[""{name=1, anchor=center, inner sep=0}, "{\widehat n}"'{inner sep=.8ex}, "\shortmid"{marking}, from=2-1, to=2-3]
        \arrow["\leq"{description}, draw=none, from=0, to=1]
      \end{tikzcd}\qquad &(\text{Compactness})\\
      & = \quad 
      % https://q.uiver.app/#q=WzAsNCxbMCwwLCJ4XFx0aW1lcyB4Il0sWzQsMCwiMSJdLFs0LDEsIjEiXSxbMCwxLCJ4XFx0aW1lcyB4Il0sWzAsMSwiXFxjb3Byb2Rfe1xcRGVsdGEsMX0gKFxcRGVsdGEsMSleKigocFxcdGltZXMgXFxpZClcXG9kb3RcXERlbHRhXFxvZG90IHhfISkiLDAseyJzdHlsZSI6eyJib2R5Ijp7Im5hbWUiOiJiYXJyZWQifX19XSxbMSwyLCIiLDAseyJsZXZlbCI6Miwic3R5bGUiOnsiaGVhZCI6eyJuYW1lIjoibm9uZSJ9fX1dLFswLDMsIiIsMix7ImxldmVsIjoyLCJzdHlsZSI6eyJoZWFkIjp7Im5hbWUiOiJub25lIn19fV0sWzMsMiwiXFx3aWRlaGF0IG4iLDIseyJzdHlsZSI6eyJib2R5Ijp7Im5hbWUiOiJiYXJyZWQifX19XSxbNCw3LCJcXGxlcSIsMSx7InNob3J0ZW4iOnsic291cmNlIjoyMCwidGFyZ2V0IjoyMH0sInN0eWxlIjp7ImJvZHkiOnsibmFtZSI6Im5vbmUifSwiaGVhZCI6eyJuYW1lIjoibm9uZSJ9fX1dXQ==
      \begin{tikzcd}[ampersand replacement=\&]
        {x\times x} \&\&\&\& 1 \\
        {x\times x} \&\&\&\& 1
        \arrow[""{name=0, anchor=center, inner sep=0}, "{\coprod_{\Delta,1} (\Delta,1)^*((p\times \id)\odot\Delta\odot x_!)}"{inner sep=.8ex}, "\shortmid"{marking}, from=1-1, to=1-5]
        \arrow[equals, from=1-1, to=2-1]
        \arrow[equals, from=1-5, to=2-5]
        \arrow[""{name=1, anchor=center, inner sep=0}, "{\widehat n}"'{inner sep=.8ex}, "\shortmid"{marking}, from=2-1, to=2-5]
        \arrow["\leq"{description}, draw=none, from=0, to=1]
      \end{tikzcd} \qquad &(\text{\cref{eqn:transpose-local-product-with-identity}})\\ 
      &\leftrightarrow \quad
      % https://q.uiver.app/#q=WzAsNCxbMCwwLCJ4XFx0aW1lcyB4Il0sWzMsMCwiMSJdLFszLDEsIjEiXSxbMCwxLCJ4XFx0aW1lcyB4Il0sWzAsMSwiKFxcRGVsdGEsMSleKigocFxcdGltZXMgXFxpZClcXG9kb3RcXERlbHRhXFxvZG90IHhfISkiLDAseyJzdHlsZSI6eyJib2R5Ijp7Im5hbWUiOiJiYXJyZWQifX19XSxbMSwyLCIiLDAseyJsZXZlbCI6Miwic3R5bGUiOnsiaGVhZCI6eyJuYW1lIjoibm9uZSJ9fX1dLFswLDMsIiIsMix7ImxldmVsIjoyLCJzdHlsZSI6eyJoZWFkIjp7Im5hbWUiOiJub25lIn19fV0sWzMsMiwiKFxcRGVsdGEsMSleKihcXHdpZGVoYXQgbikiLDIseyJzdHlsZSI6eyJib2R5Ijp7Im5hbWUiOiJiYXJyZWQifX19XSxbNCw3LCJcXGxlcSIsMSx7InNob3J0ZW4iOnsic291cmNlIjoyMCwidGFyZ2V0IjoyMH0sInN0eWxlIjp7ImJvZHkiOnsibmFtZSI6Im5vbmUifSwiaGVhZCI6eyJuYW1lIjoibm9uZSJ9fX1dXQ==
      \begin{tikzcd}[ampersand replacement=\&]
        {x\times x} \&\&\& 1 \\
        {x\times x} \&\&\& 1
        \arrow[""{name=0, anchor=center, inner sep=0}, "{(\Delta,1)^*((p\times \id)\odot\Delta\odot x_!)}"{inner sep=.8ex}, "\shortmid"{marking}, from=1-1, to=1-4]
        \arrow[equals, from=1-1, to=2-1]
        \arrow[equals, from=1-4, to=2-4]
        \arrow[""{name=1, anchor=center, inner sep=0}, "{(\Delta,1)^*(\widehat n)}"'{inner sep=.8ex}, "\shortmid"{marking}, from=2-1, to=2-4]
        \arrow["\leq"{description}, draw=none, from=0, to=1]
      \end{tikzcd} \qquad &(\text{Right Adjoint})\\
      &\leftrightarrow \quad
      % https://q.uiver.app/#q=WzAsNCxbMCwwLCJ4XFx0aW1lcyB4Il0sWzMsMCwiMSJdLFszLDEsIjEiXSxbMCwxLCJ4XFx0aW1lcyB4Il0sWzAsMSwiKChwXFx0aW1lcyBcXGlkKVxcb2RvdFxcRGVsdGFcXG9kb3QgeF8hKSIsMCx7InN0eWxlIjp7ImJvZHkiOnsibmFtZSI6ImJhcnJlZCJ9fX1dLFsxLDIsIiIsMCx7ImxldmVsIjoyLCJzdHlsZSI6eyJoZWFkIjp7Im5hbWUiOiJub25lIn19fV0sWzAsMywiIiwyLHsibGV2ZWwiOjIsInN0eWxlIjp7ImhlYWQiOnsibmFtZSI6Im5vbmUifX19XSxbMywyLCJcXHByb2Rfe1xcRGVsdGEsMX0oXFxEZWx0YSwxKV4qKFxcd2lkZWhhdCBuKSIsMix7InN0eWxlIjp7ImJvZHkiOnsibmFtZSI6ImJhcnJlZCJ9fX1dLFs0LDcsIlxcbGVxIiwxLHsic2hvcnRlbiI6eyJzb3VyY2UiOjIwLCJ0YXJnZXQiOjIwfSwic3R5bGUiOnsiYm9keSI6eyJuYW1lIjoibm9uZSJ9LCJoZWFkIjp7Im5hbWUiOiJub25lIn19fV1d
      \begin{tikzcd}[ampersand replacement=\&]
        {x\times x} \&\&\& 1 \\
        {x\times x} \&\&\& 1
        \arrow[""{name=0, anchor=center, inner sep=0}, "{((p\times \id)\odot\Delta\odot x_!)}"{inner sep=.8ex}, "\shortmid"{marking}, from=1-1, to=1-4]
        \arrow[equals, from=1-1, to=2-1]
        \arrow[equals, from=1-4, to=2-4]
        \arrow[""{name=1, anchor=center, inner sep=0}, "{\prod_{\Delta,1}(\Delta,1)^*(\widehat n)}"'{inner sep=.8ex}, "\shortmid"{marking}, from=2-1, to=2-4]
        \arrow["\leq"{description}, draw=none, from=0, to=1]
      \end{tikzcd} \qquad &(\text{Right Adjoint})\\
      &\leftrightarrow\quad 
      % https://q.uiver.app/#q=WzAsNCxbMCwwLCJ4Il0sWzMsMCwieCJdLFszLDEsIngiXSxbMCwxLCJ4Il0sWzAsMSwicCIsMCx7InN0eWxlIjp7ImJvZHkiOnsibmFtZSI6ImJhcnJlZCJ9fX1dLFsxLDIsIiIsMCx7ImxldmVsIjoyLCJzdHlsZSI6eyJoZWFkIjp7Im5hbWUiOiJub25lIn19fV0sWzAsMywiIiwyLHsibGV2ZWwiOjIsInN0eWxlIjp7ImhlYWQiOnsibmFtZSI6Im5vbmUifX19XSxbMywyLCJcXG92ZXJsaW5le1xccHJvZF97XFxEZWx0YSwxfShcXERlbHRhLDEpXiooXFx3aWRlaGF0IG4pfSIsMix7InN0eWxlIjp7ImJvZHkiOnsibmFtZSI6ImJhcnJlZCJ9fX1dLFs0LDcsIlxcbGVxIiwxLHsic2hvcnRlbiI6eyJzb3VyY2UiOjIwLCJ0YXJnZXQiOjIwfSwic3R5bGUiOnsiYm9keSI6eyJuYW1lIjoibm9uZSJ9LCJoZWFkIjp7Im5hbWUiOiJub25lIn19fV1d
      \begin{tikzcd}[ampersand replacement=\&]
        x \&\&\& x \\
        x \&\&\& x
        \arrow[""{name=0, anchor=center, inner sep=0}, "p"{inner sep=.8ex}, "\shortmid"{marking}, from=1-1, to=1-4]
        \arrow[equals, from=1-1, to=2-1]
        \arrow[equals, from=1-4, to=2-4]
        \arrow[""{name=1, anchor=center, inner sep=0}, "{\overline{\prod_{\Delta,1}(\Delta,1)^*(\widehat n)}}"'{inner sep=.8ex}, "\shortmid"{marking}, from=2-1, to=2-4]
        \arrow["\leq"{description}, draw=none, from=0, to=1]
      \end{tikzcd}\qquad &(\text{Compactness})
    \end{align*}
  which proves that $-\wedge\id$ has a right adjoint, so that $\id$ is exponentiable as in \cref{def:exponentiable-and-cartesianclosed}, as claimed.
\end{proof}

Some notation for the main result needs to be established. Use `$\id\supset n$' as notation for each right adjoint as in \cref{lemma:identity-proarrow-exponentiable}. Thus, we have the adjunction statement $p\wedge\id\leq n$ iff $p\leq \id\supset n$. Likewise, for $m\colon x\proto y$ and its tabulator $Tm$, we write $m^*(n)$ for the proarrow domain of the restriction cell 
  \begin{equation*}
    % https://q.uiver.app/#q=WzAsNCxbMCwxLCJ4Il0sWzEsMSwieSJdLFswLDAsIlRtIl0sWzEsMCwiVG0iXSxbMCwxLCJuIiwyLHsic3R5bGUiOnsiYm9keSI6eyJuYW1lIjoiYmFycmVkIn19fV0sWzIsMCwiZCIsMl0sWzIsMywibV4qKG4pIiwwLHsic3R5bGUiOnsiYm9keSI6eyJuYW1lIjoiYmFycmVkIn19fV0sWzMsMSwiYyJdLFs2LDQsIlxcbWF0aHJte3Jlc30iLDEseyJzaG9ydGVuIjp7InNvdXJjZSI6MjAsInRhcmdldCI6MjB9LCJzdHlsZSI6eyJib2R5Ijp7Im5hbWUiOiJub25lIn0sImhlYWQiOnsibmFtZSI6Im5vbmUifX19XV0=
    \begin{tikzcd}
      Tm & Tm \\
      x & y
      \arrow[""{name=0, anchor=center, inner sep=0}, "{m^*(n)}"{inner sep=.8ex}, "\shortmid"{marking}, from=1-1, to=1-2]
      \arrow["d"', from=1-1, to=2-1]
      \arrow["c", from=1-2, to=2-2]
      \arrow[""{name=1, anchor=center, inner sep=0}, "n"'{inner sep=.8ex}, "\shortmid"{marking}, from=2-1, to=2-2]
      \arrow["{\mathrm{res}}"{description}, draw=none, from=0, to=1]
    \end{tikzcd}
  \end{equation*}
Likewise $\prod_m$ and $\Sigma_m$ denote the right- and left-adjoints to substitution along the domain and codomain morphisms coming with $Tm$. As in \cref{def:tabulator} write $t_m$ for the span formed by the legs of the tabulator of $m$. 

\begin{theorem} \label{theo:locally-cartesian-closed}
  If $\dbl D$ is a $\prod$-double olog with strong tabulators, then each hom category $\dbl D(x,y)$ is cartesian closed and the formulas
    \begin{equation*}
      m\wedge n = \coprod_m(\id_{Tm}\wedge m^*(n))\qquad\qquad m\supset p = \prod_m(\id_{Tm}\supset m^*(p))
    \end{equation*}
  compute the product and exponential.
\end{theorem}
\begin{proof}
  The local product formula is immediate by strength and Frobenius:
    \begin{equation*}
      m\wedge n = (\coprod_m\id_{Tm})\wedge n\cong \coprod_m(\id_{Tm}\wedge m^*(n)).
    \end{equation*}
  Apply \cref{lemma:identity-proarrow-exponentiable} to the local product in the coproduct on the right. This gives:
    \begin{equation*}
      m\wedge n \leq p \qquad \text{iff} \qquad n\leq \prod_m(\id_{Tm}\supset m^*(p))
    \end{equation*}
  as claimed.
\end{proof}

\begin{corollary}
  Any $\prod$-double olog with strong tabulators is locally cartesian closed as a fibration (cf. \cite[\S 1.8]{jacobs1999}).
\end{corollary}
\begin{proof}
  Any `double category of relations' satisfies Beck-Chevalley and Frobenius reciprocity. The previous shows each fiber is cartesian closed. Reindexing preserves the closed structure by Beck-Chevalley.
\end{proof}

\begin{remark}
  The formulas in \cref{theo:locally-cartesian-closed} closely resemble and in fact are inspired by, or even based upon, those appearing in \cite[Proposition 10.5.4]{jacobs1999}. However, the formulas in the reference appear in the entirely different context of \emph{closed comprehension categories} and provide fiberwise cartesian closed structure as a result of assumptions placed on that structure. Now, an original version of \cref{theo:locally-cartesian-closed} mimicked this development by asking for a highly structured comprehension scheme \cite[\S 8]{lambert2022} via tabulators and valued in spans as in \cite{niefield2012}. The argument of \cite[Proposition 10.5.4]{jacobs1999} runs almost exactly: local products are almost seen to be obtained as a successive application of left adjoints. There is one issue, however, namely, that without in some way inserting a diagonal, the sequence of adjunction bijections does not end up producing the local product in the receiving structure of the comprehension scheme. In the reference this is the crucial step that shows that the left adjoint formula for the local product works. Due to the \emph{split contexts} of proarrows in double categories, this necessary move fails here. What was recognized is that the needed effect can be achieved by intersecting with an identity predicate, hence the local product with identity proarrows considered in \cref{lemma:identity-proarrow-exponentiable}. Now, to make this work, however, the identity proarrows must actually behave like identity predicates. The former technical key for this is the assumption of \emph{unit-purity}. However, this has been made redundant as it is implied by discreteness in `double categories of relations'. Thus, altogether, the result of \cref{theo:locally-cartesian-closed} is that the \emph{native structure} of a $\prod$-double olog with strong tabulators is \emph{always homwise cartesian closed} without the additional immense structuring assumptions on a comprehension scheme. Note also that the formulas in \cref{theo:locally-cartesian-closed} are almost exactly those in \cite[Proposition 10.5.4]{jacobs1999} except on the product side adjusted by interesection with an identity and on the implication side adjusted by the identity antecedent within the dependent product. Notice also that we have not crossed notation. Unit-purity implies that the general implication formula applied to the former case recovers the identity exponential.
\end{remark}

\section{Cocartesianness \& Negation}
\label{section:cocartesian-negation}

The definition of a \emph{cocartesian double category} is an evident one on the pattern established by instantiating the abstract 2-categorical definition of a (co)cartesian object in the 2-category of double categories, pseudo double functors and tight transformations \cite{aleiferi2018}. It has been stated in an equivalent form and studied extensively from the standpoint of hypergraph categories and decorated cospans in \cite{patterson2023}. For our purposes, just as cartesianness leads to local products/conjunction $\wedge$ and local terminals/truth $\top$, cocartesianness leads to local disjunction $\vee$ and logical falsity/absurdity $\bot$.

\begin{definition}\label{def:cocartesian}
  A double category $\dbl D$ is \textbf{cocartesian} if the double functors 
    \begin{equation*}
      \Delta\colon\dbl D\to\dbl D\times\dbl D \qquad\text{and}\qquad !\colon \dbl D\to 1 
    \end{equation*}
  have left adjoints in the 2-category of double categories, pseudo double functors and tight transformations denoted by $+$ and $\emptyset$, respectively.
\end{definition}

\begin{remark}
  As observed in \cite{lambert2024} for cartesian double categories, by duality a cocartesian double category will have finite coproducts in each of the two categories $\dbl D_0$ and $\dbl D_1$. These should be preserved appropriately by the external structure functors of the double category. The tight codiagonals and initials will be denoted 
    \begin{equation*}
      \nabla_x\colon x+x\to x\qquad\text{and}\qquad \init\colon\emptyset\to x
    \end{equation*}
  Additionally, a cocartesian equipment will have local coproducts and initials in the following way.
\end{remark}

\begin{lemma} \label{lemma:local-coproducts}
  A cocartesian equipment has finite coproducts locally.
\end{lemma}
\begin{proof}
  The construction is dual to that of local products. These are given by the following extensions:
    \begin{equation*}
      % https://q.uiver.app/#q=WzAsNCxbMCwwLCJ4K3giXSxbMSwwLCJ5K3kiXSxbMSwxLCJ5Il0sWzAsMSwieCJdLFswLDEsIm0rbiIsMCx7InN0eWxlIjp7ImJvZHkiOnsibmFtZSI6ImJhcnJlZCJ9fX1dLFsxLDIsIlxcbmFibGFfeSJdLFswLDMsIlxcbmFibGFfeCIsMl0sWzMsMiwibVxcdmVlIG4iLDIseyJzdHlsZSI6eyJib2R5Ijp7Im5hbWUiOiJiYXJyZWQifX19XSxbNCw3LCJcXG1hdGhybXtleHR9IiwxLHsic2hvcnRlbiI6eyJzb3VyY2UiOjIwLCJ0YXJnZXQiOjIwfSwic3R5bGUiOnsiYm9keSI6eyJuYW1lIjoibm9uZSJ9LCJoZWFkIjp7Im5hbWUiOiJub25lIn19fV1d
      \begin{tikzcd}
        {x+x} & {y+y} \\
        x & y
        \arrow[""{name=0, anchor=center, inner sep=0}, "{m+n}"{inner sep=.8ex}, "\shortmid"{marking}, from=1-1, to=1-2]
        \arrow["{\nabla_x}"', from=1-1, to=2-1]
        \arrow["{\nabla_y}", from=1-2, to=2-2]
        \arrow[""{name=1, anchor=center, inner sep=0}, "{m\vee n}"'{inner sep=.8ex}, "\shortmid"{marking}, from=2-1, to=2-2]
        \arrow["{\mathrm{ext}}"{description}, draw=none, from=0, to=1]
      \end{tikzcd}\qquad\qquad 
      % https://q.uiver.app/#q=WzAsNCxbMCwwLCJcXGVtcHR5c2V0Il0sWzEsMCwiXFxlbXB0eXNldCJdLFsxLDEsInkiXSxbMCwxLCJ4Il0sWzAsMSwiXFxpZF9cXGVtcHR5c2V0IiwwLHsic3R5bGUiOnsiYm9keSI6eyJuYW1lIjoiYmFycmVkIn19fV0sWzEsMiwiISJdLFswLDMsIiEiLDJdLFszLDIsIlxcYm90IiwyLHsic3R5bGUiOnsiYm9keSI6eyJuYW1lIjoiYmFycmVkIn19fV0sWzQsNywiXFxtYXRocm17ZXh0fSIsMSx7InNob3J0ZW4iOnsic291cmNlIjoyMCwidGFyZ2V0IjoyMH0sInN0eWxlIjp7ImJvZHkiOnsibmFtZSI6Im5vbmUifSwiaGVhZCI6eyJuYW1lIjoibm9uZSJ9fX1dXQ==
      \begin{tikzcd}
        \emptyset & \emptyset \\
        x & y
        \arrow[""{name=0, anchor=center, inner sep=0}, "{\id_\emptyset}"{inner sep=.8ex}, "\shortmid"{marking}, from=1-1, to=1-2]
        \arrow["{\init}"', from=1-1, to=2-1]
        \arrow["{\init}", from=1-2, to=2-2]
        \arrow[""{name=1, anchor=center, inner sep=0}, "\bot"'{inner sep=.8ex}, "\shortmid"{marking}, from=2-1, to=2-2]
        \arrow["{\mathrm{ext}}"{description}, draw=none, from=0, to=1]
      \end{tikzcd}
    \end{equation*}
  Coproduct injections to $m\vee n$ are induced from those into the coproduct $m+n$. 
\end{proof}

\begin{definition} \label{def:negation-operator}
  For objects $x$ and $y$ in a locally cartesian closed equipment with initial objects locally, the corresponding \textbf{negation operator} is the implication operator 
    \begin{equation*}
      (-)\Rightarrow\bot\colon \dbl D(x,y) \to \dbl D(x,y) \qquad\qquad p\;\;\mapsto\;\; \neg p := p\Rightarrow \bot.
    \end{equation*}
  into the local initial $\bot$.
\end{definition}

\begin{example} \label{example:negation}
  We end the section with some examples. In $\Rel$, negation is the usual \emph{relational complement}. That is for $R\colon X\proto Y$, negation $R\Rightarrow\bot$ is calculated as 
    \begin{equation*}
      R\Rightarrow\bot = \lbrace (x,y)\mid \neg (xRy)\rbrace
    \end{equation*}
  that is, the set of pairs \emph{not} related under $R$. As an example type of query, consider the table 
    \begin{equation*}
      \begin{tabular}{| l | l | }
          \hline\multicolumn{2}{| c |}{\bf Inventory}\\
          \hline {\bf Apothecary }&{\bf Ingredient}\\
          \hline Arcadia's Cauldron & giant's toe \\
          \hline Arcadia's Cauldron & wheat \\
          \hline Arcadia's Cauldron & void salts \\
          \hline Elgrim's Elixers & giant's toe \\
          \hline Elgrim's Elixers & swamp fungal pod \\
          \hline The White Phial & Daedra heart \\
          \hline
      \end{tabular} 
    \end{equation*}
  instancing again the Inventory olog. Suppose that we are willing to work for our wheat but we still need the giant's toe. Suppose also that we are deathly allergic to swamp fungal pod. We will not even go in any shops carrying it. If we just want to know which shops do not carry swamp fungal pod, we restrict along the appropriate global element:
    \begin{equation*}
      % https://q.uiver.app/#q=WzAsNCxbMCwwLCJcXGZib3h7QXBvdGhlY2FyeX0iXSxbNCwwLCIxIl0sWzQsMSwiXFxmYm94e0luZ3JlZGllbnR9Il0sWzAsMSwiXFxmYm94e0Fwb3RoZWNhcnl9Il0sWzAsMSwiXFx0ZXh0cm17U2hvcHMgd2l0aCBzd2FtcCBmdW5nYWwgcG9kfSIsMCx7InN0eWxlIjp7ImJvZHkiOnsibmFtZSI6ImJhcnJlZCJ9fX1dLFsxLDIsIlxcdGV4dHtzd2FtcCBmdW5nYWwgcG9kfSJdLFswLDMsIiIsMix7ImxldmVsIjoyLCJzdHlsZSI6eyJoZWFkIjp7Im5hbWUiOiJub25lIn19fV0sWzMsMiwiXFxtYXRocm17SW52ZW50b3J5fSIsMix7InN0eWxlIjp7ImJvZHkiOnsibmFtZSI6ImJhcnJlZCJ9fX1dLFs0LDcsIlxcbWF0aHJte3Jlc30iLDEseyJzaG9ydGVuIjp7InNvdXJjZSI6MjAsInRhcmdldCI6MjB9LCJzdHlsZSI6eyJib2R5Ijp7Im5hbWUiOiJub25lIn0sImhlYWQiOnsibmFtZSI6Im5vbmUifX19XV0=
      \begin{tikzcd}
        {\fbox{Apothecary}} &&&& 1 \\
        {\fbox{Apothecary}} &&&& {\fbox{Ingredient}}
        \arrow[""{name=0, anchor=center, inner sep=0}, "{\textrm{Shops with swamp fungal pod}}"{inner sep=.8ex}, "\shortmid"{marking}, from=1-1, to=1-5]
        \arrow[equals, from=1-1, to=2-1]
        \arrow["{\text{swamp fungal pod}}", from=1-5, to=2-5]
        \arrow[""{name=1, anchor=center, inner sep=0}, "{\mathrm{Inventory}}"'{inner sep=.8ex}, "\shortmid"{marking}, from=2-1, to=2-5]
        \arrow["{\mathrm{res}}"{description}, draw=none, from=0, to=1]
      \end{tikzcd}
    \end{equation*}
  then negate and restrict:
    \begin{equation*}
      % https://q.uiver.app/#q=WzAsNCxbMCwxLCJcXGZib3h7QXBvdGhlY2FyeX0iXSxbNSwxLCJcXGZib3h7SW5ncmVkaWVudH0iXSxbNSwwLCIxIl0sWzAsMCwiXFxmYm94e0Fwb3RoZWNhcnl9Il0sWzAsMSwiXFxuZWcoXFx0ZXh0cm17U2hvcHMgd2l0aCBzd2FtcCBmdW5nYWwgcG9kfSkiLDIseyJzdHlsZSI6eyJib2R5Ijp7Im5hbWUiOiJiYXJyZWQifX19XSxbMiwxLCJcXHRleHR7U3dhbXAgZnVuZ2FsIHBvZH0iXSxbMywwLCIiLDIseyJsZXZlbCI6Miwic3R5bGUiOnsiaGVhZCI6eyJuYW1lIjoibm9uZSJ9fX1dLFszLDIsIlxcdGV4dHJte1Nob3BzIHdpdGhvdXQgc3dhbXAgZnVuZ2FsIHBvZH0iLDAseyJzdHlsZSI6eyJib2R5Ijp7Im5hbWUiOiJiYXJyZWQifX19XSxbNyw0LCJcXG1hdGhybXtyZXN9IiwxLHsic2hvcnRlbiI6eyJzb3VyY2UiOjIwLCJ0YXJnZXQiOjIwfSwic3R5bGUiOnsiYm9keSI6eyJuYW1lIjoibm9uZSJ9LCJoZWFkIjp7Im5hbWUiOiJub25lIn19fV1d
      \begin{tikzcd}
        {\fbox{Apothecary}} &&&&& 1 \\
        {\fbox{Apothecary}} &&&&& {\fbox{Ingredient}}
        \arrow[""{name=0, anchor=center, inner sep=0}, "{\textrm{Shops without swamp fungal pod}}"{inner sep=.8ex}, "\shortmid"{marking}, from=1-1, to=1-6]
        \arrow[equals, from=1-1, to=2-1]
        \arrow["{\text{Swamp fungal pod}}", from=1-6, to=2-6]
        \arrow[""{name=1, anchor=center, inner sep=0}, "{\neg(\textrm{Shops with swamp fungal pod})}"'{inner sep=.8ex}, "\shortmid"{marking}, from=2-1, to=2-6]
        \arrow["{\mathrm{res}}"{description}, draw=none, from=0, to=1]
      \end{tikzcd}
    \end{equation*}
  returning the table 
    \begin{equation*}
      \begin{tabular}{| l |}
          \hline\multicolumn{1}{| c |}{\bf Shops without Swamp Fungal Pod}\\
          \hline Arcadia's Cauldron \\
          \hline The White Phial \\
          \hline
      \end{tabular}. 
    \end{equation*}
  Note that the negation proarrow as the codomain of the second restriction cell above actually returns a collection of pairs, namely, those shop-ingredient pairs not in the original relation. So, the second filter is actually for those shops \emph{not related} to swamp fungal pod. In other words, we form the relation return those shops \emph{with} the ingredient, form the relational complement of that, and then filter for those shops not carrying the ingredient we want to avoid.

  Now, this is indeed a little elaborate. Alternatively, we can form a conjoint and then negate and compose: 
    \begin{equation*}
      % https://q.uiver.app/#q=WzAsMyxbNiwwLCIxIl0sWzIsMCwiXFxmYm94e0luZ3JlZGllbnR9Il0sWzAsMCwiXFxmYm94e0Fwb3RoZWNhcnl9Il0sWzEsMCwiXFxuZWcoXFx0ZXh0e1N3YW1wIGZ1bmdhbCBwb2R9KV4qIiwwLHsic3R5bGUiOnsiYm9keSI6eyJuYW1lIjoiYmFycmVkIn19fV0sWzIsMSwiXFxtYXRocm17SW52ZW50b3J5fSIsMCx7InN0eWxlIjp7ImJvZHkiOnsibmFtZSI6ImJhcnJlZCJ9fX1dXQ==
      \begin{tikzcd}
        {\fbox{Apothecary}} && {\fbox{Ingredient}} &&&& 1
        \arrow["{\mathrm{Inventory}}"{inner sep=.8ex}, "\shortmid"{marking}, from=1-1, to=1-3]
        \arrow["{\neg(\text{Swamp fungal pod})^*}"{inner sep=.8ex}, "\shortmid"{marking}, from=1-3, to=1-7]
      \end{tikzcd}
    \end{equation*}
  which returns the same result in a simpler fashion but does not utilize in the same way the native functional aspects of the olog. Note that the effect all the way on the right is instanced by exactly all the ingredients other than swamp fungal pod. The composite with its implied existential quantification returns those shops having an ingredient on the list of everything but swamp fungal pod. Now, we can combine this easily with a local conjunction to execute the compound query where we look also for giant's toe:
    \begin{equation*}
      % https://q.uiver.app/#q=WzAsNSxbOCwxLCIxIl0sWzMsMSwiXFxmYm94e0luZ3JlZGllbnR9XFx0aW1lcyBcXGZib3h7SW5ncmVkaWVudH0iXSxbMCwxLCJcXGZib3h7QXBvdGhlY2FyeX1cXHRpbWVzIFxcZmJveHtBcG90aGVjYXJ5fSJdLFswLDAsIlxcZmJveHtBcG90aGVjYXJ5fSJdLFs4LDAsIjEiXSxbMSwwLCJcXG5lZyhcXHRleHR7U3dhbXAgZnVuZ2FsIHBvZH0pXipcXHRpbWVzXFx0ZXh0e0dpYW50J3MgdG9lfSIsMix7InN0eWxlIjp7ImJvZHkiOnsibmFtZSI6ImJhcnJlZCJ9fX1dLFsyLDEsIlxcbWF0aHJte0ludmVudG9yeX1cXHRpbWVzXFxtYXRocm17SW52ZW50b3J5fSIsMix7InN0eWxlIjp7ImJvZHkiOnsibmFtZSI6ImJhcnJlZCJ9fX1dLFszLDIsIlxcRGVsdGEiLDJdLFs0LDAsIiIsMix7ImxldmVsIjoyLCJzdHlsZSI6eyJoZWFkIjp7Im5hbWUiOiJub25lIn19fV0sWzMsNCwiXFx0ZXh0e1Nob3BzIHdpdGggR2lhbnQncyB0b2UgYnV0IHdpdGhvdXQgU3dhbXAgZnVuZ2FsIHBvZH0iLDAseyJzdHlsZSI6eyJib2R5Ijp7Im5hbWUiOiJiYXJyZWQifX19XSxbOSwxLCJcXG1hdGhybXtyZXN9IiwxLHsic2hvcnRlbiI6eyJzb3VyY2UiOjIwfSwic3R5bGUiOnsiYm9keSI6eyJuYW1lIjoibm9uZSJ9LCJoZWFkIjp7Im5hbWUiOiJub25lIn19fV1d
      \resizebox{\textwidth}{!}{%
        \begin{tikzcd}[ampersand replacement=\&]
        {\fbox{Apothecary}} \&\&\&\&\&\&\&\& 1 \\
        {\fbox{Apothecary}\times \fbox{Apothecary}} \&\&\& {\fbox{Ingredient}\times \fbox{Ingredient}} \&\&\&\&\& 1
        \arrow[""{name=0, anchor=center, inner sep=0}, "{\text{Shops with Giant's toe but without Swamp fungal pod}}"{inner sep=.8ex}, "\shortmid"{marking}, from=1-1, to=1-9]
        \arrow["\Delta"', from=1-1, to=2-1]
        \arrow[equals, from=1-9, to=2-9]
        \arrow["{\mathrm{Inventory}\times\mathrm{Inventory}}"'{inner sep=.8ex}, "\shortmid"{marking}, from=2-1, to=2-4]
        \arrow["{\neg(\text{Swamp fungal pod})^*\times\text{Giant's toe}}"'{inner sep=.8ex}, "\shortmid"{marking}, from=2-4, to=2-9]
        \arrow["{\mathrm{res}}"{description}, draw=none, from=0, to=2-4]
        \end{tikzcd}
        }
    \end{equation*}
  returning exactly the table
    \begin{equation*}
      \begin{tabular}{| l |}
          \hline\multicolumn{1}{| c |}{\bf Shops with Giant's Toe but without Swamp Fungal Pod}\\
          \hline Arcadia's Cauldron \\
          \hline
      \end{tabular}. 
    \end{equation*}
  of the single shop that has giant's toe and does not have swamp fungal pod. Note that, as a result of the construction of local products, this process is much like a promonoidal \emph{Day convolution} of the two effects. 
\end{example}
\begin{example} 
  This is also evident in a disjunction query, where for example we are happy to patronize any shop having either of two desired ingredients:
    \begin{equation*}
      % https://q.uiver.app/#q=WzAsNSxbNiwwLCIxIl0sWzMsMCwiXFxmYm94e0luZ3JlZGllbnR9XFx0aW1lcyBcXGZib3h7SW5ncmVkaWVudH0iXSxbMCwwLCJcXGZib3h7QXBvdGhlY2FyeX0rIFxcZmJveHtBcG90aGVjYXJ5fSJdLFswLDEsIlxcZmJveHtBcG90aGVjYXJ5fSJdLFs2LDEsIjEiXSxbMSwwLCJcXHRleHR7V2hlYXR9K1xcdGV4dHtHaWFudCdzIHRvZX0iLDAseyJzdHlsZSI6eyJib2R5Ijp7Im5hbWUiOiJiYXJyZWQifX19XSxbMiwxLCJcXG1hdGhybXtJbnZlbnRvcnl9XFx0aW1lc1xcbWF0aHJte0ludmVudG9yeX0iLDAseyJzdHlsZSI6eyJib2R5Ijp7Im5hbWUiOiJiYXJyZWQifX19XSxbMiwzLCJcXERlbHRhIiwyXSxbNCwwLCIiLDIseyJsZXZlbCI6Miwic3R5bGUiOnsiaGVhZCI6eyJuYW1lIjoibm9uZSJ9fX1dLFszLDQsIlxcdGV4dHtTaG9wcyB3aXRoIFdoZWF0IG9yIEdpYW50J3MgdG9lfSIsMix7InN0eWxlIjp7ImJvZHkiOnsibmFtZSI6ImJhcnJlZCJ9fX1dLFsxLDksIlxcbWF0aHJte2V4dH0iLDEseyJzaG9ydGVuIjp7InNvdXJjZSI6MjB9LCJzdHlsZSI6eyJib2R5Ijp7Im5hbWUiOiJub25lIn0sImhlYWQiOnsibmFtZSI6Im5vbmUifX19XV0=
      \begin{tikzcd}
        {\fbox{Apothecary}+ \fbox{Apothecary}} &&& {\fbox{Ingredient}\times \fbox{Ingredient}} &&& 1 \\
        {\fbox{Apothecary}} &&&&&& 1
        \arrow["{\mathrm{Inventory}+\mathrm{Inventory}}"{inner sep=.8ex}, "\shortmid"{marking}, from=1-1, to=1-4]
        \arrow["\nabla"', from=1-1, to=2-1]
        \arrow["{\text{Wheat}+\text{Giant's toe}}"{inner sep=.8ex}, "\shortmid"{marking}, from=1-4, to=1-7]
        \arrow[""{name=0, anchor=center, inner sep=0}, "{\text{Shops with Wheat or Giant's toe}}"'{inner sep=.8ex}, "\shortmid"{marking}, from=2-1, to=2-7]
        \arrow[equals, from=2-7, to=1-7]
        \arrow["{\mathrm{ext}}"{description}, draw=none, from=1-4, to=0]
      \end{tikzcd}
    \end{equation*}
  This is a local disjunction query, returning the table
    \begin{equation*}
      \begin{tabular}{| l |}
          \hline\multicolumn{1}{| c |}{\bf Shops with Wheat or Giant's Toe}\\
          \hline Arcadia's Cauldron \\
          \hline Elgrim's Elixers \\
          \hline
      \end{tabular}
    \end{equation*}
  of exactly those shops having either one but not necessarily both. We now leave it to the reader's ingenuity to construct other combinations and examples of interest.
\end{example}

\section{Distributivity}
\label{section:distributivity}

Now that we have cocartesian structure, it is worth asking about its interaction with the cartesian structure. In particular, it is natural to ask for a \emph{distributive law}. Phrased for the two operations in the ambient 2-category, it is a \emph{global} distributive law describing the interaction between the two operations in total. A main result of the present work is that global distributivity is stable under localization. That is, suppose we have an equipment that is both cartesian and cocartesian with Frobenius reciprocity for local products. If distributivity holds for this equipment as a cartesian and cocartesian object in the 2-category of double categories, then the induced local products and coproducts are distributive too. In this sense global distributivity localizes. The result requires some set up and notation.

\begin{construction}[Global Distributor] \label{construction:global-distributor}
   Let $\dbl D$ denote a fixed equipment that is both cartesian and cocartesian. Suppose that local products satisfy Frobenius. Now, induced via universal properties are canonical comparison morphisms for both distributivity and interchange: 
    \begin{equation*}
      \gamma \colon (x\times y)+(x\times z) \to x\times (y+z)\qquad\qquad
      \delta\colon (x\times y)+(x\times z) \to (x+x) \times (y+z)
    \end{equation*}
  Subscript indexing will be added as needed. These morphisms are of course related at least by codiagonals:
    \begin{equation*}
      \delta(\nabla\times 1) = \gamma
    \end{equation*}
  as can be checked using the universal properties of the involved arrows. Now additionally, $\delta$ commutes with diagonal and codiagonals in the sense that 
    \begin{equation*}
      \delta(\Delta_x+\Delta_x) = \Delta_{x+x}\qquad\qquad (\nabla_x\times\nabla_x)\delta = \nabla_{x\times x}
    \end{equation*}
  both hold. The interchanger $\delta$ fits into the following commutative diagram:
    \begin{equation*}
      % https://q.uiver.app/#q=WzAsNSxbMCwxLCIoeFxcdGltZXMgeCkrKHhcXHRpbWVzIHgpIl0sWzIsMSwiKHgreClcXHRpbWVzICh4K3gpIl0sWzAsMCwieCt4Il0sWzEsMCwieCJdLFsyLDAsInhcXHRpbWVzIHgiXSxbMCwxLCJcXGRlbHRhIiwyXSxbMiwwLCJcXERlbHRhK1xcRGVsdGEiLDJdLFsyLDMsIlxcbmFibGEiXSxbMyw0LCJcXERlbHRhIl0sWzEsNCwiXFxuYWJsYVxcdGltZXNcXG5hYmxhIiwyXV0=
      \begin{tikzcd}
        {x+x} & x & {x\times x} \\
        {(x\times x)+(x\times x)} && {(x+x)\times (x+x)}
        \arrow["\nabla", from=1-1, to=1-2]
        \arrow["{\Delta+\Delta}"', from=1-1, to=2-1]
        \arrow["\Delta", from=1-2, to=1-3]
        \arrow["\delta"', from=2-1, to=2-3]
        \arrow["{\nabla\times\nabla}"', from=2-3, to=1-3]
      \end{tikzcd}
    \end{equation*}
  and in this sense each object $x$ is much like a bialgebra for the global rig structure of $\dbl D$. Now, there is a canonical comparison cell of the form 
    \begin{equation*}
      % https://q.uiver.app/#q=WzAsNCxbMCwwLCIoeFxcdGltZXMgeSkrKHhcXHRpbWVzIHopIl0sWzIsMCwiKHgnXFx0aW1lcyB5JykrKHgnXFx0aW1lcyB6JykiXSxbMCwxLCJ4XFx0aW1lcyAoeSt6KSJdLFsyLDEsIngnXFx0aW1lcyAoeScreicpIl0sWzAsMSwiKG1cXHRpbWVzIHApKyhtXFx0aW1lcyBxKSAiLDAseyJzdHlsZSI6eyJib2R5Ijp7Im5hbWUiOiJiYXJyZWQifX19XSxbMCwyLCJcXGdhbW1hIiwyXSxbMiwzLCJtXFx0aW1lcyAocCtxKSIsMix7InN0eWxlIjp7ImJvZHkiOnsibmFtZSI6ImJhcnJlZCJ9fX1dLFsxLDMsIlxcZ2FtbWEiXSxbNCw2LCJcXEdhbW1hIiwxLHsic2hvcnRlbiI6eyJzb3VyY2UiOjIwLCJ0YXJnZXQiOjIwfSwic3R5bGUiOnsiYm9keSI6eyJuYW1lIjoibm9uZSJ9LCJoZWFkIjp7Im5hbWUiOiJub25lIn19fV1d
      \begin{tikzcd}
        {(x\times y)+(x\times z)} && {(x'\times y')+(x'\times z')} \\
        {x\times (y+z)} && {x'\times (y'+z')}
        \arrow[""{name=0, anchor=center, inner sep=0}, "{(m\times p)+(m\times q) }"{inner sep=.8ex}, "\shortmid"{marking}, from=1-1, to=1-3]
        \arrow["\gamma"', from=1-1, to=2-1]
        \arrow["\gamma", from=1-3, to=2-3]
        \arrow[""{name=1, anchor=center, inner sep=0}, "{m\times (p+q)}"'{inner sep=.8ex}, "\shortmid"{marking}, from=2-1, to=2-3]
        \arrow["\Gamma"{description}, draw=none, from=0, to=1]
      \end{tikzcd}
    \end{equation*}
  for any such triple of objects of $\dbl D$. This has associated to it globular cells by \emph{sliding}, namely, 
    \begin{equation*}
      % https://q.uiver.app/#q=WzAsNixbMCwwLCIoeFxcdGltZXMgeSkrKHhcXHRpbWVzIHopIl0sWzIsMCwiKHgnXFx0aW1lcyB5JykrKHgnXFx0aW1lcyB6JykiXSxbMSwxLCJ4XFx0aW1lcyAoeSt6KSJdLFszLDAsIngnXFx0aW1lcyAoeScreicpIl0sWzAsMSwiKHhcXHRpbWVzIHkpKyh4XFx0aW1lcyB6KSJdLFszLDEsIngnXFx0aW1lcyAoeScreicpIl0sWzAsMSwiKG1cXHRpbWVzIHApKyhtXFx0aW1lcyBxKSAiLDAseyJzdHlsZSI6eyJib2R5Ijp7Im5hbWUiOiJiYXJyZWQifX19XSxbMSwzLCJcXGdhbW1hXyEiLDAseyJzdHlsZSI6eyJib2R5Ijp7Im5hbWUiOiJiYXJyZWQifX19XSxbMCw0LCIiLDEseyJsZXZlbCI6Miwic3R5bGUiOnsiaGVhZCI6eyJuYW1lIjoibm9uZSJ9fX1dLFs0LDIsIlxcZ2FtbWFfISIsMix7InN0eWxlIjp7ImJvZHkiOnsibmFtZSI6ImJhcnJlZCJ9fX1dLFszLDUsIiIsMix7ImxldmVsIjoyLCJzdHlsZSI6eyJoZWFkIjp7Im5hbWUiOiJub25lIn19fV0sWzIsNSwibVxcdGltZXMgKHArcSkiLDIseyJzdHlsZSI6eyJib2R5Ijp7Im5hbWUiOiJiYXJyZWQifX19XSxbNiwxMSwiXFxHYW1tYV8hIiwxLHsic2hvcnRlbiI6eyJzb3VyY2UiOjIwLCJ0YXJnZXQiOjIwfSwic3R5bGUiOnsiYm9keSI6eyJuYW1lIjoibm9uZSJ9LCJoZWFkIjp7Im5hbWUiOiJub25lIn19fV1d
      \begin{tikzcd}
        {(x\times y)+(x\times z)} && {(x'\times y')+(x'\times z')} & {x'\times (y'+z')} \\
        {(x\times y)+(x\times z)} & {x\times (y+z)} && {x'\times (y'+z')}
        \arrow[""{name=0, anchor=center, inner sep=0}, "{(m\times p)+(m\times q) }"{inner sep=.8ex}, "\shortmid"{marking}, from=1-1, to=1-3]
        \arrow[equals, from=1-1, to=2-1]
        \arrow["{\gamma_!}"{inner sep=.8ex}, "\shortmid"{marking}, from=1-3, to=1-4]
        \arrow[equals, from=1-4, to=2-4]
        \arrow["{\gamma_!}"'{inner sep=.8ex}, "\shortmid"{marking}, from=2-1, to=2-2]
        \arrow[""{name=1, anchor=center, inner sep=0}, "{m\times (p+q)}"'{inner sep=.8ex}, "\shortmid"{marking}, from=2-2, to=2-4]
        \arrow["{\Gamma_!}"{description}, draw=none, from=0, to=1]
      \end{tikzcd}
    \end{equation*}
  and 
    \begin{equation*}
      % https://q.uiver.app/#q=WzAsNixbMSwwLCIoeFxcdGltZXMgeSkrKHhcXHRpbWVzIHopIl0sWzMsMCwiKHgnXFx0aW1lcyB5JykrKHgnXFx0aW1lcyB6JykiXSxbMCwwLCJ4XFx0aW1lcyAoeSt6KSJdLFsyLDEsIngnXFx0aW1lcyAoeScreicpIl0sWzAsMSwieFxcdGltZXMgKHkreikiXSxbMywxLCIoeCdcXHRpbWVzIHknKSsoeCdcXHRpbWVzIHonKSJdLFswLDEsIihtXFx0aW1lcyBwKSsobVxcdGltZXMgcSkgIiwwLHsic3R5bGUiOnsiYm9keSI6eyJuYW1lIjoiYmFycmVkIn19fV0sWzIsMCwiXFxnYW1tYV4qIiwwLHsic3R5bGUiOnsiYm9keSI6eyJuYW1lIjoiYmFycmVkIn19fV0sWzQsMywibVxcdGltZXMgKHArcSkiLDIseyJzdHlsZSI6eyJib2R5Ijp7Im5hbWUiOiJiYXJyZWQifX19XSxbMiw0LCIiLDIseyJsZXZlbCI6Miwic3R5bGUiOnsiaGVhZCI6eyJuYW1lIjoibm9uZSJ9fX1dLFszLDUsIlxcZ2FtbWFeKiIsMix7InN0eWxlIjp7ImJvZHkiOnsibmFtZSI6ImJhcnJlZCJ9fX1dLFsxLDUsIiIsMCx7ImxldmVsIjoyLCJzdHlsZSI6eyJoZWFkIjp7Im5hbWUiOiJub25lIn19fV0sWzYsOCwiXFxHYW1tYV4qIiwxLHsic2hvcnRlbiI6eyJzb3VyY2UiOjIwLCJ0YXJnZXQiOjIwfSwic3R5bGUiOnsiYm9keSI6eyJuYW1lIjoibm9uZSJ9LCJoZWFkIjp7Im5hbWUiOiJub25lIn19fV1d
      \begin{tikzcd}
        {x\times (y+z)} & {(x\times y)+(x\times z)} && {(x'\times y')+(x'\times z')} \\
        {x\times (y+z)} && {x'\times (y'+z')} & {(x'\times y')+(x'\times z')}
        \arrow["{\gamma^*}"{inner sep=.8ex}, "\shortmid"{marking}, from=1-1, to=1-2]
        \arrow[equals, from=1-1, to=2-1]
        \arrow[""{name=0, anchor=center, inner sep=0}, "{(m\times p)+(m\times q) }"{inner sep=.8ex}, "\shortmid"{marking}, from=1-2, to=1-4]
        \arrow[equals, from=1-4, to=2-4]
        \arrow[""{name=1, anchor=center, inner sep=0}, "{m\times (p+q)}"'{inner sep=.8ex}, "\shortmid"{marking}, from=2-1, to=2-3]
        \arrow["{\gamma^*}"'{inner sep=.8ex}, "\shortmid"{marking}, from=2-3, to=2-4]
        \arrow["{\Gamma^*}"{description}, draw=none, from=0, to=1]
      \end{tikzcd}
    \end{equation*}
  each of which is invertible if the original $\Gamma$ is invertible \cite[Lemma A.7]{patterson2026}. These are called the \emph{commutor} and \emph{cocommutor} cells associated to $\Gamma$. The data of the arrows $\gamma$ and cells $\Gamma$ organize into a tight transformation of double categories 
    \begin{equation*}
      % https://q.uiver.app/#q=WzAsNixbMCwwLCJcXGRibCBEXntcXHRpbWVzM30iXSxbMCwxLCJcXGRibCBEXntcXHRpbWVzNH0iXSxbMCwyLCJcXGRibCBEXntcXHRpbWVzNH0iXSxbMiwyLCJcXGRibCBEXntcXHRpbWVzMn0iXSxbMywyLCJcXGRibCBEIl0sWzMsMCwiXFxkYmwgRF57XFx0aW1lczJ9Il0sWzAsMSwiXFxEZWx0YVxcdGltZXMgMVxcdGltZXMgMSIsMl0sWzEsMiwiMVxcdGltZXMgXFx0YXVcXHRpbWVzIDEiLDJdLFsyLDMsIigtKz0pXFx0aW1lcyAoPys/PykiLDJdLFszLDQsIlxcdGltZXMiLDJdLFswLDUsIjFcXHRpbWVzICgtKz0pIl0sWzUsNCwiXFx0aW1lcyAiXSxbMiw1LCJcXEdhbW1hIiwxLHsic2hvcnRlbiI6eyJzb3VyY2UiOjIwLCJ0YXJnZXQiOjIwfSwibGV2ZWwiOjJ9XV0=
      \begin{tikzcd}
        {\dbl D^{\times3}} &&& {\dbl D^{\times2}} \\
        {\dbl D^{\times4}} \\
        {\dbl D^{\times4}} && {\dbl D^{\times2}} & {\dbl D}
        \arrow["{1\times (-+=)}", from=1-1, to=1-4]
        \arrow["{\Delta\times 1\times 1}"', from=1-1, to=2-1]
        \arrow["{\times }", from=1-4, to=3-4]
        \arrow["{1\times \tau\times 1}"', from=2-1, to=3-1]
        \arrow["\Gamma"{description}, between={0.2}{0.8}, Rightarrow, from=3-1, to=1-4]
        \arrow["{(-+=)\times (?+??)}"', from=3-1, to=3-3]
        \arrow["\times"', from=3-3, to=3-4]
      \end{tikzcd}
    \end{equation*}
  where $\tau$ is the canonical morphism interchanging the factors of the cartesian product. The (co)cartesian structure is said to be \textbf{globally distributive} if the arrow and cell components of $\Gamma$ viewed as a transformation are all invertible. Under this assumption of invertibility note that we thus have an invertible cell 
    \begin{equation*} 
      % https://q.uiver.app/#q=WzAsMTEsWzEsMCwiKHhcXHRpbWVzIHkpKyh4XFx0aW1lcyB6KSJdLFszLDAsIih4J1xcdGltZXMgeScpKyh4J1xcdGltZXMgeicpIl0sWzAsMCwieFxcdGltZXMgKHkreikiXSxbMiwxLCJ4J1xcdGltZXMgKHknK3onKSJdLFswLDEsInhcXHRpbWVzICh5K3opIl0sWzMsMSwiKHgnXFx0aW1lcyB5JykrKHgnXFx0aW1lcyB6JykiXSxbNCwwLCJ4J1xcdGltZXMgKHknK3onKSJdLFs0LDEsIngnXFx0aW1lcyAoeScreicpIl0sWzIsMiwieCdcXHRpbWVzICh5Jyt6JykiXSxbNCwyLCJ4J1xcdGltZXMgKHknK3onKSJdLFswLDIsInhcXHRpbWVzICh5K3opIl0sWzAsMSwiKG1cXHRpbWVzIHApKyhtXFx0aW1lcyBxKSAiLDAseyJzdHlsZSI6eyJib2R5Ijp7Im5hbWUiOiJiYXJyZWQifX19XSxbMiwwLCJcXGdhbW1hXioiLDAseyJzdHlsZSI6eyJib2R5Ijp7Im5hbWUiOiJiYXJyZWQifX19XSxbNCwzLCJtXFx0aW1lcyAocCtxKSIsMix7InN0eWxlIjp7ImJvZHkiOnsibmFtZSI6ImJhcnJlZCJ9fX1dLFsyLDQsIiIsMix7ImxldmVsIjoyLCJzdHlsZSI6eyJoZWFkIjp7Im5hbWUiOiJub25lIn19fV0sWzMsNSwiXFxnYW1tYV4qIiwyLHsic3R5bGUiOnsiYm9keSI6eyJuYW1lIjoiYmFycmVkIn19fV0sWzEsNSwiIiwwLHsibGV2ZWwiOjIsInN0eWxlIjp7ImhlYWQiOnsibmFtZSI6Im5vbmUifX19XSxbMSw2LCJcXGdhbW1hXyEiLDAseyJzdHlsZSI6eyJib2R5Ijp7Im5hbWUiOiJiYXJyZWQifX19XSxbNSw3LCJcXGdhbW1hXyEiLDIseyJzdHlsZSI6eyJib2R5Ijp7Im5hbWUiOiJiYXJyZWQifX19XSxbNiw3LCIiLDEseyJsZXZlbCI6Miwic3R5bGUiOnsiaGVhZCI6eyJuYW1lIjoibm9uZSJ9fX1dLFszLDgsIiIsMSx7ImxldmVsIjoyLCJzdHlsZSI6eyJoZWFkIjp7Im5hbWUiOiJub25lIn19fV0sWzgsOSwiXFxpZCIsMix7InN0eWxlIjp7ImJvZHkiOnsibmFtZSI6ImJhcnJlZCJ9fX1dLFs3LDksIiIsMSx7ImxldmVsIjoyLCJzdHlsZSI6eyJoZWFkIjp7Im5hbWUiOiJub25lIn19fV0sWzQsMTBdLFsxMCw4LCJtXFx0aW1lcyAocCtxKSIsMl0sWzExLDEzLCJcXEdhbW1hXioiLDEseyJzaG9ydGVuIjp7InNvdXJjZSI6MjAsInRhcmdldCI6MjB9LCJzdHlsZSI6eyJib2R5Ijp7Im5hbWUiOiJub25lIn0sImhlYWQiOnsibmFtZSI6Im5vbmUifX19XSxbMTcsMTgsIjEiLDEseyJzaG9ydGVuIjp7InNvdXJjZSI6MjAsInRhcmdldCI6MjB9LCJzdHlsZSI6eyJib2R5Ijp7Im5hbWUiOiJub25lIn0sImhlYWQiOnsibmFtZSI6Im5vbmUifX19XSxbMTMsMjQsIjEiLDEseyJsYWJlbF9wb3NpdGlvbiI6NjAsInNob3J0ZW4iOnsic291cmNlIjoyMCwidGFyZ2V0IjoyMH0sInN0eWxlIjp7ImJvZHkiOnsibmFtZSI6Im5vbmUifSwiaGVhZCI6eyJuYW1lIjoibm9uZSJ9fX1dLFs1LDIxLCJcXGNvbmciLDEseyJsYWJlbF9wb3NpdGlvbiI6NDAsInNob3J0ZW4iOnsidGFyZ2V0IjoyMH0sInN0eWxlIjp7ImJvZHkiOnsibmFtZSI6Im5vbmUifSwiaGVhZCI6eyJuYW1lIjoibm9uZSJ9fX1dXQ==
      \begin{tikzcd}[column sep=small]
        {x\times (y+z)} & {(x\times y)+(x\times z)} && {(x'\times y')+(x'\times z')} & {x'\times (y'+z')} \\
        {x\times (y+z)} && {x'\times (y'+z')} & {(x'\times y')+(x'\times z')} & {x'\times (y'+z')} \\
        {x\times (y+z)} && {x'\times (y'+z')} && {x'\times (y'+z')}
        \arrow["{\gamma^*}"{inner sep=.8ex}, "\shortmid"{marking}, from=1-1, to=1-2]
        \arrow[equals, from=1-1, to=2-1]
        \arrow[""{name=0, anchor=center, inner sep=0}, "{(m\times p)+(m\times q) }"{inner sep=.8ex}, "\shortmid"{marking}, from=1-2, to=1-4]
        \arrow[""{name=1, anchor=center, inner sep=0}, "{\gamma_!}"{inner sep=.8ex}, "\shortmid"{marking}, from=1-4, to=1-5]
        \arrow[equals, from=1-4, to=2-4]
        \arrow[equals, from=1-5, to=2-5]
        \arrow[""{name=2, anchor=center, inner sep=0}, "{m\times (p+q)}"'{inner sep=.8ex}, "\shortmid"{marking}, from=2-1, to=2-3]
        \arrow[from=2-1, to=3-1]
        \arrow["{\gamma^*}"'{inner sep=.8ex}, "\shortmid"{marking}, from=2-3, to=2-4]
        \arrow[equals, from=2-3, to=3-3]
        \arrow[""{name=3, anchor=center, inner sep=0}, "{\gamma_!}"'{inner sep=.8ex}, "\shortmid"{marking}, from=2-4, to=2-5]
        \arrow[equals, from=2-5, to=3-5]
        \arrow[""{name=4, anchor=center, inner sep=0}, "{m\times (p+q)}"', from=3-1, to=3-3]
        \arrow[""{name=5, anchor=center, inner sep=0}, "\id"'{inner sep=.8ex}, "\shortmid"{marking}, from=3-3, to=3-5]
        \arrow["{\Gamma^*}"{description}, draw=none, from=0, to=2]
        \arrow["1"{description}, draw=none, from=1, to=3]
        \arrow["1"{description, pos=0.6}, draw=none, from=2, to=4]
        \arrow["\cong"{description, pos=0.4}, draw=none, from=2-4, to=5]
      \end{tikzcd}
    \end{equation*}
  since $\gamma$, being invertible, induces a proarrow adjoint equivalence. This cell of course collapses to an identity if $\dbl D$ is also locally posetal. Thus, in the loose bicategory of $\dbl {D}$, proarrow distributivity takes the form of what we think of as a \emph{conjugacy rule}. This will be used in the calculation of the following proposition, the main result of the section.
\end{construction}

\begin{theorem} \label{theorem:local-distributivity}
  Let $\dbl D$ denote a cocartesian `double category of relations'. If $\dbl D$ is globally distributive, then local products distribute over local coproducts in the sense that
    \begin{equation*}
      (m\wedge p)\vee (m\wedge q) = m\wedge (p\vee q)
    \end{equation*}
  holds in $\dbl D(x,y)$. In this sense $\dbl D$ enjoys \textbf{local distributivity}.
\end{theorem}
\begin{proof}
  The argument is inspired by the reverse direction of \cite[Proposition 9.2.3]{jacobs1999} inasmuch as Frobenius reciprocity is used to bring the companions and conjoints forming the conjunction as an extension cell outside the local product to let the global distributivity conjugacy isos do the rest of the work. Throughout we use the notation and terminology of \cref{construction:global-distributor} above. In detail, first observe that by Frobenius reciprocity and the construction of local products and coproducts by restrictions and extensions:
    \begin{align*}
      m\wedge(p\vee q) &= (m\wedge\coprod_{\nabla,\nabla}(p+q))\\
      &=\coprod_{\nabla,\nabla}((\nabla,\nabla)^*(m)\wedge(p+q))\\
      &= \nabla^*\odot((\nabla_!\odot m \odot \nabla^*)\wedge (p+q)) \odot \nabla_!\\
      &= \nabla^*\odot(\Delta_!\odot ((\nabla_!\odot m \odot \nabla^*)\times (p+q))\odot\Delta^*) \odot \nabla_!.
    \end{align*}
  Note that we are suppressing the indexing subscripts. Now, working from the last line, we can pull the most imbedded codiagonals outside the product, and then use distributivity as proarrow conjugacy:
    \begin{align*}
      &\nabla^*\odot(\Delta_!\odot ((\nabla_!\odot m \odot \nabla^*)\times (p+q))=\Delta^*) \odot \nabla_! \\
      = \;&\nabla^*\odot\Delta_!\odot(\nabla_!\times 1)\odot (m\times (p+q))\odot (\nabla^*\times 1)\odot \Delta^*\odot\nabla_!\\
      =\;&\nabla^*\odot\Delta_!\odot(\nabla_!\times 1)\odot \gamma^*\odot ((m\times p)+ (m+q))\odot \gamma_!\odot (\nabla^*\times 1)\odot \Delta^*\odot\nabla_!.
    \end{align*}
  Now, from the last line, we use the relationship between $\gamma$ and the interchanger $\delta$, to get 
    \begin{align*}
      &\nabla^*\odot\Delta_!\odot(\nabla_!\times 1)\odot \gamma^*\odot ((m\times p)+ (m+q))\odot \gamma_!\odot (\nabla^*\times 1)\odot \Delta^*\odot\nabla_!\\
      =\;&\nabla^*\odot\Delta_!\odot\delta^*\odot ((m\times p)+ (m+q))\odot \delta_!\odot \Delta^*\odot\nabla_!.
    \end{align*}
  But the interchanger commutes with the diagonals and thus commutes with their companions and conjoints since, being invertible, it induces a proarrow adjoint equivalence. So, the last line immediately above is calculated to give the desired disjunction of conjunctions:
    \begin{align*}
      &\nabla^*\odot\Delta_!\odot\delta^*\odot ((m\times p)+ (m+q))\odot \delta_!\odot \Delta^*\odot\nabla_!\\
      =\;&\nabla^*\odot(\Delta_!+\Delta_!) \odot ((m\times p)+ (m+q))\odot (\Delta^*+\Delta^*)\odot\nabla_!\\
      =\;&\nabla^* \odot (\Delta_!\odot (m\times p)\odot\Delta^*+ \Delta_!\odot(m+q)\odot\Delta^*)\odot\nabla_!\\
      =\;&\nabla^* \odot ((m\wedge p) + (m\wedge q))\odot\nabla_!\\
      =\;&(m\wedge p) \vee (m\wedge q).
    \end{align*}
  Stringing all the calculations together shows that distributivity does hold in $\dbl D(x,y)$, as required.
\end{proof}

\begin{remark}
  The argument above can actually be carried out in the case where $\dbl D$ is merely a cartesian and cocartesian equipment satisfying Frobenius reciprocity and global distributivity. The equalities above are all replaced by isomorphisms and tracking the typing shows that they are all globular. Nonetheless, we phrase the argument in terms of `double categories of relations' since that is the presented forum for double ologs, and required is a fuller account of the relationship between general loose compactness, Frobenius reciprocity, and distributivity, which we leave to a future investigation. Note that we have not assumed anything about the `double category of relations' having any dependent products or being cartesian closed.
\end{remark}

\begin{corollary}
  Any cocartesian `double category of relations' whose coproducts satisfy Beck-Chevalley and with global distributivity in the above sense is a \emph{coherent fibration} in the sense of \cite[Definition 4.2.1]{jacobs1999}.
\end{corollary}
\begin{proof}
  The proof is immediate from \cref{theorem:local-distributivity} and the definition of a coherent fibration.
\end{proof}

\section{Modal FOL \& Description Logic}
\label{section:FOL-DescriptionLogic}

Combining the structures encountered so far we arrive at the richest double-categorical forum for data manipulation and interpretation of logic. This is a double olog with dependent products, comprehension and distributive coproducts. Such a structure is rich enough to interpret first-order predicate logic when viewed as an equipment and satisfies several canonical modal axioms.

\begin{definition}
  A \textbf{FOML double olog} is a double olog with dependent products, strong tabulators, and cocartesian products satisfying global distributivity.
\end{definition}

Since this is a rather heavy definition, we summarize the developed structures and properties, namely, that, a FOML double olog 
  \begin{enumerate}
    \item is a $\prod$-double olog: it has first-order quantification satisfying Beck-Chevalley (\cref{def:prod-double-olog,lemma:beck-chevalley});
    \item has local products and equality satisfying Frobenius reciprocity (\cref{prop:Frobenius-reciprocity});
    \item is locally cartesian closed (\cref{theo:locally-cartesian-closed});
    \item has local disjunction and negation (\cref{lemma:local-coproducts});
    \item is locally distributive (\cref{theorem:local-distributivity});
    \item and interprets modal necessity and possibility (\cref{def:necessity-from-dep-products}).
  \end{enumerate}
This, as has been showcased throughout the development, is more than enough structure to handle classical relational algebra and the database queries developed in \cite{codd1972} as well as modal necessity (safety) and possiblity (liveness). Putting this in the context of fibrational semantics, we have the following culminating result.

\begin{corollary} \label{corollary:first-order-fibration}
  A FOML double olog is a first-order fibration \cite[Definition 4.2.1]{jacobs1999}.
\end{corollary} 
\begin{proof}
  Certianly the source-target projection is a coherent fibration, as observed above. Local cartesian closure and all dependent products are all that is additionally required.
\end{proof} 

As a consequence, any FOML double olog $\dbl D$, viewed as a fibration $\dbl D_1\to\dbl D_0\times\dbl D_0$, interprets first-order predicate logic in the manner described in \cite[\S 4.3]{jacobs1999}. This result is thus one realization of the general proposal and project of viewing structured equipments through the lense of fibrational semantics \cite{lambert2025}. More development on this point will be left to future work (see \cref{section:conclusion-compactness}). Now, with regard to the interaction between modal operators and propositional connectives we have the following. The second is the standard (K) axiom assumed of any sensible modal system.

\begin{proposition}
  In any FOML double olog, the inequalities 
    \begin{enumerate}
      \item $\Box\phi\wedge\Box\psi\leq \Box(\phi\wedge\psi)$
      \item $\Box(\phi\Rightarrow\psi)\leq \Box\phi\Rightarrow\Box\psi$
    \end{enumerate}
  both hold.
\end{proposition}
\begin{proof}
  For the first, $\Box$ is a right adjoint in any case and thus preserves products. Secondly, we have the derivation  
    \begin{prooftree}
      \AxiomC{$\phi\Rightarrow\psi \leq \phi\Rightarrow\psi$}
      \UnaryInfC{$\phi\Rightarrow\psi\wedge\phi \leq \psi$}
      \UnaryInfC{$\Box(\phi\Rightarrow \psi \wedge\phi) \leq \Box\psi$}
      \UnaryInfC{$\Box(\phi\Rightarrow \psi) \wedge\Box\phi \leq \Box\psi$}
      \UnaryInfC{$\Box(\phi\Rightarrow \psi) \leq \Box\phi\Rightarrow\Box\psi$}
    \end{prooftree}
  using cartesian closedness, product preservation, and the fact that in any case $\Box$ is a functor.
\end{proof}

Having developed the previous structures, we note in closing that FOML double ologs are rich enough to interpret description logic \cite{brachman2004,baader2010}. This is the basis of web ontology languages (OWL) used in the development of web3 and a form of knowledge representation. The standard definition of the logic and its set-theoretic semantics is the following.

\begin{definition}
    An \textbf{attributive concept language} (ACL) is given by a signature $\mathcal A = (O,A,R)$, consisting of objects $O$, atomic concepts $A$ and roles $R$, and a set of concepts, defined as the smallest set closed under the following formation rules:
        \begin{enumerate}
            \item $\top$, $\bot$ are concepts;
            \item each atomic concept is a concept;
            \item if $C$ is a concept, so is $\neg C$;
            \item if $C$ and $D$ are concepts, so are $C\cap D$ and $C\cup D$;
            \item if $C$ is a concept and $R$ is a role, then $\forall R.C$ and $\exists R. C$ are concepts.
        \end{enumerate}
An \textbf{interpretation} of an ACL consists of a set $X$ and a correspondence $\llbracket-\rrbracket$ assigning 
    \begin{enumerate}
        \item a set element $\llbracket a\rrbracket \in X$ for each object $a\in O$;
        \item to each concept $C$ a subset $\llbracket C\rrbracket\subset X$;
        \item to each role $R$ a binary relation $\llbracket R\rrbracket \subset X\times X$;
    \end{enumerate}
in such a way that  
    \begin{enumerate}
        \item $\llbracket C\cap D\rrbracket = \llbracket C\rrbracket \cap \llbracket D\rrbracket$
        \item $\llbracket C\cup D\rrbracket = \llbracket C\rrbracket \cup \llbracket D\rrbracket$
        \item $\llbracket \neg C \rrbracket = X\setminus \llbracket C\rrbracket$
        \item $\llbracket\exists R.C\rrbracket = \lbrace x\in X\mid \text{there is } y\in X \text{ such that } (x,y)\in R \text{ and } y\in C\rbrace$ 
        \item $\llbracket\forall R.C\rrbracket = \lbrace x\in X\mid \text{for all } y\in X \text{ if } (x,y)\in R \text{ then } y\in C\rbrace$. 
    \end{enumerate}
  Denote such an interpretation by $\mathcal A \to \Rel$.
\end{definition}

The notation is chosen because, technically, we are using the double-category data of relations without also using the full double-structure. This of course suggests a more general interpretation in any suitably structured `double category of relations'. Note that the interpretation of quantifiers is exactly the modal statements relative to the fixed subset/proposition summarized in \cref{fig:quant-queries}.

\begin{definition} \label{definition:interpret-description-logic}
  An \textbf{interpretation} of an ACL in an FOML double olog $\dbl D$ consists of an object $X$ and a correspondence $\llbracket-\rrbracket$ assigning 
    \begin{enumerate}
        \item a term $\llbracket a\rrbracket \colon u_a\to x$ for each object $a\in O$;
        \item to each concept $C$ a monic $\llbracket C\rrbracket\rightarrowtail x$;
        \item to each role $R$ a proarrow $\llbracket R\rrbracket \colon x\proto x$;
    \end{enumerate}
in such a way that  
    \begin{enumerate}
        \item $\llbracket C\cap D\rrbracket = \llbracket C\rrbracket \wedge \llbracket D\rrbracket$
        \item $\llbracket C\cup D\rrbracket = \llbracket C\rrbracket \vee \llbracket D\rrbracket$
        \item $\llbracket \neg C \rrbracket = \llbracket C\rrbracket\Rightarrow \bot$
        \item $\llbracket\exists R.C\rrbracket = \diamondsuit_R\llbracket C\rrbracket$ 
        \item $\llbracket\forall R.C\rrbracket = \Box_R\llbracket C\rrbracket$
    \end{enumerate}
  where now $\wedge$, $\vee$ et cetera are the operations on $\dbl D(x,x)$. Denote such an interpretation by $\mathcal A \to \dbl D$.
\end{definition}

\begin{remark}
  Note that the modal operators make the interpretation of the quantified concepts over the roles very straightforward. And again this situation of having a relation and a proposition/concept on either the source or target of the relation is exactly our Inventory olog together with the shopping list that began the paper in \cref{section-ologs}.
\end{remark}

\section{Conclusion \& Prospectus: The Role of Compactness}
\label{section:conclusion-compactness}

It has been evident to the author at least since the work on `double categories of relations' that a missing account of \emph{loose compactness} is an profound bottleneck on further development of the present theory (see the discussion of \cite[\S 12]{lambert2022}). In fact this concern goes back to a first exposure to Yoneda structures in \cite{weber2007} where a form of compactness is an essential ingredient in phrasing what it means to be a \emph{2-topos}. Owing to the need of recasting this derived loose structure with a first-class status, the author has had in mind a prospective synthetic axiomatization of profunctors as a double category leading to a notion of a \emph{double topos}. Loose compactness would thus naturally be a central feature. Developments meant to supply this lack in the double-categorical literature have since been initiated in related work \cite{patterson2024a,patterson2024b}. Especially the recent fruitful collaboration on \emph{twisted double functors} \cite{lambert2026} is meant to supply a precise framework for a missing notion of a \emph{twisted adjunction} and a \emph{loose Yoneda theory} (orthogonal to that of \cite{pare2011} in which a general account of loose compactness can be phrased. The point of this concluding section is to discuss the necessity of such a framework and the horizons its successful articulation would open. 

The necessity of a general account of double-categorical loose compactness is implicit in the original \cite{carboni1987} and glaring in \cite{lambert2022}. `Bicategories of relations' and the more general \emph{cartesian bicategories} of the former reference are an alternative to \emph{allegories} \cite{freyd1990} presenting an axiomatization of the \emph{loose bicategorical} structure of relational systems. As powerful as the framework is, however, there are two subtleties that led to the the work of \cite{lambert2022}. The first is the fact that the purely bicategorical approach, in an ironic twist relative to purely functional ologs, leaves out genuinely functional arrows as first-class entities, rather encoding them as certain special relations, namely, the \emph{maps}. As we have already noted, this makes naming, subtyping, and equational reasoning unnecessarily difficult; and the introduction of double categories elegantly solves this problem. A more fundamental issue with the framing of the approach, however, is that a `bicategory of relations' is asked to be \emph{locally posetal}. This indeed has the effect of simplifying many calculations (the derivation of the modular laws and compactness), but also more fundamentally of essentially reorienting the whole framework to that of a \emph{monoidal 2-category} since the bicategorical coherence conditions as a result all hold on the nose.

Now, from the standpoint of double category theory, this assumption thus has the effect of transposing the entire theory from the loose bicategorical direction to the tight 2-categorical direction. Of course, in essence, the entire point of double categories is to keep these as largely independent but related orthogonal directions axiomatizing first-class entities. This is why the equipment structure is necessary in \cite{lambert2022} to show that the loose bicategory of a `double category of relations' is a `bicategory of relations'. For the axioms governing diagonals in the cartesian structure \cite{aleiferi2018} are posed in the \emph{tight direction} and thus have to be transposed via the equipment structure to the loose bicategory. The happy circumstance is that the right adjoints are provided by conjoints; in other words, this generalization is impossible without the automatic symmetry of the full equipment structure. This is not to say that the definition of a `bicategory of relations' is the wrong one; it is just that it is not really about bicategories. The proper generalization should actually be closer to \emph{the underlying monoidal 2-category of a cartesian double category (with suitable further structure) is a `bicategory of relations'}. This explains the considerable effort involved in the follow up paper \cite{carboni2008} which removes the assumption of \emph{locally posetal}. Essentially, in summary, this follow-up paper is inadvertantly trying to conflate categorification and transposition. The work on doing this transposition carefully has appeared only recently in \cite{patterson2026}.

Now, the present work follows \cite{lambert2025} in treating `double categories of relations' as a basic structural framework for double ologs. This of course carries the assumption again that each such olog is locally posetal. This means that from the fibrational perspective, our double ologs are models of a \emph{logic over a type theory} rather than of a \emph{depedent type theory} where there may be more than one morphism between given proarrows. The development of the present work shows that the operative structures in double-olog-ing and querying are really minimally (1) (co)cartesian structure, (2) equipment structure with Beck-Chevalley and for maximal expressivity also (3) dependent products, (4) strong tabulators, and the derived rules of modularity, Frobenius reciprocity and compactness. Now, again, these derived rules are consequences of the `double category of relations' structure and certainly depend for their proofs and form on discreteness and the locally posetal assumption. Now, discreteness is still of interest even if locally posetal is dropped. This is like moving from equivalence relations to general groupoids. In particular, the triviality of the duality involution is specific to relations and such structures. So, on the one hand, the locally posetal assumption is a computational convenience and a simplification making a direction connection between databases and logics over type theories (extending the reflections and program of \cite[\S 2.5]{Spivak2010}. On the other hand, it seems like an artificial limitation on the databases that can be modeled, and on the full connection with dependent type theory that is lurking in the background, inasmuch as the essential features can just be axiomatized: a distributive (co)cartesian equipment with strong tabulators and some suitably weakened versions of Frobenius reciprocity and of Beck-Chevalley for dependent (co)products. 

Whether compactness is axiomatized directly or derived in this more general framework, it presents a problem in that it is not really known what to look for without the locally posetal assumption. In particular, the general case should allow a non-trivial duality involution modeled by \emph{opposite} in profunctors, or perhaps by \emph{reverse path} in directed homotopy theory \cite{grandis2003,grandis2002directed2,grandis2009directedBook}. And, in essence, although it has been studied for 2-categories \cite{shulman2018contravariance}, there is no precise \emph{general theory} recovering oppositization and its interaction with monoidal products in the loose direction for double categories. This, again, on our view, requires the framework of loose adjunctions and a loose Yoneda theory stemming from the recent machinery of twisted double functors and in particular the twisted hom functor \cite{lambert2026}. Points of inspiration for such a theory are of course the ordinary 1-categorical accounts of *-autonomy \cite{barr1999autonomous} and compact closedness \cite{kelly1980coherence} but also the bicategorical approaches to duality and involutions in \cite{daystreet1997,shulman2018contravariance}.

Owing to the locally posetal assumption and the resulting triviality of the accompanying duality involution, compactness in the case of `double categories of relations' seems an apparent triviality. And indeed it could be observed that the bookkeeping problem of remembering which objects are on the left or right side of a relation is merely a nuisance that \emph{ought} to be remedied by something like compactness. It is maintained here, however, that this view is wrong, or at least short-sighted, on at least two counts. The first is of course already apparent in the nature of double ologs. Typically, for example, we naively think of a relation such as \emph{parent of} as pointing from parents to children, and not the other way around, much less from, or to, some product of the two types. Likewise, as a technical matter, relational division in the original \cite{codd1972} is actually framed for two tables representing relations of essentially arbitrary arity. Compactness is in the background grouping by moving columns from one side to the other to make sure the division columns match without incorporating those not involved in the query. Compactness plays non-trivially in both cases; in a salutory way in the latter by making sure the division works properly; but in an artificial way in the former inasmuch as it collapses intended semantic distinctions. In either case, compactness is doing non-trivial work and needs to be better understood.

Now, the second count is of course that this proposed view refuses to see this simpler case as a testing ground for the generalized theory of compactness with non-trivial duality as discussed above. This general account reveals an endemic \emph{sidedness} in that the example of profunctors is presented with opposite on one or the other side. This is necessary also for Lawvere metric spaces \cite{lawvere2002}, due to lack of symmetry, and for order ideals \cite[\S 3.4.6]{grandis2019} and other enriched profunctors. Now, our convention is that profunctors are typed with opposite on the left. This preserves the typing of the hom functors in diagrammatic order. Thus, profunctors from the terminal (for us anyway) are essentially copresheaves and profunctors to the terminal are presheaves. Now, the former are algebraic objects and the latter are geometric. We liken the former to \emph{states} and the latter to \emph{effects} or \emph{generalized propositions} as in CQM \cite{AbramskyCoecke2004,Selinger2007,HeunenVicary2019}. Or again the former to points or certain spaces and the latter to measurements or quantities; in this sense, the adjoint situation of \emph{Isbell conjugacy} \cite[\S 7]{Lawvere2005Reprint} applied to profunctors is already a high-level statement of space-quantity duality and thus also state-effect duality. In this way, we see from a very broad philosophical perspective that although compactness with trivial duality seems to conflate states and effects, it really ought not to in a general account. And arguably this ought not to be thought of as a collapse even in the trivial case. For in the example of matrices, states are column vectors and effects are row vectors. If these are set-indexed, again we have only a trivial duality involution, and although row and column vectors are related by transposition, surely no one would simply conflate them as \emph{the same} or argue that it does not matter how they are arranged. It seems to us worthwhile, then, even in the trivial case, to preserve a semantic distinction even if the bijections say we can move freely between them without dualizing objects/types explicitly.

Moreover, this two-sidedness of \emph{split contexts} for relational or propositional typing in general double ologs is consistent with a broader view of the semantics that especially enriched profunctors already do to provide. For example, the connection between *-autonomy, enriched structures and models of linear logic is well-established \cite{seely1989,barr1991,day2004}. Likewise, profunctors have explicitly and directly been studied as models of linear logic \cite{dunn2015}. \emph{Chu spaces} \cite[Appendix]{barr1979} naturally organize into a *-autonomous category. These all have such split contexts; in paricular, we view each Chu space as a scalar-valued profunctor-like morphism from the product of an underlying space and its frame of opens. That is, in summary, profunctor-type models of linear logic already incorporate some kind of non-trivial duality and an inherent sidedness in variable contexts. Our view is that the fact that the requisite care in the treatment of split contexts appearing as an inevitability in a generalization of locally posetal ologs is not a coincidence. For these locally posetal ologs are local discretizations of more general cartesian equipment-type structures that are generally profunctor-type double categories. And the references above suggest the types of logics and type theories these structures are meant to interpret.

From a type-theoretic perspective, the generalization is represented as a move from logics over type theories to some \emph{dependent type theory} where proarrows, formerly identified as relational propositions in a split context, are now identified as dependent types in a \emph{split context} whose two sides are mediated by a type-theoretic instantiation of compactness. Sidedness is the apparently novel feature. But the technical point is to axiomatize and provide double-categorical semantics for something like \emph{opposite types} \cite{,Zeilberger2009,AgudeloAgudelo2022,AgudeloAgudelo2022a,rivera2026} resulting in a \emph{directed dependent type theory with split contexts} as the internal language of such structures. And in fact oppositization has already been likened to a way to transition between \emph{polarities} from \emph{inductive} to \emph{coinductive} types, or from \emph{positive} to \emph{negative} types, where the former has rules for constructing terms and the latter has rules for deconstructing terms \cite{levy2004cbpv,levy2005adjunction,Zeilberger2009,LicataZeilbergerHarper2008,LicataHarper2009,ahman2016}. This is much like the \emph{prover-denier duality} of game semantics which indeed brings us full-circle back to Frobenius algebras and Frobenius pairs \cite{mellies2012,mellies2016} (i.e. essentially generalized discreteness in the language of \cite{carboni1987} except again with non-trivial duality). Naturally, it is no coincidence that Chu spaces have served as models of concurrency and game semantics \cite{pratt1995}. This is all to say that a successful elementary axiomatization of profunctors as a double topos (as prognosticated in our \cite[\S 12.3]{lambert2022} has an immense payoff, naturally providing semantic models of fragments of linear logic over directed type theories and providing a common framework for discussion of data, programs, concurrency, game semantics, and ultimately, we believe, CQM, synthetic probability theory \cite{fritz2020} and the synthetic branching spacetime models discussed in \cite{lambert2024b}.

\paragraph{Acknowledgements.} Early versions of this work were presented at the 
NU Mathematics Colloquium and at FMCS 2026. Thanks are due to all the 
participants of both meetings and particularly to Marcello LanFranchi, Dorette 
Pronk, and Peter Selinger for several provocative questions. Special thanks to 
Evan Patterson for reviewing an early draft of the manuscript; and to Darien DeWolf for the invitation to speak at FMCS 2026 which prompted much of the work in this paper.

\printbibliography[heading=bibintoc]

\end{document}
